%Version 3 October 2023
% See section 11 of the User Manual for version history
%
%%%%%%%%%%%%%%%%%%%%%%%%%%%%%%%%%%%%%%%%%%%%%%%%%%%%%%%%%%%%%%%%%%%%%%
%%                                                                 %%
%% Please do not use \input{...} to include other tex files.       %%
%% Submit your LaTeX manuscript as one .tex document.              %%
%%                                                                 %%
%% All additional figures and files should be attached             %%
%% separately and not embedded in the \TeX\ document itself.       %%
%%                                                                 %%
%%%%%%%%%%%%%%%%%%%%%%%%%%%%%%%%%%%%%%%%%%%%%%%%%%%%%%%%%%%%%%%%%%%%%

%%\documentclass[referee,sn-basic]{sn-jnl}% referee option is meant for double line spacing

%%=======================================================%%
%% to print line numbers in the margin use lineno option %%
%%=======================================================%%

%%\documentclass[lineno,sn-basic]{sn-jnl}% Basic Springer Nature Reference Style/Chemistry Reference Style

%%======================================================%%
%% to compile with pdflatex/xelatex use pdflatex option %%
%%======================================================%%

%%\documentclass[pdflatex,sn-basic]{sn-jnl}% Basic Springer Nature Reference Style/Chemistry Reference Style


%%Note: the following reference styles support Namedate and Numbered referencing. By default the style follows the most common style. To switch between the options you can add or remove Numbered in the optional parenthesis. 
%%The option is available for: sn-basic.bst, sn-vancouver.bst, sn-chicago.bst%  
 
%%\documentclass[sn-nature]{sn-jnl}% Style for submissions to Nature Portfolio journals
%%\documentclass[sn-basic]{sn-jnl}% Basic Springer Nature Reference Style/Chemistry Reference Style
\documentclass[sn-mathphys-num]{sn-jnl}% Math and Physical Sciences Numbered Reference Style 
%%\documentclass[sn-mathphys-ay]{sn-jnl}% Math and Physical Sciences Author Year Reference Style
%%\documentclass[sn-aps]{sn-jnl}% American Physical Society (APS) Reference Style
%%\documentclass[sn-vancouver,Numbered]{sn-jnl}% Vancouver Reference Style
%%\documentclass[sn-apa]{sn-jnl}% APA Reference Style 
%%\documentclass[sn-chicago]{sn-jnl}% Chicago-based Humanities Reference Style

%%%% Standard Packages
%%<additional latex packages if required can be included here>

\usepackage{graphicx}%
\usepackage{multirow}%
\usepackage{amsmath,amssymb,amsfonts}%
\usepackage{amsthm}%
\usepackage{mathrsfs}%
\usepackage[title]{appendix}%
\usepackage{xcolor}%
\usepackage{textcomp}%
\usepackage{manyfoot}%
\usepackage{booktabs}%
\usepackage{algorithm}%
\usepackage{algorithmicx}%
\usepackage{algpseudocode}%
\usepackage{listings}%
%%%%
\usepackage[all]{xy}
\usepackage{todonotes}
%%%%%=============================================================================%%%%
%%%%  Remarks: This template is provided to aid authors with the preparation
%%%%  of original research articles intended for submission to journals published 
%%%%  by Springer Nature. The guidance has been prepared in partnership with 
%%%%  production teams to conform to Springer Nature technical requirements. 
%%%%  Editorial and presentation requirements differ among journal portfolios and 
%%%%  research disciplines. You may find sections in this template are irrelevant 
%%%%  to your work and are empowered to omit any such section if allowed by the 
%%%%  journal you intend to submit to. The submission guidelines and policies 
%%%%  of the journal take precedence. A detailed User Manual is available in the 
%%%%  template package for technical guidance.
%%%%%=============================================================================%%%%

\newcommand\tP{{\mathcal P}}
\newcommand\tQ{{\mathcal Q}}
\newcommand{\tX}{{\mathcal X}}
\newcommand\tY{{\mathcal Y}}
\newcommand\C{{\mathcal C}}
\newcommand\F{{\mathcal F}}
\newcommand\G{{\mathcal G}}
\newcommand\tS{{\mathcal S}}
\newcommand{\Oo}{\mathcal O}
\newcommand{\Uu}{\mathcal U}
\newcommand{\Ff}{\mathcal F}
\newcommand{\Ee}{\mathcal E}
\newcommand{\Ll}{\mathcal L}
\newcommand{\Zz}{{\mathcal Z}}
\newcommand{\Mmm}{\mathcal M}
\newcommand{\I}{\mathcal I}
\newcommand{\N}{\mathcal N}
\newcommand{\Rr}{\mathcal R}
\newcommand{\D}{\mathcal D}
\newcommand{\T}{\mathcal T}
\newcommand{\tRHom}{{\mathcal R\mathit{Hom}}}

%\newcommand{\RHom}{{\mathrm{hom}}}
\newcommand{\subsetq}{\subset}
%\newcommand{\lest}{\prec}
%\newcommand{\leqt}{\preceq}
%\newcommand{\get}{\succ}
%\newcommand{\geqt}{\succeq}
\newcommand{\Tor}{\mathrm{Tor}}
\newcommand{\RR}{\mathrm R}
\newcommand{\LL}{\mathrm L}
\newcommand{\rep}{{\mathrm{rep}}}
\newcommand{\Char}{{\mathrm{Char}}}
\newcommand{\ind}{{\mathrm{ind}}}
\newcommand{\hgt}{{\mathrm {ht}}}
\newcommand{\res}{{\mathrm{res}}}
\newcommand{\triv}{{\mathrm{triv}}}
\newcommand{\Rind}{{\mathrm{Rind}}}
\newcommand{\hull}{{\mathrm{hull}}}
\newcommand{\id}{{\mathrm{id}}}
\newcommand{\red}{{\mathrm{red}}}
\newcommand{\pre}{{\mathrm{pre}}}
\newcommand{\Ext}{{\mathrm{Ext}}}
%\newcommand{\RHom}{{\mathrm{RHom}}}
\newcommand{\Rep}{{\mathrm{Rep}}}
\newcommand{\Perf}{{\mathrm{Perf}}}
\newcommand{\fg}{{\mathrm {fg}}}
\newcommand{\fr}{{\mathrm {fr}}}
\newcommand*{\Dd}{\mathop{\mathrm D\kern0pt}\nolimits}
\newcommand*{\RGamma}{\mathop{\mathrm {R\Gamma}}\nolimits}
\newcommand*{\Hom}{\mathop{\mathrm {Hom}}\nolimits}
\newcommand*{\RHom}{\mathop{\mathrm {RHom}}\nolimits}
\newcommand*{\Rf}{\mathop{\mathrm {Rf}}\nolimits}
%\newcommand*{\R\pi}{\mathop{\mathrm R{\pi}}\nolimits}
\newcommand*{\Lf}{\mathop{\mathrm {Lf}}\nolimits}
%\newcommand*{\Ext}{\mathop{\mathrm {Ext}}\nolimits}
\newcommand*{\R}{\mathop{\mathrm R}\nolimits}
%\newcommand{\fmod}{{\amod^{\text{{ T\tiny fg}}}}}
%\newcommand{\amod}{{\text{\rm -mod}}}
%\newcommand*{\Ind}{\mathop{\mathrm {ind}}\nolimits}
%\newcommand*{\Res}{\mathop{\mathrm {res}}\nolimits}

\newcommand*{\Coh}{\mathop{\sf {Coh}}\nolimits}
\newcommand{\QCoh}{\mathsf{QCoh}}
\newcommand{\coker}{{\mathsf{coker}}}
\newcommand{\cone}{{\mathsf{cone}}}
\newcommand{\kernel}{{\mathsf{ker}}}
\newcommand{\sd}{\sf D}
\newcommand{\Ss}{\sf S}
\newcommand{\Mm}{\sf M}
\newcommand{\A}{{\sf A}}
\newcommand{\B}{{\sf B}}
\newcommand{\sk}{{\sf k}}
\newcommand{\sK}{{\sf K}}
\newcommand{\sX}{{\sf X}}
\newcommand{\sH}{{\sf H}}

\newcommand{\onto}{{\twoheadrightarrow}}
%\newcommand{\hk}{{\rceil}}
%\newcommand{\sq}{{\square}}
\newcommand{\Gm}{\mathbb G}
\newcommand{\Bm}{\mathbb B}
\newcommand{\Tm}{\mathbb T}
\newcommand{\Pm}{\mathbb P}
\newcommand{\Hm}{\mathbb H}
\newcommand{\Pp}{\mathbb P}
\newcommand{\Q}{{\mathbb Q}}
\newcommand{\Z}{{\mathbb Z}}
\newcommand{\cC}{{\mathbb C}}

\newcommand{\kmod}{\sk\textrm{-}\mathrm{mod}}
\newcommand{\kMod}{\sk\textrm{-}\mathrm{Mod}}
\newcommand*{\vectk}{\mathop{\sf {vect}\textrm-\sk}\nolimits}

\newcommand{\Cc}{\mathbf C}
\newcommand{\bg}{\mathbf G}
\newcommand{\fb}{\mathbf B}
\newcommand{\bt}{\mathbf T}
\newcommand{\bp}{\mathbf P}
\newcommand{\bh}{\mathbf H}
\newcommand{\bu}{\mathbf U}
\newcommand{\bl}{\mathbf L}

\newcommand{\fX}{{ X}}

\newcommand{\oa}{\omega_\alpha}
\newcommand{\ob}{\omega_\beta}

\def\text#1{\mbox{#1}}
%% as per the requirement new theorem styles can be included as shown below
%\theoremstyle{thmstyleone}%
\newtheorem{theorem}{Theorem}[section]
\newtheorem{corollary}[theorem]{Corollary}
\newtheorem{lemma}[theorem]{Lemma}
\newtheorem{proposition}[theorem]{Proposition}
\newtheorem{conjecture}[theorem]{Conjecture}
%\newtheorem{theorem}{Theorem}%  meant for continuous numbers
%%\newtheorem{theorem}{Theorem}[section]% meant for sectionwise numbers
%% optional argument [theorem] produces theorem numbering sequence instead of independent numbers for Proposition
%\newtheorem{proposition}[theorem]{Proposition}% 
%%\newtheorem{proposition}{Proposition}% to get separate numbers for theorem and proposition etc.
\theoremstyle{definition}
%\theoremstyle{thmstyletwo}%
%\newtheorem{example}{Example}%
%\newtheorem{remark}{Remark}%
%\newtheorem{definition}[theorem]{Definition}


%\theoremstyle{thmstylethree}%
\newtheorem{definition}[theorem]{Definition}%
\newtheorem{example}[theorem]{Example}
\newtheorem{remark}[theorem]{Remark}
\newtheorem{aside}[theorem]{Aside}
\newtheorem{notations}[theorem]{Notation}
\newtheorem{notation}[theorem]{Notation}
\raggedbottom
%%\unnumbered% uncomment this for unnumbered level heads
\long\def\comment#1{}
\begin{document}
\setlength{\baselineskip}{12.5pt}

\title[Highest weight category structures and full exceptional collections]{Highest weight category structures on Rep($\fb$) and full exceptional collections on generalized flag varieties over $\mathbb Z$}

\author*[1]{\fnm{Alexander} \sur{Samokhin}}\email{alexander.samokhin@math.uni-bielefeld.de}

\author[2]{\fnm{Wilberd} \spfx{van der}\sur{Kallen}}\email{W.vanderKallen@uu.nl}
%\equalcont{These authors contributed equally to this work.}

%\author[1,2]{\fnm{Third} \sur{Author}}\email{iiiauthor@gmail.com}
%\equalcont{These authors contributed equally to this work.}

\affil[1]{\orgdiv{Fakult\"at f\"ur Mathematik}, \orgname{Universit\"at Bielefeld}, \orgaddress{%\street{}, 
\postcode{33501}
\city{Bielefeld},  \country{Germany}}}

\affil[2]{\orgdiv{Mathematical Institute}, \orgname{Utrecht University}, \orgaddress{\street{P.O.Box 80.010}, \postcode{3508 TA} \city{Utrecht},  \country{The Netherlands}}}

\abstract{Given a split reductive Chevalley group scheme $\Gm$ over $\mathbb Z$ and a parabolic subgroup scheme $\Pm\subset \Gm$, this paper constructs
$\Gm$-linear semiorthogonal decompositions of the bounded derived category of noetherian representations of $\Pm$ with each semiorthogonal component being equivalent to the bounded derived category of noetherian representations of $\Gm$. The  $\Gm$-linear semiorthogonal decompositions in question are compatible with the Bruhat order on cosets of the Weyl group of $\Pm$ in the Weyl group of $\Gm$. %The terms of full exceptional collections are a priori complexes of coherent sheaves; 
%the question whether they are actually coherent sheaves is to be investigated further. 
Their construction builds upon the foundational results on $\Bm$-modules from the works of Mathieu, Polo, and van der Kallen, and upon properties of the Steinberg basis of the $ \Tm$-equivariant $\sK$-theory of $ \Gm/\Bm$. As a corollary, we obtain full exceptional collections in the bounded derived category of coherent sheaves on generalized flag schemes $\Gm/\Pm$ over $\mathbb Z$.}


\keywords{$B$-modules, equivariant $K$-theory, Steinberg basis, flag variety, derived category, semiorthogonal decompositions, exceptional collections}

%%\pacs[JEL Classification]{D8, H51}

%%
\pacs[MSC2020 Classification]{14F08, 14L30, 14M15, 19L47, 20G05, 20G10}

\maketitle

\tableofcontents

\section{Introduction}

%--------------------------------------------------------------
\subsection{The context}
%--------------------------------------------------------------
Let $\Gm$ be a split semi-simple simply connected Chevalley group scheme and $\Bm\subset \Gm$ a Borel subgroup scheme. 
Associated to it is the flag scheme $\Gm/\Bm$ over $\mathbb Z$. More generally, let $\Pm\subset \Gm$ be a parabolic subgroup scheme containing $\Bm$. Generalized flag varietes for the group scheme $\Gm$ are the quotient schemes
$\Gm/\Pm$.

The goal of this paper is to initiate a study of $\Dd ^b(\rep (\Bm))$, the bounded derived category of noetherian representations of $\Bm$ (and, more generally, of the category $\Dd ^b(\rep (\Pm))$). Our main results, which are Theorems \ref{th:semi-orthogonal} and \ref{th:Psemi-orthogonal}, construct some particular $\Gm$-linear semiorthogonal decompositions of $\Dd ^b(\rep (\Bm))$ (resp., $\Gm$-linear semiorthogonal decompositions of $\Dd ^b(\rep (\Pm))$) that can be considered as categorifications of the classical results 
\cite{BGG} and \cite{Dem}. As a corollary to those theorems, we obtain in Theorem \ref{th:coh-semi-orthogonal} (resp., in Theorem \ref{th:Pcoh-semi-orthogonal}) full exceptional collections in the bounded derived category of coherent sheaves on $\Gm/\Bm$ over $\mathbb Z$ (resp., in the bounded derived category of coherent sheaves on $\Gm/\Pm$). 

We believe that the categorical decompositions of $\Dd ^b(\rep (\Bm))$ will have further applications in geometric representation theory (see Section \ref{sec:disc} below). Chronologically, we arrived at Theorems \ref{th:semi-orthogonal} and \ref{th:Psemi-orthogonal} starting off with the question whether full exceptional collections on generalized flag varietes can be obtained using representation theory of a Borel subgroup scheme $\Bm$. We now explain this path in a greater detail.

The study of exceptional collections on generalized flag varieties has a long history; more recently, the paper \cite{KuzPol} set out an approach that influenced many papers on the subject. For a very recent and comprehensive survey of the works that followed {\it loc.cit.}, we refer the reader to \cite{Fon} and the references therein. In this non-technical part of the introduction we emphasise some of the features that make the approach of the present paper different from the previous constructions.

Given a smooth proper Noetherian scheme over a field $\sk$, a full exceptional collection in the bounded derived category of coherent sheaves $\Dd ^b(\Coh(X))$ is, informally speaking, a way to break up $\Dd ^b(\Coh(X))$ into elementary pieces, each equivalent to $\Dd ^b(\vectk)$ which are glued to each other in a non-trivial way. 
Upon decategorification, a full exceptional collection gives rise to a basis of the Grothendieck group ${\sK}^0(X)$.
 If one is interested only in the $\sK$-theoretic information, for many schemes of interest it is easy to produce a basis of ${\sK}^0(X)$: for instance, if a scheme $X$ has a ${\mathbb G}_m$-action with isolated fixed points then by Bia\l{}ynicki-Birula's theorem $X$ has a stratification into locally closed subschemes such that the structure sheaves of their closures are a basis of ${\sK}^0(X)$. 
 Flag schemes $\Gm/\Bm$ furnish a classical example of such a stratification: in this case, the closures of algebraic cells of the  Bia\l{}ynicki-Birula's stratification are Schubert schemes via the Bruhat decomposition. 

Given a scheme $X$ as above and assuming it has a natural geometric basis of the group ${\sK}^0(X)$, one can ask whether such a basis can be lifted to a full exceptional collection in $\Dd ^b(\Coh(X))$. 
Speaking less loosely, a full exceptional collection must satisfy two conditions: it should generate the category $\Dd ^b(\Coh(X))$ in a suitable sense and some cohomological vanishing should hold. 
Generation doesn't pose a problem: for schemes as above, i.e.\ having an algebraic stratification into locally closed subschemes, which are isomorphic to affine spaces, the direct sum of structure sheaves of their closures is a classical generator of $\Dd ^b(\Coh(X))$. 
But there is no chance that for such a set of generators the cohomological vanishing (known as ``semiorthogonality'' condition) would hold. 
On the other hand, there is more to flag schemes that one would want to take into account: there is a partial order on Schubert schemes (the Bruhat order) and the order on the sought-for full exceptional collection is expected to be compatible with the Bruhat order on Schubert schemes. 
Thus, what one is after is finding a correct ``lift'' of the natural $\sK$-theoretic basis to the derived category. 
In the present paper, we provide such a lift. 
As explained below, there are two major inputs into our approach: one is representation-theoretic and comes from the highest weight category structures on the category of rational representations of the Borel group scheme $\Bm$, and the other one is combinatorial and comes from the distinguished basis of the ($\mathbb \bt$-equivariant) $\sK$-group of
the flag scheme $\Gm/\Bm$, called the Steinberg basis, \cite{St}.

There is a basic way to relate algebraic representations of $\Bm$ to $\Gm$-equivariant vector bundles on $\Gm/\Bm$ via the ``associated sheaf'' construction. 
Note that the latter construction was the starting point in \cite{KuzPol} (see Section \ref{sec:disc} below for more on that). Let $\fb$ (resp., $\bg$) denote the group schemes over $\sk$ obtained  
by base change from $\Bm$ (resp., $\Gm$) along ${\rm Spec}(\sk)\rightarrow {\rm Spec}(\mathbb Z)$, where $\sk$ is a field.
The category $\rep(\bg)$ of rational modules over the group scheme $\bg$ that are finite-dimensional over $\sk$ has a highest weight category structure, \cite{CPSHW}, \cite {Donkin1}, and it is these  structures that are expected to manifest themselves through conjectural full exceptional collections. 
In this respect, a model example of such semiorthogonal decompositions/full exceptional collections is given by a theorem of Efimov's \cite{Ef} that solves the problem for the Grassmannians $\sf Gr_{k,n}$ in the best possible way. 
Precisely, one has (cf. also \cite{BLVdB} for a different approach over a field $\sk$):\\

\begin{theorem}[{\cite[Theorem 1.8]{Ef}}]\label{th:Ef}
Let  $\sf Gr_{k,n}$ be the Grassmannian of $k$-dimensional vector subspaces of $n$-dimensional space, defined over $\mathbb Z$. 
There exists a tilting vector bundle on  $\sf Gr_{k,n}$ such that its endomorphism algebra has two natural structures of a split quasi--hereditary algebra over $\mathbb Z$.\\
\end{theorem}


The above theorem \ref{th:Ef} and, more generally, results of \cite{Ef} and \cite{BLVdB} rely considerably on representation theory of the full linear group $GL_n$ and on the highest weight category structure on $\rep(GL_n)$.
Remarkably, through the works \cite{Math}, \cite{Polo3}, \cite{WvdK} and \cite{WvdKTata}, highest weight category structures have also been known to exist on the category $\rep( \fb)$ for a Borel subgroup $ \fb\subset  \bg$. Here $\rep( \fb)$ stands for the category of rational $\fb$-modules that are finite-dimensional over $\sk$.
It seems, however, that the category $\rep( \fb)$ as a highest weight category has not yet received the full attention that it undoubtedly deserves. 
The following quote from Stephen Donkin, \cite[Remark 4.7]{Donsur} had served as a particular impetus for the first named author: ``In fact ${\mathsf k}[{\fb}]$ is a quasi--hereditary coalgebra. 
This follows from van der Kallen's paper \cite{WvdK}. 
This is a deep and sophisticated work which generalizes the above (restriction gives a full embedding of the category of $ \bg$-modules into the category of $\fb$-modules) and whose consequences, to the best of my knowledge, have so far not been investigated or exploited.''

In the present paper, we demonstrate the force of highest weight category structures on the category $\Rep(\fb)$ and, in accordance with Donkin's suggestion, show some of its consequences. Our two main theorems are:

\

\begin{theorem}[Theorem \ref{th:semi-orthogonal}]\label{th:1}
The category $\D=\Dd^b (\rep ({\fb}))$ has a $ \bg$-linear semiorthogonal decomposition 

\begin{equation}
\D= \langle \sf X _v\rangle _{v\in W}
\end{equation}
with respect to  a total order $\prec$ on the Weyl group $W$ that refines the Bruhat order. Each subcategory $\sf X _v$ is equivalent to $\Dd ^b(\rep( \bg))$.
\end{theorem}

\

As it will become clear, the case of Borel subgroups is the most important. With all the setup developed for proving Theorem \ref{th:1}, we are in a position to prove the parabolic version of it. Let $\bp$ be a parabolic subgroup containing $\fb$. (If $\sk=\Z$, then $\bp$ means $\mathbb P$.) Let $W_\bp$ be the parabolic Weyl group corresponding to $\bp$, and $W^\bp$ be the set of minimal coset representatives of $W/W_\bp$. Let $\prec _{\bp}$ denote the restriction to $W^\bp$ of the chosen total order $\prec$ on $W$ from Theorem \ref{th:1} above. Then:

\

\begin{theorem}[Theorem \ref{th:Psemi-orthogonal}]\label{th:2}
%-----------------------------------------------------------------
The category $\D=\Dd^b (\rep ({\bp}))$ has a $ \bg$-linear semiorthogonal decomposition 
\begin{equation}
\D= \langle {\hat{\sf X}} _v\rangle _{v\in W^\bp}
\end{equation}
with respect to the order $\prec _{\bp}$ on  $W^\bp$. Each subcategory ${\hat{\sf X}} _v$ is equivalent to $\Dd ^b(\rep( \bg))$.
\end{theorem}

\

Theorems \ref{th:semi-orthogonal} and \ref{th:coh-semi-orthogonal} have their counterparts over $\mathbb Z$ if one replaces $\bg$ (resp., $\fb$) by $\Gm$ (resp., by $\Bm$). For more details, see Section \ref{sec:SODofB-modoverG-mod} and Section \ref{sec:FEConGmodB}, respectively.

\

Theorems \ref{th:1} and \ref{th:2} have some immediate applications. 

\

\begin{theorem}[Theorem \ref{th:coh-semi-orthogonal}]\label{th:3}
Let $\prec$ be the same total order on $W$ as in Theorem \ref{th:1}, 
and let $\D=\Dd^b (\Coh( \bg/\fb))$. Let $v,w$ denote elements of  $W$. Then there are objects $\tX_v\in \D$ such that 
\begin{enumerate}
\item $\Hom_\D(\tX_v,\tX_v[i])=\begin{cases}\sk &\text{if $i=0$}\\
0&\text{else.}\end{cases}$.

\

\

\item If $w\succ  v$ then $\Hom_\D(\tX_v,\tX_w[i])=0$ for all $i$.\\
\item The triangulated hull of $\{\tX_v\mid v \in W\}$ is $\D$.
\end{enumerate}
In other words, the collection of objects $(\tX_v)_{v \in W}$ is a full exceptional collection in $ \D$.
\end{theorem}

\

The next Theorem \ref{th:4} is in the same position with respect to the previous Theorem \ref{th:3} as Theorem \ref{th:2} is with respect to  Theorem \ref{th:1}: 

\

\begin{theorem}[Theorem \ref{th:Pcoh-semi-orthogonal}]\label{th:Pcoh-semi-orthogonal}\label{th:4}
Let $\prec _{\bp}$ be the same total order on $W^{\bp}$ as in Theorem \ref{th:2}, and let $\D=\Dd^b (\Coh( \bg/\bp))$. Let $v,w\in W^\bp$. Then there are objects $\hat \tX_v\in \D$ such that 
\begin{enumerate}
\item $\Hom_\D(\hat\tX_v,\hat\tX_v[i])=\begin{cases}\sk &\text{if $i=0$},\\
0&\text{else.}\end{cases}$\\



\item If $w\succ _{\bp}v$ then $\Hom_\D(\hat\tX_v,\hat\tX_w[i])=0$ for all $i$.\\
\item The triangulated hull of $\{\hat\tX_v\mid v \in W^\bp\}$ is $\D$.
\end{enumerate}
In other words, the collection of objects $(\hat\tX_v)_{v \in W^\bp}$ is a full exceptional collection in $ \D$.
\end{theorem}

\

In Sections \ref{sec:FEConGmodB} and \ref{sec:caseP} we explain how Theorem \ref{th:3} (resp., Theorem \ref{th:4}) follows immediately from Theorem \ref{th:1} (resp., Theorem \ref{th:2}). Still, our line of thought worked backwards and we conclude this section with a few preliminary remarks explaining our logic. By its very design, the highest weight category structure on an abelian category gives rise to two distinguished collections of objects (standard and costandard objects) that satisfy the $\Ext$-vanishing close to the one that is required for exceptional collections. 
It would therefore be reasonable to start off with one of such collections in $\Rep(\fb)$, convert them into a collection of equivariant vector bundles, and try to ensure the cohomological vanishing conditions in the category $\Dd ^b(\Coh({\bg}/{\fb}))$. 
It turns out that this doesn't work on the nose, but it does work with the input provided by the combinatorics of the Steinberg basis, \cite{St}. Coming back to the problem of lifts of $\sf K$-theoretic classes to the derived category level, we have already mentioned the classical works \cite{BGG} and \cite{Dem} that compute the singular cohomology and, respectively, the $\sf K$-theory of flag varieties $\bg/\fb$ in terms of natural operators called nowadays BGG-Demazure operators. These operators, in different guises, will appear throughout the paper: one of the highest weight category structures on $\rep( \fb)$ is defined in terms of those, see Section \ref{sec:rappelB-mod}. 
We now proceed to the second part of the introduction in which we outline the main steps of the argument.

%--------------------------------------------------------------
\subsection{Outline of the proof}
%--------------------------------------------------------------

Evidence for the principal construction of this paper comes from special full exceptional collections on flag varieties of rank two groups over ${\mathbb Z}[\frac{1}{6}]$, the proofs of which were sketched in \cite{Sam}. We provide the complete details of the rank two case in Section \ref{sec:rank_2}. 
For the purposes of introduction, assume for simplicity that $\Pm=\Bm$ and that
we are working over a field $\sk$  whose characteristic  is not too small.
Thus, $\fb$ and $ \bg$ denote the base change of the group schemes $\Bm$ and $\Gm$ along the morphism ${\rm Spec}(\sk)\rightarrow {\rm Spec}(\mathbb Z)$, \cite[Section 2.6]{Donkin2}. 
Given a finite dimensional $\fb$-module $M$, let $\Ll(M)$ denote the associated (locally free) coherent sheaf on $ \bg/ \fb$. 
%\todo[inline]{
The full exceptional collections from Section \ref{sec:rank_2} are given by locally free sheaves that are obtained via the forgetful functor $\Rep( \fb)\rightarrow {\Coh}({ \bg/\fb})$ and therefore have underlying $\fb$-module structure. That the terms of those collections are $\bg$-equivariant vector bundles is in accordance with a general fact on exceptional objects on smooth projective varieties acted upon by a linear algebraic group, \cite[Lemma 2.2]{Polishchuk}. More importantly, the collections in question are obtained with the help of Demazure functors (see Section \ref{subsec:Demfun}). That points out to the fact that the highest weight category structures on $\rep(\fb)$ can be relevant: as was already mentioned, one of those structures is defined using Demazure operators, see Section \ref{subsec:excord}. 

With those rank two examples in hand, it has become apparent that the $\fb$-module theory developed in \cite{WvdK} would be vital in generalizing the approach to arbitrary rank. 
One way to summarize the results of {\it loc.cit.} is saying that the abelian category $\rep( \fb)$ has two highest weight category structures; recall that simple modules in $\rep(\fb)$ are parametrized by $X( \bt)$, the weight lattice of a maximal torus. 
Now the costandard objects with respect to one structure, which is defined by the {\it excellent} order on $X( \bt)$, (Definition \ref{def:excord}) are denoted $P(\lambda)$, $\lambda \in X( \bt)$ and are related to {\it excellent filtrations} on $\fb$-modules, see Section \ref{subsec:excord}.  
The costandard objects with respect to another highest weight category structure, which is defined by the {\it antipodal excellent} order on  $X( \bt)$, (Definition \ref{def:antipodal}) are denoted $Q(\lambda)$, $\lambda \in X( \bt)$ and are related to {\it relative 
Schubert filtrations} on $\fb$-modules, see Section \ref{subsec:antipodexcord}. It is these modules that will play a pivotal role in the construction of full exceptional collections on generalized flag schemes. 
We refer the reader to \cite{WvdK} for the highest weight category approach to $\rep( \fb)$; in \cite{WvdKTata} those structures are somewhat hidden in the background. The lectures \cite{WvdKTata} are devoted to the breakthrough paper \cite{MathG} by Mathieu; that breakthrough will be used in the guise of Corollary \ref{cor:PQnabla}.

In the next two paragraphs, we assume again that the group $\bg$ is of rank two. We choose a total order $\prec$ on the corresponding Weyl group that refines the Bruhat order.
For notational convenience, let $\langle X_v\rangle, v\in W$ denote the set of $\fb$-modules that underlie the full exceptional collection of $\bg$-equivariant vector bundles on $\bg/\fb$ ordered with respect to $\prec$. Recall that $\rep(\fb)$ denotes the category of finitely generated $\fb$-modules over $\sk$. A detailed inspection of the modules $\langle X_v\rangle, v\in W$ led to the two key observations that made it possible to see the pattern. The first one is that in each type in rank two the $\fb$-modules $X_v$'s have relative Schubert filtrations by modules $Q(\lambda)$, while the $\sk$-linear dual $\fb$-modules $X_v^*$'s have excellent filtrations by modules $P(\lambda)$. 
The second observation is the appearance of 
the Steinberg weights $e_v \in X( \bt)$ in weight filtrations on the modules $X_v$'s; by \cite{St}, those weights form a basis $\{\sk_{e_v}\}_{v\in W}$ (the Steinberg basis) of the representation ring $R( \fb)=\sK_0(\rep( \fb))$ as a module over the representation ring $R( \bg)$, see Section \ref{subsec:Steinberg}. 

More precisely, in the set of modules $Q(\lambda)$ that occur in the filtration of $X_v$'s, the weight $\lambda$ of the $\fb$-socle of $Q(\lambda)$ is always a Steinberg weight $e_w$ with $w\geq v$ in the Bruhat order on $W$. 
On the other hand, the  modules $P(-\lambda)^*$
that occur in the other filtration of $X_v$'s are of the form $\nabla_\mu\otimes P(-e_w)^*$ where $\nabla_\mu$ is a dual Weyl module for a  dominant weight $\mu\in X(\bt)$ and $w\leq v$ in the Bruhat order. As the corresponding coherent sheaves $\Ll(X_v), v\in W$ on $\bg/\fb$ form a full exceptional collection in rank two by Section \ref{sec:rank_2}, this imposes a semiorthogonality condition on the sheaves $\Ll(P(-e_w)^*), w\in W$ and $\Ll(Q(e_v)), v\in W$ with respect to the Bruhat order. Furthermore, the $\fb$-socle of the module $Q(e_v)$ is ${\sk }_{e_v}$, while the same character ${\sk }_{e_v}$ is also the $\fb$-head of the module $P(-e_v)^*$. These facts in rank two will serve as the basis for producing the sought-for semiorthogonal decompositions for flag varieties in arbitrary rank; it starts with Theorems \ref{th:indSteinberg_vanishing} and \ref{th:indPQ} (the reader is also invited at this point to keep consulting Section \ref{sec:Leitfaden} for an easier navigation in the key statements across the text). It is also at this stage that the interplay among three orders - one on the Weyl group (the Bruhat order) and the other two on $X(\bt)$ (the excellent and antipodal excellent orders) becomes crucial in our considerations.

%We show in Theorem \ref{th:generation} that one still gets a basis of $R(\fb)$ when replacing a few $\sk_{e_v}$'s with $Q(e_v)$ or $P(-e_v)^*$, see Corollary \ref{cor:Steinberg_alternative}.
%Section \ref{sec:B-generation}.

It should be noted that an earlier paper \cite{AAGZ} unveiled the role played  by the Steinberg basis in the setting of exceptional collections. At the same time, it showed its limitations already for the group ${\bf G}_2$ if the Steinberg basis is taken literally. As our paper explains, the crucial missing ingredient was that the Steinberg weights should have been seen through the lens of the category $\rep(\fb)$ -- that is as socles (resp., heads) of costandard (resp. standard) objects for the highest weight category structure on $\rep(\fb)$ given by the antipodal excellent order. 
The remaining parts of the costandard modules $Q(e_v), v\in W$ (resp., of   the standard modules $P(-e_v)^*, v\in W$) that lie above the socle  (resp., below the head) of their common Steinberg weight ${\sk}_{e_v}$ are responsible for the eventual semiorthogonality properties.

Now, as we have the Steinberg basis $\{\sk_{e_v}\}_{v\in W}$, and the 
$\fb$-modules $P(-e_v)^*$ and $Q(e_v)$ with $v\in W$ at our disposal for a general semisimple group $\bg$, let us take a closer look at the semiorthogonality condition on $\Ll (P(-e_w)^*)$ and $\Ll (Q(e_v))$ and explain how we arrived at those. Precisely, this is about Theorems \ref{th:indSteinberg_vanishing} and \ref{th:indPQ} (triangularity of the transition matrices at the $\sf K$-theoretical level). 
The preparatory technical work for those theorems is done in Section \ref{sec:triangular} and proceeds as follows. Let $\fb^+$ be the opposite Borel subgroup to $\fb$. The $\bt$-equivariant $\sf K$-theory of a flag variety $\bg/\fb^+$ for a Borel subgroup $\fb^+$
has two natural bases: one consisting of the classes $\{[\Oo _{X_w}],w\in W\}$ of the closures of $\fb^+$-orbits on $\bg/\fb^+$  (the ``Schubert basis''), and another consisting of the classes 
$\{[\Oo_{\fX^w}(-\partial \fX^w)]\}_{w\in W}$ of the $\fb ^{-}$-orbits on $\bg/\fb^+$ for the opposite Borel subgroup $\fb ^{-}$ to $\fb^+$ (the ``opposite Bruhat cell basis''). The characters of the modules $Q(e_v)$ and $P(-e_w)$ can be computed with Demazure operators, and with computer assisted computations in rank three we have found a close relation in the $ \bt$-equivariant $\sK$-group $\sK_{ \bt}( \bg/ \fb)$ between the classes $\{[\Ll(P(-e_w))], w\in W\}$ and the Schubert basis $\{[\Oo _{X_w}],w\in W\}$. More precisely, the $R(\bt)$-module ${\sf K}^0_{\bt}(\bg/\fb)$ is equipped with a natural $R(\bt)$-valued paring, and a result from \cite{GrahamKumar}, which is attributed to Knutson in {\it loc.cit.}, asserts that the Schubert and the opposite Bruhat cell bases are orthogonal to each other with respect to that pairing. Using this very important input, in Section \ref{sec:triangular} we show in Theorem \ref{th:triang} that the transition matrix between the basis $\{\Ll(P(-e_w))\}_{w\in W}$ and the Schubert basis of $\sK_{ \bt}( \bg/ \fb)$ is always invertible and triangular, up to a permutation of rows and columns. Similarly, the transition matrix from the basis $\{\Ll(Q(e_v))\}_{v\in W}$ to the opposite Bruhat cell basis is also invertible and triangular, up to a permutation of rows and columns. Finally, much of the cohomology vanishing required by the statement of Theorem \ref{th:indSteinberg_vanishing} is assured by the results explained in \cite{WvdKTata}, e.g.\ Theorems \ref{th:highInd} and \ref{th:highIndZ}, which are based on \cite{MathG}. These allow to reduce the proof of Theorem \ref{th:indSteinberg_vanishing} to the Euler characteristic computations, and this is precisely what Section \ref{sec:triangular} does with the help of all the input explained above in this paragraph. The triangularity results of Theorems \ref{th:indSteinberg_vanishing} and \ref{th:indPQ} turn out to be the key to the desired $\Ext$-vanishing required for exceptional collections. Arguably,  Theorems \ref{th:indSteinberg_vanishing} and \ref{th:indPQ} are the most important cohomological statements of our work.

Section \ref{sec:triangular}, discussed in the previous paragraph, concerns with the structure of the $R(\bt)$-module ${\sf K}^0(\bg/\fb)$. Similarly, the representation ring $R( \fb)=\sK_0(\rep( \fb))$ is a module over the representation ring $R(\bg)$; in both cases, the $R(\bt)$-module structure for the former module (resp., the $R(\bg)$-module structure for the latter) are essential for our purposes. More generally, the category $\rep( \fb)$ has a module structure over the base abelian category $\rep( \bg)$; that module structure is given by the restriction functor $\res_{\fb}^{ \bg}$,  cf. the aforementioned quote from Donkin. Passing to derived categories, the category $\Dd ^b(\rep( \fb))$ receives a module  structure over the base category $\Dd ^b(\rep( \bg))$ (since that module structure comes from a monoidal action at the level of abelian categories, it doesn't require further compatibilities). In this framework, thanks to Theorems \ref{th:indSteinberg_vanishing} and \ref{th:indPQ} the basic examples of Section \ref{sec:rank_2} can be reformulated as saying that for a group $\bg$ of rank two and a Borel subgroup $\fb\subset \bg$, there are semi-orthogonal decompositions of the category $\Dd ^b(\rep( \fb))$ that are linear over $\Dd ^b(\rep( \bg))$ ($\bg$-linear semiorthogonal decompositions for short, see Definition \ref{def:G-linear-SOD}) with each semiorthogonal component being equivalent to $\Dd ^b(\rep( \bg))$. 

Given a base scheme $S$, there are robust notions of $S$-linear triangulated categories and of $S$-linear semiorthogonal decompositions of a given $S$-linear triangulated category, \cite{Kuzbasech}. The notions of a $\bg$-linear triangulated category and of a $\bg$-linear semiorthogonal decomposition are in complete parallel and could be put on an equal footing with those from {\it loc.cit.} if we used the language of (quotient) stacks that we deliberately avoid. For that reason, we recall in Sections \ref{subsec:G-linear-triang_cat} and \ref{subsec:G-linear-sod} all the necessary definitions and statements concerning $\bg$-linear triangulated categories in the classical language of equivariant (quasi)-coherent sheaves.


Having set up the necessary framework, we proceed to constructing $ \bg$-linear semiorthogonal decompositions of $\Dd ^b(\rep( \fb))$ in arbitrary rank. For that we need to choose a total order $\prec$ on the Weyl group $W$ that refines the Bruhat order. For any element $p\in W$ of the Weyl group we will first construct in Section \ref{subsec:cut_at_p} a $\bg$-linear semiorthogonal decomposition of $\Dd ^b(\rep( \fb))$ (called ``cut at $p$'') into two admissible subcategories 
\begin{equation}\label{eq:p-cutSOD}
\langle\;\hull(\{\nabla_\lambda\otimes Q(e_v)\}_{v\succ p,\lambda\in X( \bt)_+}),\;\hull(\{\nabla_\lambda\otimes P(-e_v)^*\}_{v\preceq p,\lambda\in X( \bt)_+})\;\rangle,
\end{equation}

Here $\hull(-)$ denotes the triangulated envelope of a given set of objects, and $\nabla_\lambda$ are dual Weyl modules. Strikingly, the $\fb$-modules $P(-e_v)^*, v\in W$ and $Q(e_v), v\in W$ will serve for both semiorthogonality condition in (\ref{eq:p-cutSOD}), which will follow from Theorem \ref{th:indPQ} and Corollary \ref{cor:left orthogonal}, and for the generating property. Precisely, in Theorem \ref{th:derived_generation} 
we prove that the union of the two triangulated hulls in (\ref{eq:p-cutSOD}) is the whole $\Dd ^b(\rep( \fb))$. That proof, which takes Section \ref{subsec:D^brepB-generation}, can be considered as a categorical upgrade of 
\cite[Theorem 2]{Ana} which is essentially the same statement as in Theorem \ref{th:derived_generation}, but at the $\sf K$-theory level. The reader may prefer to consult \cite[Theorem 2 and Section 4]{Ana} as a starting point before proceeding to Section \ref{subsec:D^brepB-generation}. As a byproduct of the results of Section \ref{sec:B-generation}, we obtain an alternative proof and categorification of a theorem of Steinberg \cite{St}, see Theorem \ref{th:Pgeneration}, which is similar to \cite[Theorem 2]{Ana} of Ananyevskiy.

Semiorthogonal decompositions (\ref{eq:p-cutSOD}) are the core part of our paper as they allow to construct objects of $\Dd ^b(\rep( \fb))$ that will eventually give rise to full exceptional collections on $\bg/\fb$. More precisely, for an element $p\in W$ we define in Section \ref{subsec:X_p-Y_p} two objects $X_p$ and $Y_p$ of $\Dd ^b(\rep( \fb))$. Specifically, given the semiorthogonal decomposition (\ref{eq:p-cutSOD}) associated to $p\in W$, we define an object $X_p\in \Dd ^b(\rep( \fb))$ to be the image of $P(-e_p)^*$ under the left adjoint of the inclusion of the triangulated hull of $\{\nabla_\lambda\otimes Q(e_v)\}_{v\succeq  p,\lambda\in X( \bt)_+}$ into $\Dd ^b(\rep( \fb))$. Similarly, we define $Y_p$ be the image of $Q(e_p)$ under the right adjoint of the inclusion of the triangulated hull of $\{\nabla_\lambda\otimes P(-e_v)^*\}_{v\preceq p,\lambda\in X( \bt)_+}$ 
into $\Dd ^b(\rep( \fb))$. The reason of defining the objects $X_p$ and $Y_p$ as above was prompted by the insight coming from the examples in rank two: the exceptional objects $\tX_v$ on $\bg/\fb$ for rank two groups lie inside the intersection  of the hull of $\{\Ll(P(-e_v)^*)\}_{v\preceq p}$ with the hull of $\{\Ll(Q(e_v))\}_{v\succeq p}$. 

The final step of this core part is showing that for any $p\in W$ the objects $X_p$ and $Y_p$ are isomorphic. This is done in Sections \ref{subsec:X_p-Y_p-morphisms} and \ref{sec:X_p&Y_p-isom}. To this end, we need a refinement of our categorification of Steinberg's theorem, which is Theorem \ref{th:cone generation}. This refinement, implicit in \cite{Ana}, shows that the Steinberg weights serve as a curious meeting ground of the Bruhat order and the antipodal excellent order $<_a$ on the weight lattice $X(\bt)$ from Definition \ref{def:antipodal} (compare also Remark \ref{rem:Steinberg from anti}). At this stage we learn from Lemma \ref{lem:X_p-Y_p-isom} that the objects $X_p$'s we have constructed are $\fb$-exceptional, i.e. $\RHom _{\Dd^b(\rep(\fb))}(X_p,X_p)=\sk$; that is, they are exceptional in $\Dd ^b(\rep( \fb))$ considered as a $\sf k$-linear category.

But now, thanks to Theorem \ref{th:indPQ} and to the very construction of objects $X_p, p\in W$, there is much more to cohomological properties of the $X_p$'s than asserts Lemma \ref{lem:X_p-Y_p-isom}: these objects turn out to be exceptional in $\Dd ^b(\rep( \fb))$ considered as a $\bg$-linear category, see Theorem \ref{th:X_p-G-exceptional-coll} for the precise statement. Now, by Proposition \ref{prop:G-linear_exceptional}
each object $X_p$, $p\in W$ gives rise to a $ \bg$-linear functor 
$\Phi _p\colon \Dd^b(\rep ( \bg))\rightarrow \Dd^b(\rep(\fb))$. 
Cohomological vanishing statements from Section \ref{sec:b-cohom_vanish} combined with Theorem \ref{th:X_p-G-exceptional-coll} prove that each $ \bg$-linear functor $\Phi _p:\Dd^b(\rep ( \bg))\rightarrow \Dd^b(\rep(\fb))$ is full and faithful. 
Results of Sections \ref{subsec:G-linear-triang_cat} and \ref{subsec:G-linear-sod} then allow to establish that the collection of full triangulated subcategories ${\rm Im}(\Phi _p)\subset \Dd^b(\rep(\fb))$, $p\in W$, forms a $ \bg$-linear semiorthogonal decomposition of the 
category $\D= \Dd^b(\rep(\fb))$ with respect to the {chosen} total order $\prec$ on $W$. 
This is the statement of Theorem \ref{th:semi-orthogonal}.


Now the proof of Theorem \ref{th:coh-semi-orthogonal} goes as follows.
We define objects $\tX_p\in \Dd^b(\Coh ( \bg/ \fb))$ as the images of $X_p$, $p\in W$ from Section \ref{sec:X_p&Y_p} under the composition of the sheafification functor $\Ll$ with the forgetful functor ${\sf For}:{\Dd ^b}(\Coh ^{ \bg}( \bg/ \fb))\rightarrow \Dd^b(\Coh ( \bg/ \fb))$. 
Theorem \ref{th:semi-orthogonal} implies that the collection of objects $\tX_p, p\in W$ is an exceptional collection in $\Dd^b(\Coh ( \bg/ \fb))$, while the generating property of $\tX_p, p\in W$ is assured by the results of Section \ref{sec:B-generation}. Remark \ref{rem:basechangeG-to-point} draws a parallel of Theorem \ref{th:coh-semi-orthogonal} with a base change type of statement for semiorthogonal decompositions, \cite{Kuzbasech}.

Section \ref{sec:caseP} treats the case of a parabolic subgroup $\bp\supset \fb$. Its main statements, which are Theorems \ref{th:Psemi-orthogonal} and \ref{th:Pcoh-semi-orthogonal}, follow essentially the same path that has been set out in Theorems \ref{th:semi-orthogonal} and \ref{th:coh-semi-orthogonal}. For the reasons that are explained below in this paragraph, it is natural to expect that the objects ${\hat{\sf X}} _v, {v\in W^\bp}$ that give rise to semiorthogonal decompositions of $\Dd ^b(\rep(\bp))$ as a $\bg$-linear category (resp., the objects $\hat \tX_v\in \Dd ^b(\Coh(\bg/\bp), v\in W^\bp$ giving full exceptional collections in $\Dd ^b(\Coh(\bg/\bp)$) are contained among the objects ${\sf X} _v, {v\in W}$ of Theorem \ref{th:semi-orthogonal} (resp., among the objects $\tX_v, v\in W$ of $\Dd ^b(\Coh(\bg/\fb)$ of Theorem \ref{th:Pcoh-semi-orthogonal}). One has therefore to recognize those objects among ${\sf X} _v, {v\in W}$ (resp., among $\tX_v, v\in W$) that are obtained by the restriction functor $\res_{\bg}^{\bp}:\Dd ^b(\rep(\bp))\rightarrow \Dd ^b(\rep(\fb))$ (resp., by the pullback $\pi _{\bp}^{\ast}:\Dd ^b(\Coh(\bg/\bp)\rightarrow \Dd ^b(\Coh(\bg/\fb)$ along the projection $\pi _{\bp}:\bg/ \fb\rightarrow  \bg/ \bp$). The fundamental fact that both functors $\res_{\bg}^{\bp}$ and $\pi _{\bp}^{\ast}$ are $t$-exact and fully faithful on the respective derived categories makes it possible to recognize the sought-for exceptional objects on $\bg/ \bp$ by applying the induction functor $\Rind _{\fb}^{\bp}$ to 
appropriate objects ${\sf X} _v, {v\in W}$ (resp., the pushforward $R{\pi _{\bp}}_{\ast}$ to $\tX_v, v\in W$). It turns out the Steinberg weights $e_v, v\in W^\bp$ for a given parabolic $\bp$ behave nicely with respect to the induction functor $\Rind _{\fb}^{\bp}$ suggesting a natural parabolic analogue of the key $\fb$-modules from Section \ref{sec:rappelB-mod}. The cohomological properties of those parabolic analogues are given by Theorem \ref{th:PindPQ}, a parabolic counterpart of Theorem \ref{th:indPQ}. That allows to further apply the arguments of Sections \ref{sec:B-generation} and \ref{sec:X_p&Y_p} in the parabolic case obtaining Theorems \ref{th:Psemi-orthogonal} and \ref{th:Pcoh-semi-orthogonal}.


%--------------------------------------------------------------
\subsection{Discussion}\label{sec:disc}
%--------------------------------------------------------------
The forgetful functor $\rep( \fb)\rightarrow {\Coh}( \bg/ \fb)$ underlines Kuznetsov--Polishchuk's construction of exceptional collections on maximal Grassmannians, \cite{KuzPol}.
Starting off from a full exceptional collection in the equivariant category and suitably mutating the collection, they arrive at an exceptional collection of objects ${\mathcal E}^{\lambda}$ in the non-equivariant derived category where $\lambda$ runs over a certain collection of weights in $X( \bt)$. 
It is an interesting question to understand the relation of the objects ${\mathcal E}^{\lambda}$ from {\it loc.cit.} to the the objects $\hat \tX_v, v\in W^{\bp}$ from Theorem \ref{th:Pcoh-semi-orthogonal}. 

Furthermore, the objects ${\mathcal E}^{\lambda}$ of \cite{KuzPol} are {\it a priori} complexes of coherent sheaves. 
Quoting from {\it loc.cit.}: ``Note that a priori the exceptional objects ${\mathcal E}^{\lambda}$  constructed above are complexes. 
However, we have the following Conjecture 4.1: For any exceptional block $B\subset P_L^+$ and $\lambda \in B$ the object ${\mathcal E}^{\lambda}$ is a vector bundle.'' 
At the moment, it is not clear whether the objects $\tX_p$, $p\in W$ from Theorem \ref{th:coh-semi-orthogonal} (resp., $\hat \tX_v, v\in W^{\bp}$ from Theorem \ref{th:Pcoh-semi-orthogonal}) are pure with respect to the standard $t$-structure on $\Dd^b(\Coh ( \bg/ \fb))$ (resp., on  $\Dd^b(\Coh ( \bg/ \bp))$). Evidence for their purity comes from the rank two cases, see Section \ref{sec:rank_2}. In Section \ref{sec:dualFEC_conj}, based on further evidence coming from low rank cases in ranks up to three, we state some conjectural cohomological statements that shall eventually give the description of dual exceptional collections to $\tX_v, v\in W$ in $\Dd ^b(\Coh(\bg/\fb)$. If true, that description must provide further evidence for purity of the objects $\tX_v, v\in W$ up to a shift.

Presumably, Theorems \ref{th:semi-orthogonal} and \ref{th:Psemi-orthogonal} have further extensions to larger settings than the one of 
semisimple algebraic groups of our paper. Much of the representation-theoretic input that we used in proving Theorems \ref{th:semi-orthogonal} and \ref{th:Pcoh-semi-orthogonal} also exists in the quantum setting, \cite{APW}. This suggest natural generalizations
of the said theorems to the quantum case. Representation theory of parahoric subalgebras of affine Kac-Moody Lie algebras studied in \cite{FKM} suggests another extension.

A very interesting question is to understand a relation of our results to baric structures/staggered $t$-structures on the derived categories of equivariant coherent sheaves from \cite{Achar} and \cite{AchTr}. 
Further, full exceptional collections on flag varieties appear through the computations related to the dual Steinberg basis of ${\sf K}^0_{ \bg}( \bg/ \fb)$, \cite{Daw}. 
These, in turn, are related to Lusztig's asymptotic affine Hecke algebra and to Lusztig's canonical basis in the $ \bt$-equivariant $\sf K$-theory of Springer fibres. 
Ultimately, the problem is to relate the objects $X_p$, $p\in W$ of the present paper to Lusztig's canonical basis. 



\subsection{Leitfaden}\label{sec:Leitfaden}

\vspace{2cm}
\xymatrix @R=2pc @C=-2.5pc{*+[F]\txt{\bf Corollary \ref{cor:PQnabla} \\ 
($\fb$-cohomology vanishing \\ for $P(\lambda)\otimes Q(\mu)\otimes \nabla_\nu$)} \ar[d]  \ar[drr]
 & &  *+[F]\txt{\rm {\bf Theorem \ref{th:triang}} (Triangularity)} \ar[d]   &  \\
 *+[F]\txt{\bf Theorem \ref{th:highInd}\\
($\RR^i\ind_{\fb}^ \bg (P(\lambda)\otimes Q(\mu))=0$ \\ for $i>0$)} \ar[drr]  
& &  *+[F]\txt{{\bf Theorem \ref{th:indSteinberg_vanishing} (\ref{th:indPQ})} \\ 
(Computing $\ind_{\fb}^ \bg (P(-e_v)\otimes Q(e_w))$ \\ 
 for $w\succeq v$)}  \ar[d]  &  \\
&  *+[F]\txt{\bf Theorem \ref{th:generation} \\
(Generation)} \ar[r] \ar@{.>}[ddl]^{Refined \ version \ of \ generation} &  *+[F]\txt{{\bf Section \ref{subsec:cut_at_p}} (``Cut at $p\in W$'')} \ar[dd]  \\
& & \\
 *+[F]\txt{\bf Section \ref{sec:X_p&Y_p-isom}\\ 
(Vanishing of the cone of $X_p\rightarrow Y_p$)
} \ar[ddr]        
&  & *+[F]\txt{\bf Section \ref{subsec:X_p-Y_p} \\
(Construction of $X_p$'s and $Y_p$'s)} \ar[ddl] \ar[ll] &  \\
\\
&   *+[F]\txt{\bf Section \ref{sec:SODofB-modoverG-mod} \\ 
(Constructing $\bf G$-semi-orthogonal \\\ decompositions of $\Dd ^b(\rep(\fb)$) \\ with the help of objects $X_v,v\in W$} \ar[dd]  & \\
\\
&   *+[F]\txt{\bf Sections \ref{sec:FEConGmodB} and \ref{sec:caseP} \\ 
(Full exceptional collections \\ on ${\mathbb G/\mathbb P}$)}  &
 }

\

%--------------------------------------------------------------
\subsection*{Acknowledgements}
%--------------------------------------------------------------
A.S. is grateful to Henning Krause for his unwavering support of this work and for his and Charles Vial's joint invitation to the BIREP group of the University of Bielefeld for a long-term visit in 2023-2025. 
The research of A.S. was supported by the Deutsche Forschungsgemeinschaft (SFB-TRR 358/1 2023 - 491392403). He also acknowledges the hospitality of the Mathematical Institute of the University of Utrecht during his visits in May 2024, and in March and May 2025. 

W.v.d.K. thanks BIREP for the opportunity to visit Bielefeld in September 2023.

%--------------------------------------------------------------
\subsection{Generalities}\label{subsec:Notation}
%--------------------------------------------------------------

%--------------------------------------------------------------
\subsubsection{Semisimple groups}\label{subsec:reductive}
%--------------------------------------------------------------
Let ${\mathbb G}\rightarrow \mathbb Z$ be a split simply connected semisimple Chevalley group scheme (a smooth split affine group scheme over ${\rm Spec}(\mathbb Z)$ whose geometric fibres are connected simply connected
semisimple algebraic groups), and ${\mathbb G}/{\mathbb B}\rightarrow \mathbb Z$ be the corresponding Chevalley flag scheme (resp., 
the corresponding
generalized flag scheme ${\Gm}/{  \Pp}\rightarrow \Z$ for a standard parabolic subgroup scheme ${\mathbb P}\subset \mathbb G$ over ${\mathbb Z}$). 
Then ${\mathbb G/\mathbb P}\rightarrow {\rm Spec}({\mathbb Z})$ is flat and any line bundle $\mathbb L$ on ${\Gm/  \Pp}$ also comes from a line bundle  on ${\mathbb G/\mathbb B}$. 

\begin{remark}
    For the proofs we require that ${\mathbb G}$ is simply connected.
    But that is fine, because otherwise one just views the generalized flag scheme as a homogeneous space for the simply connected cover of ${\mathbb G}$.
\end{remark}

Given an affine algebraic group $\bh$, flat over a noetherian base ring $\sk$, an $\bh$-module is a $\sk[\bh]$-comodule (a rational representation). 
The abelian category of $\bh$-modules is denoted $\Rep(\bh)$ and $\rep(\bh)$  denotes the full abelian subcategory of representations which are finitely generated over the base ring $\sk$. (cf. \cite[Section 1]{AndChev}).

\begin{lemma}(\cite[I 2.7(4), I 4.2(1), I 4.4]{Jan})
    Let $M$, $N\in \Rep(\bh)$ with $M$
    finitely generated and projective over $\sk$. Then $M^*=\Hom_\sk(M,\sk)$ is an $\bh$-module and 
    $$\Ext^n_\bh(M,N)=H^n(\bh,M^*\otimes N).$$
\end{lemma}


\begin{theorem}[Universal coefficient Theorem]\label{th:universal coefficients}
Let $\sk$ be a Dedekind ring and let $\bg$ be an affine algebraic group, flat over $\sk$. Let $\sk'$ be commutative $\sk$-algebra and $n\geq0$.
\begin{enumerate}
    \item There is for any $\bg$-module $M$, flat over $\sk$, an exact sequence
    $$0\to H^n(\bg,M)\otimes \sk'\to H^n(\bg_{\sk'},M)\to \Tor_1^\sk(H^{n+1}(G,M),\sk')\to 0.$$
    \item Let  $X$ be a closed reduced
    $\Bm_\sk$-invariant subscheme, flat over $\sk$, of the generalized flag variety $(\Gm/\Pp)_\sk$. For any  $\Pp_\sk$-module $N$, finitely generated and flat over $\sk$, we have an exact 
    sequence 
    \end{enumerate}
    $$0\to H^n(X,\Ll(N))\otimes \sk'\to H^n(X_{\sk'},\Ll(N))\to \Tor_1^\sk(H^{n+1}(X,\Ll(N),\sk')\to 0,$$
where $\Ll(N)$ is the vector bundle associated to $N$ 
    (\cite[I 5.8]{Jan}).
\end{theorem}
\begin{proof}
        The first part is \cite[Proposition I 4.18]{Jan}.
        The second part is a very special case,  because by \cite[I 5.8(2)]{Jan} these sheaf cohomologies may be computed through algebraic group cohomology.
        The second part is tailored to
        our needs. One may also prove the second part with \v Cech cohomology.
        Of course we allow $\Pp=\Bm$.
    \end{proof}

    \

    
Our base ring $\sk$ is usually a field or $\Z$. We will discuss the case $\sk=\Z$ separately, or leave the necessary modifications to the reader.
For instance, when $M$ is a finitely generated free module over $\sk=\Z$, then $\dim_\sk(M)$ means its rank. 
In our references the field is often assumed algebraically closed, but this assumption is irrelevant in our context because of
the straightforward behaviour of the various cohomology groups under field extensions of $\sk$.


Let  $\bg$  be  obtained from ${\mathbb G}$ by base change along ${\rm Spec}(\sk)\rightarrow {\rm Spec}(\mathbb Z)$.
When $\sk=\Z$ this just means that $\bg=\Gm$. 

Now let $\sk$ be a field of arbitrary characteristic.
So $ \bg$ is a split semisimple simply connected affine algebraic group. 
Let $ \bt$ be a split maximal torus of $ \bg$, obtained by base change from a split maximal torus $\Tm$ in $\Bm$.
Let $W = N_{ \bg}({\bt})/ \bt$ be the Weyl group and let $\fb\supseteq \bt$ correspond with $\Bm\supseteq \Tm$. 
Let $X = X( \bt)$ be the weight lattice and $\Phi$ be the root system. We choose the system $\Phi ^+$ of positive roots which makes $\fb$ the negative Borel and let $\Pi$ be the set of simple roots. 
Let $(~,~)$ be a non-singular, symmetric, positive definite $W$-invariant form on ${\mathbb R}\otimes_{\mathbb Z} X$ and let $X(\bt)_+=\{\lambda \in X: (\lambda,\alpha ^{\vee})\geq 0$ for all $\alpha \in \Pi$\}, be
the set of dominant weights, where $\alpha^\vee =\frac{2\alpha}{(\alpha ,\alpha)}$ for $\alpha \in \Phi$.\\

We want to emphasize once again the fact that our convention (which is that of \cite{Jan}) is that ${\fb}$ corresponds 
to the {\it negative roots} (as opposed to the conventions of \cite{GrahamKumar}, \cite{MathG}, \cite{Polo3}, and \cite{WvdK}). 
For that reason, the translation of various results from {\it loc.cit.} that are used extensively in this paper requires some care.\\ 

Let $\Pi=\{\alpha _1,\dots ,\alpha _r\}$ be the set of simple roots relative to our choice of $\Phi ^+$. The fundamental weights $\omega_i$ satisfy $(\omega_i,\alpha_j^\vee)=\delta_{ij}$.
A parabolic subgroup of $ \bg$ is usually denoted by $ \bp$; in particular, for a simple root $\alpha \in \Pi$, denote ${ \bp}_{\alpha}$ the minimal parabolic subgroup of $ \bg$ associated to $\alpha$. 

Given a $\fb$-module $M$ we associate to it as in \cite[I 5.8]{Jan} a sheaf $\Ll(M)$ on $ \bg/ \fb$ with fibre $M$ at the point $ \fb/ \fb$. Given a weight $\lambda \in  {\rm X}({\bt})$, denote $\Ll _{\lambda}$ the corresponding line bundle on ${ \bg}/{\fb}$.
Here we still follow  \cite{Jan}, so that $\Ll_\lambda$ or $\Ll(\lambda)$ is associated to the one dimensional $\fb$-module $\sk_\lambda$ of weight $\lambda$, not associated to the dual of that representation.\\

The weight lattice $X$ has a natural partial order $\geq_d$, known as the dominance order: 
for $\lambda,\mu \in X$ we write $\lambda \geq_d \mu$ if $\lambda - \mu$ is a sum of positive roots, with repetitions allowed\footnote{That is, $\lambda -\mu$ is of the form
$\sum _{\alpha\in \Phi ^+}n_{\alpha}\alpha$, for non-negative integers $n_{\alpha}$.}. 
We denote $w_0$ the longest element of the Weyl group $W$ and let $\lambda ^{\ast} = -w_0\lambda$ for $\lambda \in X(\bt)_+$ (dual or contragredient, cf. \cite[Section 2.2]{Hummodular}).\\
 
For $\lambda \in X$ let %${\sk}_{\lambda}$ be the one dimensional %$\fb$-module on which $ \bt$ acts via $\lambda$ and let 
$\nabla _{\lambda}$ or $\nabla(\lambda)$ be the induced module $\ind_{\fb}^{ \bg}({\sk}_{\lambda})$ \cite[I, Section 3.3]{Jan}. 
It is finite-dimensional and non-zero if and only if $\lambda$ is dominant. 
For $\lambda \in X(\bt)_+$ we denote by $\Delta (\lambda)$ the Weyl module $(\nabla _{\lambda^{\ast}})^{\ast}$. Then $\nabla _{\lambda}$  has simple socle $L(\lambda)$ and $\Delta (\lambda)$ has simple head $L(\lambda)$. 
If $\sk=\Z$ then  $\nabla _{\lambda}$, $\Delta (\lambda)$ are finitely generated and free over $\sk$, but heads and socles make less sense.\\

\begin{proposition}\label{prop:cohinduced}(\cite[II Proposition 4.13, B.4]{Jan})
Let $\lambda\in X(\bt)_+$. 
\begin{equation*}
\Ext ^i_{\bg}(\sk,\nabla(\lambda)) = \Ext ^i_{\fb}(\sk,\sk_\lambda)=\left\{
   \begin{array}{l}
    {\sf k}, \quad {\rm for} \quad i=0, \  \lambda=0,\\
    0, \quad {\rm otherwise}. \\
   \end{array}
  \right.
\end{equation*}
\end{proposition}
%--------------------------------------------------------------
\subsubsection{Triangulated categories}\label{subsec:triangulated}
%--------------------------------------------------------------

Let $\sk$ be a field. Given a $\sk$-linear triangulated category $\D$, equipped with a shift functor $[1]\colon \D \rightarrow \D$ and two
objects $A,B\in \D$, we let denote $\Hom ^{\bullet}_{\D}(A,B)$ 
the graded $\sk$-vector space $\bigoplus _{i\in \mathbb Z}{\Hom} _{\D}(A,B[i])$. A full triangulated subcategory $\A \subset \D$ is a full subcategory which is closed under shifts and taking cones.\\
A strictly full subcategory is a full subcategory that contains with every object also the objects isomorphic to it.

Given an algebraic variety $X$ over a field $\sk$ (an integral scheme \dots ) denote 
$\Dd _{qc}(X)$ (resp., $\Dd _{qc}^+(X)$, $\Dd _{qc}^-(X)$) the unbounded derived category (resp., the bounded below, the bounded above) of quasi-coherent sheaves and $\Perf(X)$ the category of perfect complexes on $X$. 
The bounded derived category of coherent sheaves on $X$ is denoted $\Dd ^b(X)$ (it coincides with $\Perf(X)$ when $X$ is smooth).
Given a morphism $f:X\rightarrow Y$ between two varieties (integral schemes \dots), we write $f_{\ast}$ (resp, $f^{\ast}$) for the corresponding derived functors of push--forwards and pull--backs between $\Dd _{qc}(X)$ and $\Dd _{qc}(Y)$,
 and ${\rm R}^if_{\ast}$, ${\rm L}^if^{\ast}$ for corresponding abelian functors between the abelian categories of coherent sheaves. 
 We write $(-)\otimes(-)$ for the derived functor of tensor product $\Dd _{qc}(X)\times \Dd _{qc}(X)\rightarrow \Dd _{qc}(X)$. We use $\RHom _X(-,-)$  (resp., 
$R{\mathcal Hom}_X(-,-)$) for the right derived functor of $\Hom_X(-,-)$ (resp., 
local ${\mathcal Hom}$) and $\RGamma (X,-)$ for the right derived functor of global sections $\Gamma (X,-)$.

%--------------------------------------------------------------
\section{Rappels: $\fb$-modules}\label{sec:rappelB-mod}
%--------------------------------------------------------------

%--------------------------------------------------------------
\subsection{Highest weight category structure on $\Rep( \bg)$}\label{subsec:hwc_repG}
%--------------------------------------------------------------
Let $\sk$ be a field.

Recall that relative to the choice of the system of positive roots $\Phi ^+$ there is a partial order $\leq_d$ on $X( \bt)$: 
given two weights $\lambda,\mu \in X( \bt)$ we write $\lambda \leq_d\mu$ if $\mu -\lambda $ is a nonnegative integer combination of elements of $\Phi ^+$.
Simple modules in $\Rep( \bg)$ are parametrized by the dominant weights $X(\bt)_+$, \cite[II, Section 2]{Jan}. 
By \cite{CPSHW} and \cite{Donkin1}, there is a highest weight category structure on $\Rep( \bg)$ with the weight poset 
$(X(\bt)_+,\leq_d)$: the standard modules in this structure are the Weyl modules $\Delta _{\lambda},\lambda \in X(\bt)_+$ and the costandard modules are the dual Weyl modules $\nabla _{\lambda},\lambda \in X(\bt)_+$.

%--------------------------------------------------------------
\subsection{Joseph-Demazure modules}\label{subsec:J-D}
%--------------------------------------------------------------

Let $ \bh\subset  \bg$ be a flat subgroup scheme of $ \bg$. 
Associated to it are the restriction and induction functors ${\res}_{ \bh}^{ \bg}$ and ${\ind}_{ \bh}^{ \bg}$, \cite[I, Section 3.1]{Jan}. Recall the definition of Joseph's functors ${\sH}_w, w\in W$, following \cite{WvdK}. 


\begin{definition}\label{def:Joseph_functor}
Let $M$ be a $\fb$-module and  $w\in W$. 
Consider the Schubert variety $\fX_w=\overline{\fb w\fb}/\fb$ associated to $w$. Then the functor ${\sH}_w:{\Rep}( \fb)\rightarrow {\Rep}( \fb)$ is given by $M\rightarrow {\rm H}^0(\fX_w,\Ll (M))$. 
\end{definition}

Equivalently, by \cite[Proposition 2.2.5]{WvdKTata}, the functor ${\sH}_w$ can be described as follows.
If $\alpha\in \Pi$ with corresponding simple reflection $s$, let $\bp_s$ or $\bp_\alpha$ denote the minimal parabolic subgroup generated by $\fb$ and ${ \bu}_{\alpha}$. 
Then $\sH_s={\res}_{\fb}^{{ \bp}_{s}}{\ind}_{\fb}^{{ \bp}_{s}}$. And if  $s_{1}s_{2}\dots s_{n}$ is   a reduced expression for $w$, then
$${\sH}_w:=\sH_{s_1}\circ\dots\circ \sH_{s_n}.$$
 

The functor ${\sH}_w$ is left exact; let ${\rm R}{\sH}_w$ denote the right derived functor $\Dd ^+({\Rep}( \fb))\rightarrow \Dd ^+({\Rep}( \fb))$. It restricts to a functor between bounded derived categories  $\Dd ^b({\rep}( \fb))\rightarrow \Dd ^b({\rep}( \fb))$ (see Section \ref{subsec:G-linear-triang_cat}).

\begin{definition}\label{def:dualJosephmod}(Dual Joseph modules $P(\lambda)$).
Let $\lambda \in X( \bt)$. Let $\lambda^+$ be the dominant weight in the Weyl group orbit of $\lambda$ and let $w$ be minimal so that $\lambda=w\lambda^+$. The dual Joseph module $P(\lambda)$ is set to be ${\sH}_w(\lambda ^+)$. Its $\fb$-socle is of weight $\lambda$ \cite[Lemma 2.2.9]{WvdKTata}.
\end{definition}

%--------------------------------------------------------------
\subsection{Relative Schubert modules}\label{subsec:Schubfilt}
%--------------------------------------------------------------

%--------------------------------------------------------------
\begin{definition}[Relative Schubert modules $Q(\lambda)$]\label{def:defineQ}
%--------------------------------------------------------------

Let $\lambda=w\lambda^+$ as above. 
The relative Schubert module $Q(\lambda)$ is set to be the kernel of the (surjective) restriction homomorphism $P(\lambda)\rightarrow {\rm H}^0(\partial \fX_w,\Ll (\lambda ^+))$, where the boundary $\partial \fX_w$ is the union of the $\fX_v$ that are strictly contained in $\fX_w$.
\end{definition}

%\begin{remark}
 %   The  notation $X_v$ is also used for certain objects of $\Dd^b (\rep ({\fb}))$. Confusion with Schubert varieties seems unlikely.
%\end{remark}
%--------------------------------------------------------------
\subsection{$\fb$-cohomological duality}\label{subsec:key_orthogonality}
%--------------------------------------------------------------

Let $P(\lambda)$ be the dual Joseph module with socle of weight $\lambda$. The modules $Q(\mu)$ (minimal relative Schubert module with socle $\sk_{\mu}$) satisfy:
\begin{equation}\label{eq:P-Q-duality}
\Ext ^i_{\fb}(P(\lambda)^{\ast},Q(\mu)) = \left\{
   \begin{array}{l}
    {\sf k}, \quad {\rm for} \quad i=0, \  \lambda=-\mu,\\
    0, \quad {\rm otherwise}. \\
   \end{array}
  \right.
\end{equation}
See also Theorem \ref{th:Tata Theorem 3.2.6}.
%--------------------------------------------------------------
\subsection{Excellent order}\label{subsec:excord}
%--------------------------------------------------------------

Recall that we have fixed a Weyl group invariant inner product $(~,~)$ on $X({\bt})\otimes _{\mathbb Z}\mathbb R$.

\begin{definition}(\cite{WvdK})\label{def:excord}
Let $\lambda, \mu \in X({\bt})$. Define that $\lambda$ is less than $\mu$ in the excellent order, notation $\lambda \leq_e \mu$, if either $(\lambda, \lambda)<(\mu,\mu)$ or $\lambda = w\nu$, $\mu=z\nu$ for some $\nu \in X({\bt})_+$, $w,z\in W$ with $w\leq z$ (in the Bruhat order on $W$).
\end{definition}


%--------------------------------------------------------------
\subsection{Antipodal excellent order}\label{subsec:antipodexcord}
%--------------------------------------------------------------

\begin{definition}\label{def:antipodal}(\cite{WvdK})
Define that $\lambda \leq_a \mu$ in the antipodal excellent order if $-\lambda \leq_e -\mu$ in the excellent order.
\end{definition}

%--------------------------------------------------------------
\subsection{Two highest weight category structures on $\Rep( \fb)$}\label{subsec:hwc_repB}
%--------------------------------------------------------------

\begin{theorem}(\cite[Theorem 1.6, (i)]{WvdK})\label{th:Pexc}
The category  ${\mathrm{Rep} }( \fb)$ of rational representations  of $ \fb
$ is a highest weight category with respect to the excellent order.
The $P(\lambda)$ are the costandard modules
for this order.
\end{theorem}


\begin{theorem}(\cite[Theorem 1.6, (ii)]{WvdK})
The category  ${\mathrm{Rep} }( \fb)$ of rational representations  of $\fb$ is a highest weight category with respect to the antipodal excellent order.
The $Q(\lambda)$ are the costandard modules
for this order.
\end{theorem}

\begin{remark}\label{rem:anti-standard}
The terminology costandard is not used in \cite{WvdK}.  In the usual terminology nowadays, cf.\ \cite[\S 2]{Donsur}, $P(\lambda)$ is the costandard module for $<_e$ with socle $\sk_\lambda$ and $Q(\lambda)$ is the costandard module for $<_a$ with socle $\sk_\lambda$. Dually, $P(\lambda)^*$ is the standard module for $<_a$ with head $\sk_\lambda^*=\sk_{-\lambda}$ and $Q(\lambda)^*$ is the standard module for $<_e$ with head $\sk_\lambda^*=\sk_{-\lambda}$.
So equation (\ref{eq:P-Q-duality}) is about the interaction of standard and costandard modules for $<_a$, as in \cite[Proposition 3.12]{Ef}.
\end{remark}

%--------------------------------------------------------------
\subsection{Cohomology vanishing for $\fb$-modules}
%------------------------------------------------------

\begin{theorem}(\cite[Theorem 2.20(i)]{WvdK}, \cite[Theorem 3.2.6]{WvdKTata})\label{th:Tata Theorem 3.2.6}
Let $\lambda, \mu\in X({\bt})$. Then $H^p( \fb,P(\lambda)\otimes Q(\mu))=0$ for $p>0$.
\end{theorem}

This will be subsumed by Corollary  \ref{cor:PQnabla}.
%--------------------------------------------------------------
\subsection{Filtrations on $\fb$-modules}
%--------------------------------------------------------------

We recall here the fundamental statements concerning filtrations on $\fb$-modules. 
As we emphasized in Section \ref{subsec:reductive}, one has to take some care when citing the literature, 
because our $\fb$ is not $ \fb^+$. We have to convert to our conventions.
That means that dominant often becomes antidominant and vice versa. 

We begin with recalling the definition of good filtration on a $ \bg$-module, \cite{Donkin1}.

\begin{definition}
A rational $ \bg$-module $M$ is said to admit a good filtration provided that there exists a finite or infinite increasing filtration
$$
0=M_0\subsetq M_1\subsetq M_2\subsetq \dots \subsetq M_n\subsetq \dots ; \bigcup _{i>0} M_i = M
$$
such that each $M_i/M_{i-1}=\nabla _{\lambda_i}$ for some $\lambda_i \in X(\bt)_+$.
\end{definition}
\begin{aside}
An unfortunate side effect of this definition is that if $M$ admits a good filtration, its dimension is at most countable, even if $\bg$ acts trivially on $M$. 
This makes results like Theorem \ref{th:G criterion} a bit complicated.
If one wants to keep representation theory clean, it is better to use \cite[Definition 4.1.1]{WvdKTata}, 
which is similar to the next definition.
\end{aside}

\begin{definition}\cite[Definition 2.3.6]{WvdKTata}
A $\fb$-module $M$ is said to have an excellent filtration if and only if there exists a filtration $0\subset F_0\subset F_1\subset \dots $ by $\fb$-modules such that $\cup F_i= M$ and $F_i/F_{i-1}=\oplus P(\lambda _i)$ for some $\lambda _i\in X( \bt)$. Here $\oplus$ stands for any number of copies, ranging from zero copies to infinitely many.
\end{definition}

\begin{remark}
Modules possessing excellent filtrations are called {\it strong} modules in \cite{Math}.
\end{remark}

\begin{definition}\cite[Definition 2.3.8]{WvdKTata}
A $\fb$-module $M$ is said to have a relative Schubert filtration if and only if there exists a filtration $0\subset F_0\subset F_1\subset \dots $ by $\fb$-modules such that $\cup F_i= M$ and $F_i/F_{i-1}=\oplus Q(\lambda _i)$ for some $\lambda _i\in X( \bt)$. Here $\oplus$ stands for any number of copies, ranging from zero copies to infinitely many.
\end{definition}

\begin{remark}
Modules possessing filtrations by relative Schubert modules are called {\it weak} modules in \cite{Math}.
\end{remark}
\begin{definition}
A weight $\lambda$ of a $\fb$-module $M$ is called \emph{extremal} if all other weights $\mu$
of $M$ are shorter or equally long:
$(\mu,\mu)\leq(\lambda,\lambda)$.   
\end{definition}

\begin{proposition}\cite[Proposition 2.3.11]{WvdKTata}\label{prop:multone}
For any $\fb$-invariant closed subset $S\subset   \bg / \fb$ and $\lambda \in X( \bt)_+$, the module $H^0(S,\Ll (\lambda))$ has a relative Schubert filtration with layers $Q(w\lambda)$. Moreover, $Q(w\lambda)$ occurs if and only if $w\lambda$ is an extremal weight of $H^0(S,\Ll (\lambda))$ and then has multiplicity one. 
\end{proposition}

\begin{corollary}\cite[Corollary 2.3.12]{WvdKTata}
The modules ${\sH}_w(\lambda)$ have relative Schubert filtration for all $w\in W$ and $\lambda \in X( \bt)_+$.
\end{corollary} 

\begin{theorem}(\cite[Corollary 5.2.3]{WvdKTata}, 
\cite[\S5, Corollary 2]{MathG})\label{th:lambda_tensor_Q-excellent}
Let $\lambda \in X( \bt)_+$ be a regular dominant weight and $\mu \in X( \bt)$ be an arbitrary weight. Then $\lambda \otimes Q(\mu)$ has excellent filtration.
\end{theorem}

\begin{theorem}\cite[Corollary 5.2.7]{WvdKTata}\label{th:Joseph_conj}
Let $\lambda \in X({\bt})$ and $\mu \in X({\bt})_+$. Then $P(\lambda)\otimes \nabla_\mu$ has excellent filtration. Note that $\nabla_\mu=P(w_0\mu)$.
\end{theorem}
\begin{proof} This relies on the main results of \cite{MathG}.
   For completeness we recall the argument. Let $\lambda^+$ be the dominant weight in the Weyl group orbit of $\lambda$ and let $w$ be minimal so that $\lambda=w\lambda^+$. By repeated application of  the  Tensor Identity \cite[I Proposition 4.8]{Jan}, we have $P(\lambda)\otimes \nabla_\mu={\sH}_w(\sk_{\lambda^+}\otimes \nabla_\mu)$. Now $k_{\lambda^+}\otimes \nabla_\mu$
   has excellent filtration by \cite[\S 5, Corollary 1]{MathG}.
   The result in the Theorem thus follows from \cite[Lemma 2.11]{WvdK}.
\end{proof}

The next result is a key ingredient in our proofs. 

\begin{corollary}\label{cor:PQnabla}
    Let $\lambda,\mu\in X( \bt)$, $\nu\in X( \bt)_+$. Then $H^p( \fb,P(\lambda)\otimes Q(\mu)\otimes \nabla_\nu)=0$ for $p>0$.
\end{corollary}
\begin{proof}
Combine Theorem \ref{th:Tata Theorem 3.2.6}
with Theorem \ref{th:Joseph_conj}.
\end{proof}

\begin{remark}
The vanishing statement of Corollary \ref{cor:PQnabla} exhibits a remarkable interplay among the three highest weight category structures: one on the category ${\mathrm{Rep} }( \bg)$ from Section \ref{subsec:hwc_repG}, 
and the other two on ${\mathrm{Rep} }( \fb)$ from  Section \ref{subsec:hwc_repB}.
Note also that $\sk = P(0)=Q(0)=\nabla(0)$ is costandard in all three highest weight category structures.
\end{remark}
\comment{
For completeness we mention

\begin{theorem}\cite[Theorem 2]{Math}\label{th:Mathieu_1st_th}
Let $\sk$ be a field of characteristic zero.
Let $n$ be an integer and $M_1,\dots M_n$ be $n$ weak modules (i.e.\ each $M_i$ has a filtration by $Q(\lambda _i)$'s, $i=1,\dots ,n)$. Suppose $\lambda -n\rho$ is dominant. Then the $\fb$-module $M_1\otimes \dots \otimes M_n\otimes {\sf k}_{\lambda}$ is strong (i.e.\ has a filtration by dual Joseph modules $P(\lambda)$).
\end{theorem}
}
%--------------------------------------------------------------
\subsection{Cohomological criteria}
%-------------------------------------------------------------

\begin{theorem}\cite[Theorem 1.1']{Donkin1}
If $\lambda \in X(\bt)_+$ and $M$ is a $ \bg$-module with a good  
filtration then ${\rm H}^i( \bg,\nabla _{\lambda}\otimes M)=0$ for $i>0$.
\end{theorem}

\begin{theorem}\cite[Theorem 4, Corollary 7]{Frid}\label{th:G criterion}
Let $M$ be a finitely/countably generated rational $ \bg$-module satisfying ${\rm H}^1( \bg,\nabla _{\mu}\otimes M)=0$ for all $\mu \in X(\bt)_+$. Then $M$ admits a good  filtration.

\end{theorem}

\begin{theorem}\label{th:good tensor}\cite[II Proposition 4.21, B.9]{Jan}
Let $\lambda$, $\mu\in X(\bt)_+$. Then $\nabla_\lambda\otimes\nabla_\mu$ admits a good  filtration.
\end{theorem}
%--------------------------------------------------------------

%-------------------------------------------------------------
\begin{theorem}\cite[Theorem 3.2.7]{WvdK}
    Let $M$ be a $\fb$-module such that for every $\lambda\in X(\bt)$, $H^1( \fb,M\otimes Q(\lambda))=0$. Then, $M$ has an excellent filtration.
\end{theorem}

\subsection{Forms of the $P(\lambda)$, $Q(\mu)$ over the integers}
For the results over $\Z$ we will use  $\Z$-forms $P(\lambda)_\Z$, $Q(\mu)_\Z$ of  $P(\lambda)$, $Q(\mu)$ respectively.
They are constructed and studied in \cite[Chapter 7]{WvdKTata}. Both  $P(\lambda)_\Z$ and $Q(\mu)_\Z$ are finitely generated and free over $\Z$. 
They share many properties with their counterparts over fields.
We often drop the $\Z$ from  the notation $P(\lambda)_\Z$, $Q(\mu)_\Z$.
One could
follow \cite[section 3]{Ef}
to discuss in what sense one has highest weight category structures on the exact category of $\Bm$-modules that are finitely generated and free over $\Z$.
 Instead we will invoke the Universal coefficient Theorem for cohomology to move back and forth between the case that $\sk$ is a field and the case $\sk=\Z$. We need that tool anyway.
 
%--------------------------------------------------------------
\section{Rappels: Weyl groups}\label{sec:Weylgroups}
%--------------------------------------------------------------

%--------------------------------------------------------------
\subsection{The Demazure product}\label{subsec:Demazureproduct}
%--------------------------------------------------------------

Let $(W,S)$ be the Coxeter presentation of the Weyl group of a simple simply connected group $ \bg$ of rank $r$. 
The Demazure product on $W$ is defined by replacing the relation $s_{\alpha}\cdot s_{\alpha}=e$ for a simple root $\alpha$ with $s_{\alpha}\star s_{\alpha}=s_{\alpha}$. Equivalently, for any $w\in W$ and any simple root $\alpha$, define $w\star s_{\alpha}$ to be the longer of $ws_{\alpha}$ or $w$. A sequence of simple reflections is reduced if and only if the usual product $s_{i}s_{2}\dots s_{k}$ is equal to the Demazure product $s_{i}\star \dots \star s_{{k}}$.
If $v,w\in W$, then $\overline{\fb v\fb}~\overline{\fb w\fb}=\overline{\fb (v\star w)\fb}$.% and ${\sH}_v\circ {\sH}_w={\sH}_{v\star w}$ (see Section \ref{subsec:Demres} and Theorem \ref{th:Demazure_kernel}, cf. also \cite[Section 4.3]{dCHL}).

\

\begin{lemma}\label{lem:Hwstar}
    If $v,w\in W$, then ${\sH}_v\circ {\sH}_w={\sH}_{v\star w}$.
\end{lemma}
\begin{proof}
    See \cite[1.3., Corollary and Theorem 3.1]{CPS} or \cite[Proposition 2.2.5]{WvdKTata}.
\end{proof}
%--------------------------------------------------------------
%\subsection{The Bruhat order and its refinement}\label{subsec:refinements_Bruhat}
%--------------------------------------------------------------

%--------------------------------------------------------------
\subsection{Total order $\prec$ on $W$, refining the Bruhat order}\label{subsec:totalorderonW}
%--------------------------------------------------------------

  Choose a total order $\prec$ on $W$  that refines the Bruhat order $<$.
Thus $w\leq v$ implies $w\preceq v$. And $w\succ v$ implies $w\nleq v$.
To get a nice fit with Theorem \ref{th:triang} one may also arrange that $w\preceq v$ implies $ww_0\succeq vw_0$. This is optional.
\begin{remark}
    The order $\prec$ gives an order on the Steinberg basis, often in conflict with $<_a$ and $>_e$.
\end{remark}
%--------------------------------------------------------------
\subsection{Combinatorics}\label{subsec:comb}
%--------------------------------------------------------------

 First we recall some facts from \cite{Humphreys} about $W$ and its action on $X( \bt)$. As usual $w_0$
  denotes the longest element of $W$.
\begin{lemma}\label{lew0reverse} Let $v,w\in W$. The following are equivalent
\begin{itemize}
\item $v<w$ \item $v^{-1}<w^{-1}$ \item $w_0w<w_0v$ \item $ww_0<vw_0$.
\end{itemize}
\end{lemma}
\begin{proof}For the first three, see \cite[page 119]{Humphreys}.
Use that $(ww_0)^{-1}=w_0w^{-1}$, $(vw_0)^{-1}=w_0v^{-1}$.
\end{proof}

\begin{lemma}\label{lem:ws shorter} Let $w\in W$ and let $\alpha $ be a simple root.
Then $\ell(ws_\alpha)>\ell(w)$ if and only if $w(\alpha)>0$.
\end{lemma}
\begin{proof}See \cite[page 116]{Humphreys}.
\end{proof}

\begin{lemma}\label{lem:left fill} Let $w\in W$.
One can successively multiply $w$ on the right by simple reflections (increasing the length by 1) until this is no longer possible and $w_0$ is obtained.\end{lemma}

\begin{proof}See \cite[page 16]{Humphreys}. We usually prefer to multiply from the left.
\end{proof}

\begin{lemma}\label{lem:in w0}
Let $w,z\in W$ with $\ell(wz)=\ell(w)+\ell(z)$. Then $w\leq w_0z^{-1}$.
\end{lemma}
\begin{proof}
If $wz=w_0$, then it is clear. We argue by induction on $\ell(w_0)-\ell(wz)$.
If $\ell(w_0)>\ell(wz)$, then there is a simple $s$ with $\ell(swz)>\ell(wz)$.
We have $\ell(sw)=\ell(w)+1$ and $\ell(swz)=\ell(sw)+\ell(z)$. So by induction
hypothesis $sw\leq w_0z^{-1}$. And thus $w\leq sw\leq w_0z^{-1}$
\end{proof}


\begin{lemma}\label{lem:minimal-coset-representative}
If $\lambda$ is dominant, then its stabilizer $W_I$ in $W$ is generated by the simple reflections that fix $\lambda$. 

One puts $W^I=\{w\in W\mid \ell(ws)>\ell(w) \mbox{ for every simple reflection $s$ that fixes }\lambda\}$.

Every $w\in W$ may be written uniquely as $vz$, where $z\in W_I$ and $v\in W^I$ is the minimal coset representative of the coset $wW_I$. The assignment of $v$ to $w$ respects the Bruhat order.

\end{lemma}
\begin{proof}Compare  \cite[1.10, 1.15]{Humphreys}. For comparison with \cite[Section 3]{Ana}, note that, for $v\in W^I$, $\ell(v)$ is the number of reflecting hyperplanes that separate $v\lambda$ from $\lambda$.
\end{proof}


%\begin{lemma}\label{refine}
%If $u\leq w$ and $\ell(w)>\ell(u)+1$, then there is $v$ with $u<v<w$. \end{lemma}
%\begin{proof}
%\cite[p.122]{Humphreys}.
%\end{proof}

Recall that we denote the Demazure product on $W $ by $\star$.
 \begin{lemma}
 Let $v,w\in W$. There are $v'\leq v$, $w'\leq w$ so that $v'w=v\star w=vw'$,
$ \ell(v')+\ell(w)=\ell(v\star w)=\ell(v)+\ell(w')$.\qed
 \end{lemma}

\

\begin{lemma}\label{lem:caps}
Let $x$, $y$, $s\in W$ with $s$ simple and $x\leq y\star s$. Choose $v$ minimal so that $v\star s= x\star s$.
Then $v\leq y$.
\end{lemma}
\begin{proof}
Note that $v$ is the minimal representative of the coset $x\langle s\rangle$. If $w$ is the 
minimal representative of the coset $y\langle s\rangle$, then $v\leq w\leq y$.
\end{proof}

\begin{lemma}\label{lem:ea}
Let $\lambda$, $\mu$ be weights in the same $W$-orbit. Then $\lambda<_e\mu$ if and only if $\lambda>_a\mu$. 
\end{lemma}
\begin{proof}
If $\lambda\leq_e\mu$, then there is a dominant weight $\omega$ and there are $w, z\in W$ with $w\leq z$, $\lambda=w\omega$, $\mu=z\omega$. Then $ww_0\geq zw_0$ and $-\lambda=wz_0(-w_0\omega)$, $-\mu=zz_0(-w_0\omega)$, so $-\lambda\geq_e-\mu$, so $\lambda>_a\mu$. The converse is proved similarly.
\end{proof}

%--------------------------------------------------------------
\section{$ \bt$-equivariant $\sK$-theory of $ \bg / \fb$}\label{sec:rappelT-equivGmodB}
%--------------------------------------------------------------
In this section we introduce the theorem on triangular transition matrices.
\subsection{Two fixed points}
Let $\sk$ still be a field.
We will need results from Graham--Kumar \cite{GrahamKumar} on the $ \bt$-equivariant $\sK$-theory $\sK_{\bt}( \bg/\fb)$ of $ \bg/ \fb$.
Their $\fb$ is our $ \fb^+$. As both $\fb$ and $ \fb^+$ will be needed, let us consider the $\bg$-variety  
$\mathcal B$ of Borel subgroups.
If $x$ is a $ \bt$ fixed point in $\mathcal B$, then its stabilizer $ \fb(x)$ is a Borel subgroup and one identifies
$\mathcal B$ with $ \bg / \fb(x)$.
Let $x_+$ be the $ \bt$ fixed point the stabilizer of which is $ \fb(x_+)= \fb^+$. Similarly, let $x_-$ be the fixed point the stabilizer of which is $ \fb(x_-)= \fb$.
Choose a representative $\dot{w_0}$ of $w_0$.
One has $x_+= \dot{w_0}x_-$. We simply write $x_+= {w_0}x_-$.
Let $\phi_0$ be the isomorphism $ \fb^+\to\fb$ sending $b$ to 
$\dot{w_0}b\dot{w_0}^{-1}$. 
If $\F$ is a $ \bg $-equivariant vector bundle on $\mathcal B$, and $x$ is   fixed by $ \bt$, then the fibre $\F_x$ is a $ \fb(x)$-module.
We have $\F_{x_+}\cong \phi_0^*(\F_{x_-})$.
Recall that when $M$ is a finite dimensional $\fb$-module, we denote as in \cite{Jan} by $\Ll(M)$ the $ \bg $-equivariant vector bundle 
$\F$ with $\F_{x_-}=M$. If $M$ is a finite dimensional $ \fb^+$-module, we denote by $\Ll^+(M)$ the $ \bg $-equivariant vector bundle 
$\F$ with $\F_{x_+}=M$. 

If $\F$ is a $ \bt$-equivariant coherent sheaf on $ \bg / \fb$ and $N$ is a finite dimensional $ \bt$-module, then
$$\F\otimes_\sk N$$ is a $ \bt$-equivariant coherent sheaf.

When working with $ \fb^+$ one should define excellent filtrations and relative Schubert filtrations in terms
of $P^+(\lambda):=\phi_0^* (P(w_0\lambda))$ and
$Q^+(\lambda):= \phi_0^* (Q(w_0\lambda))$. They have $ \fb^+$-socles of weight $\lambda$ and are the costandard
modules of \cite{WvdK}.

One has 
 $$H^i( \fb^+,P^+(\lambda)\otimes_\sk Q^+(\mu))=
\begin{cases}\sk &\text{if $i=0$ and $\lambda+\mu=0$}\\
0&\text{else.}
\end{cases}
$$

%--------------------------------------------------------------
\subsection{The Steinberg basis}\label{subsec:Steinberg}
%--------------------------------------------------------------

For $v\in W$ the Steinberg weight $e_v$ is given by $$e_v = v^{-1} \sum_{v^{-1}\alpha_i<0} \omega_i.$$
The Steinberg basis $\{\sk_{e_v}\}_{v\in W}$ consists of the corresponding one dimensional
$\fb$-modules.

It provides by \cite{St} a basis of the representation ring $R( \bt)$ as a module over the representation ring $R( \bg )=\sK_0(\rep( \bg ))$,
where $\rep( \bg )$ is the category of finitely generated $ \bg $-modules.
We will find later (Corollary \ref{cor:Steinberg_alternative}) that one still gets a basis if one replaces a few $\sk_{e_v}$ by $Q(e_v)$ or $P(-e_v)^*$.

Generators of $R( \bt)$ are often written $e^\lambda$ instead of $[\sk_\lambda]$.
We will find later (Remark \ref{rem:Steinberg from anti}) that the Steinberg weights can also be described as follows: a weight $\lambda$ is a Steinberg weight if and only if $e^\lambda$ does not lie in the $R(\bg)$-submodule of $R(\fb)$
generated by the $e^\mu$ with $\mu<_a\lambda$.
 If $\lambda$ is a Steinberg weight, then $\lambda=e_w$ where $w\in W$ is the element of minimal length making $w\lambda$ dominant.

%--------------------------------------------------------------
\subsection{Schubert varieties and opposite Schubert varieties}
%--------------------------------------------------------------
\label{subsec:alphabet}
Let the Schubert variety $\fX^+_w$ be the closure of $ \fb^+wx_+$ and let the opposite Schubert variety
$\fX^w$ be the closure of 
$ \fb^-wx_+$. Its `boundary' $\partial \fX^w$ is the union of the $\fX^v$ that are strictly contained in $\fX^w$.

The $ \bt$-equivariant $\sK$-theory $\sK_{\bt}( \bg/\fb)$ of $ \bg / \fb$ is a module for the representation ring $R( \bt)$.
%Generators of $R( \bt)$ are often written $e^\lambda$ instead of $[\sk_\lambda]$.
If $M$ is a finitely generated $\bt$-module, then its class $[M]$ in $R( \bt)$
is also written $\Char(M)$, the formal
character of $M$.

If $\F$ is a $ \bt$-equivariant coherent  sheaf and $N$ is a finite dimensional $ \bt$-module, then $\F\otimes_\sk N $
represents $[N]\cdot[\F]$.

If $\sk=\cC$, then we learn from \cite{GrahamKumar} that the $R( \bt)$-module $\sK_{\bt}( \bg/\fb)$ has a Schubert basis $\{[\Oo_{\fX^+_w}]\}_{w\in W}$.
And also an `opposite Bruhat cell' basis $\{[\Oo_{\fX^w}(-\partial \fX^w)]\}_{w\in W}$.  
They are orthogonal  under the $R( \bt)$-bilinear symmetric
pairing $\langle -, -\rangle$ on $\sK_{\bt}( \bg/\fb)$ given by 
$$\langle [\F],[\G] \rangle=\sum_i(-1)^i\Char(H^i( \bg / \fb,\F\otimes \G ))\in R( \bt).$$
This pairing makes sense over any field $\sk$.

If $\F\otimes\G$ is supported on a $ \bt$ stable closed subscheme $Y$, then 
$$\langle [\F],[\G] \rangle=\sum_i(-1)^i\Char(H^i(Y,\F\otimes\G))\in R( \bt).$$

We put $$\tP_v=\Ll(P(-e_v))$$
%Notice that $v e_v$ is dominant.

and $$\tQ_v=\Ll(Q(e_v)).$$
They are $ \bg $-equivariant vector bundles, but we also view them as $ \bt$-equivariant vector bundles.

Write $\alpha_{vw}=\langle [\Oo_{\fX^+_w}],[\tQ_v] \rangle$. So
$$[\tQ_v]=\sum \alpha_{vw}[\Oo_{\fX^w}(-\partial \fX^w)]$$
when $\sk=\cC$.

Write $\beta_{vw}=\langle [\Oo_{\fX^w}(-\partial \fX^w)],[\tP_v] \rangle$.
So
$$[\tP_v]=\sum \beta_{vw}[\Oo_{\fX^+_w}]$$
when $\sk=\cC$.

Our main result concerning these matrices is
that, with a suitable reordering of rows and columns, the matrices 
$(\alpha_{wv})$ and $(\beta_{vw})$ are upper triangular and invertible. 

\begin{theorem}[Triangular transition matrices]\label{th:triang}
\
\begin{enumerate}
    \item $\beta_{vw}=\langle [\Oo_{\fX^w}(-\partial \fX^w)],[\tP_v] \rangle$ vanishes unless $v\,w_0  \leq w$
    %\preceq w$ 
    in the Bruhat order.
    \item $\alpha_{vw}=\langle [\Oo_{\fX^+_w}],[\tQ_v] \rangle$ vanishes unless 
    %$w\preceq v\;w_0$.
    $w\leq v\,w_0$ in the Bruhat order.
    \item If $v\,w_0  = w$,  then $\alpha_{vw}=[\sk_{ve_v}]$, $\beta_{vw}=[\sk_{-ve_v}]$.
\end{enumerate}
\end{theorem}




\section{Triangularity of transition matrices}\label{sec:triangular}
%--------------------------------------------------------------

In this section
we will show that, with a suitable reordering of rows and columns, the matrices 
$(\alpha_{wv})$ and $(\beta_{vw})$ from subsection \ref{subsec:alphabet} are upper triangular and invertible. 

\

Recall that a weight $\lambda$ of a $\fb$-module $M$ is called extremal if all other weights $\mu$
of $M$ are shorter or equally long:
$(\mu,\mu)\leq(\lambda,\lambda)$

A  closed subset (\emph{i.e.} reduced closed subscheme) of $\bg/\fb$ is $\fb$-invariant if and only if it is a union of Schubert varieties.
So there are only finitely many $\fb$-invariant  closed subsets of $\bg/\fb$.
\begin{lemma}\label{lem:Schubert module}
    Let $S, S'$ be  unions of Schubert varieties  in $ \bg /{\fb}$ and let $\lambda$ be dominant. 
\begin{itemize}\item The extremal weights of $\Gamma(S,\Ll(\lambda))$ are the $w\lambda$ 
with $wx_-\in S$.
\item  $\Gamma(\bg/\fb,\Ll(\lambda))\to\Gamma(S,\Ll(\lambda))$ is surjective.
\item If $S'\subset S$, then $\Gamma(S,\Ll(\lambda))\to\Gamma(S',\Ll(\lambda))$ is surjective, and its kernel has a relative Schubert filtration.
\item $H^i(S,\Ll(\lambda))=0$ for $i>0$.
\end{itemize}
\end{lemma}
\begin{proof}
    See  \cite[Proposition 2.2.15, Lemma 2.3.10, Lemma 2.2.11,  Proposition A.2.6]{WvdKTata}, \cite[Theorem 1.2.8, Chapter 2]{BK}.
\end{proof}

\begin{proposition}\label{prop:1dimQ}Let $\mu$ be a dominant weight with  $\alpha^\vee (\mu)\leq1$ for all simple roots $\alpha$.
Then
 $Q(-\mu)=\sk_{-\mu}$.
 \end{proposition}
 \begin{proof}
 
 We may assume $\sk=\cC$, because of the base change properties \cite[Chapter 7]{WvdKTata} of the $Q(\mu)$.
 Suppose $Q(-\mu)$ is larger than its socle $\sk_{-\mu}$. Let $\nu$ be a weight of the socle of $Q(-\mu)/\sk_{-\mu}$
 and let $f$ be a nonzero weight vector of $Q(-\mu)$, of weight $\nu$, 
 mapping to a vector of the socle of $Q(-\mu)/\sk_{-\mu}$.
 As the socle of $Q(-\mu)$ is its weight space $\sk_{-\mu}$, the vector  $f$ can not be fixed by the unipotent radical $\bu$
 of $\fb$. But it is well known
 that  $\bu(\cC)$ is generated by the $x_{-\alpha}(t)$ with $\alpha$ simple. So there must be such an $x_{-\alpha}(t)$ with 
 $x_{-\alpha}(t)(f)-f$ a nonzero vector in $\sk_{-\mu}$. In particular, $-\mu=\nu + n(-\alpha)$ for some integer $n$. But $\nu$
 must be strictly shorter than the unique extremal weight $-\mu$ of $Q(-\mu)$. However, there is no strictly shorter 
 weight in $-\mu +\Z\alpha$.
 \end{proof}
 
  \begin{lemma}\label{lem:Q acyclic}Let $w\in W$.
If $M$ is a finite dimensional $ \fb^+$-module with a relative Schubert filtration,
then $H^i(\overline{ \fb^+wx_+},\Ll^+(M))=0$ for $i>0$.
\end{lemma}
\begin{proof}
Compare \cite[Theorem 1.9.(a)(ii)]{WvdK}.
\end{proof}

 
 Fix $v\in W$. Let $I$ consist of the simple reflections that fix the dominant weight $ve_v$. 
 Let $s$ be simple.
 From  the definition of $e_v$ and Lemma \ref{lem:ws shorter} it follows that $\ell(v^{-1}s)>\ell(v^{-1})$ if and only if $s\in I$. 
 In particular, $v^{-1}$ is a minimal coset representative of $v^{-1}W_I$.

\begin{notation}
    If $v\in W$, then the boundary $\partial\overline{{\fb}v\fb}$ of 
    $\overline{{\fb}v\fb}$ is defined to be the union of the $\overline{{\fb}w\fb}$
    that are strictly contained in $\overline{{\fb}v\fb}$.
\end{notation}




 \begin{proposition}\label{prop:extremalQ}
 Let $u,v\in W$.
 \begin{enumerate}
 \item $-ve_v$ is an extremal weight of $\Gamma(\overline{{\fb}w_0x_-},\Ll(-w_0ve_v))$, but not of 
 $\Gamma((\partial \overline{{\fb}v\fb})\;\overline{{\fb}v^{-1}w_0x_-},\Ll(-w_0ve_v))$.
 \item If $-uve_v\neq -ve_v$, then $-uve_v$ is an extremal weight of 
 $\Gamma((\partial \overline{{\fb}v{\fb}})\;\overline{{\fb}v^{-1}w_0x_-},\Ll(-w_0ve_v))$.
 \end{enumerate}
 \end{proposition}
 \begin{proof}Part (1).
 The extremal weights of $\Gamma(\overline{{\fb}w_0x_-},\Ll(-w_0ve_v))$ are the elements in the $W$ orbit of $-w_0ve_v$,
 hence of $-ve_v$.
 Now suppose $-ve_v$ is  an extremal weight of $\Gamma((\partial \overline{{\fb}v{\fb}})\;\overline{{\fb}v^{-1}w_0x_-},\Ll(-w_0ve_v))$.
Then $-ve_v$ can be written as $y(-w_0ve_v)$ with $y\leq z\star v^{-1}w_0$ for some $z<v$.
Replacing $z$ by a lesser element if necessary, we may assume $y\leq zv^{-1}w_0$ for some $z<v$.
Put $u=yw_0$. Then $-ve_v=-uve_v$ with $u\geq zv^{-1}$. But $-ve_v=-uve_v$ implies $u\in W_I$, so 
$u\geq zv^{-1}$ implies $zv^{-1}\in W_I$. Thus $z^{-1}\in v^{-1}W_I$.
This contradicts the minimality of $v^{-1}$ in its coset $v^{-1}W_I$.

Part (2). Now consider a weight of the form $-uve_v$ with $-uve_v\neq -ve_v$. We may replace $u$ by its minimal coset representative.
As $\ell(u)\geq 1$ there is a simple reflection $s$ with $s\notin I$ and $\ell(u)= \ell(us)+1$. In particular, $u\geq s$.
Now $v^{-1}s<v^{-1}$ by the construction of $e_v$. Put $z=sv$, $y=uw_0$. Then $z<v$, $-uve_v=y(-w_0ve_v)$, with
$y\leq sw_0=z\star v^{-1}w_0$.
 \end{proof}
 
 
 \begin{proposition}\label{prop:vw_0 notless w}
 Let $v,w\in W$ such that $vw_0\nleq w$.
 Every extremal weight of $\Gamma(\overline{{\fb}(ww_0\star v^{-1}w_0)x_-},\Ll(-w_0ve_v))$ is also an 
 extremal weight of 
 $\Gamma((\partial \overline{{\fb}ww_0{\fb}})\;\overline{{\fb}v^{-1}w_0x_-},\Ll(-w_0ve_v))$.
 \end{proposition}

 \begin{proof}
 Suppose 
$\ell((ww_0)\star (v^{-1}w_0))=\ell(ww_0)+\ell(v^{-1}w_0)$. By Lemma \ref{lem:in w0}
 $(ww_0)\leq v$, so $vw_0\leq w$, contrary to our assumption.
So $\ell((ww_0)\star (v^{-1}w_0))<\ell(ww_0)+\ell(v^{-1}w_0)$ and we have $y<ww_0$ with 
$(ww_0)\star (v^{-1}w_0))=yv^{-1}w_0 $, $\ell(yv^{-1}w_0)=\ell(y)+\ell(v^{-1}w_0)$. Every extremal weight of $\Gamma(\overline{{\fb}(ww_0\star v^{-1}w_0)x_-},\Ll(-w_0ve_v))$
is of the form $z(-w_0ve_v)$ with $z\leq yv^{-1}w_0$.
And $\Gamma((\partial \overline{{\fb}ww_0{\fb}})\;\overline{{\fb}v^{-1}w_0x_-},\Ll(-w_0ve_v))$ maps onto 
$\Gamma( \overline{{\fb}yv^{-1}w_0x_-},\Ll(-w_0ve_v))$ and then onto
$\Gamma( \overline{{\fb}zx_-},\Ll(-w_0ve_v))$ by Lemma \ref{lem:Schubert module}.
 \end{proof}
%Recall that $ G \times ^{\fb} G /{\fb}\cong  G /{\fb}\times  G /{\fb}$ by way of $(g,h)\mapsto (gh{\fb},g{\fb})$.
  


\

\begin{lemma}\label{lem:capw}
Let $S'$, $S''$ be unions of Schubert varieties with inverse images $\tilde S'$, $\tilde S''$ respectively in $\bg$. Let $w\in W$. 
Then $\tilde S'(\overline{{\fb}w{\fb}}/{\fb})\cap \tilde S''(\overline{{\fb}w{\fb}}/{\fb})=(\tilde S'\cap\tilde S'')(\overline{{\fb}w{\fb}}/{\fb})$.
\end{lemma}
\begin{proof}Clearly the right hand side is contained in the left hand side.
Both sides are unions of Schubert varieties.
First let $w$ be simple,
say $w=s$. Consider $x\in W$ with $\overline{{\fb}x{\fb}}/{\fb}$ contained in the left hand side. Choose $v$ minimal so that $v\star s= x\star s$.
Use Lemma \ref{lem:caps} to see that $\overline{{\fb}(v\star s){\fb}}/{\fb}$ is contained in the right hand side. But then so is $\overline{{\fb}x{\fb}}/{\fb}$. Next, if $\ell(w)>1$, choose a reduced expression $s_1\cdots s_n$ of $w$ and write $\overline{{\fb}w{\fb}}$ as $\overline{{\fb}s_1{\fb}}\cdots  \overline{{\fb}s_n{\fb}}$. Show that $\tilde S'(\overline{{\fb}w{\fb}})\cap \tilde S''(\overline{{\fb}w{\fb}})=(\tilde S'\cap\tilde S'')(\overline{{\fb}w{\fb}})$.
\end{proof}
 \begin{notation}
Let $S$ be a $B$-invariant closed subset union of  $ \bg /{\fb}$ and let $\F$ be a $\fb$-equivariant vector bundle on $S$. If $\sk$ is a field, put
$$\chi(S,\F)=\langle[\Oo_S],[\F]\rangle=\sum_i(-1)^i\Char(H^i(S,\F))$$
 in $R(T)$.
\end{notation}  

\

  \begin{proposition}\label{prop:Leray}
  Let $\lambda$ be a dominant weight and let $w\in W$. Let  $\tilde S$ be the inverse image of $S$ in $ \bg $.
   \begin{enumerate} 
 \item $H^0(\overline{{\fb}w{\fb}} S,\Ll(\lambda))\to H^0({\overline{{\fb}w{\fb}}/{\fb}},\Ll(H^0(S,\Ll(\lambda)))$ is an isomorphism.
 \label{Leray1}
  \item  $H^i({\overline{{\fb}w{\fb}}/{\fb}},\Ll(H^0(S,\Ll(\lambda)))=0$ for $i>0$.
  \label{Leray2}
 \item If $\sk$ is a field, then
 $
     \chi(\tilde S\overline{{\fb}w{\fb}}/\fb ,\Ll(\lambda))= \chi(S,\Ll(H^0({\overline{{\fb}w{\fb}}/{\fb}},\Ll(\lambda))).
 $
 \label{Leray3}
  \end{enumerate}
 \end{proposition}
\begin{proof}
If $\sk$ is a field of finite characteristic,
then parts \ref{Leray1} and \ref{Leray2} follow from \cite[Proposition 2.24, Theorem 1.9.(a)(ii)]{WvdK}  by induction on $\ell(w)$.
By the Universal coefficient Theorem \ref{th:universal coefficients}
both parts go through when $\sk=\Z$ or when $\sk$ is a field of characteristic zero.

Proof of part \ref{Leray3}. By the previous parts the result holds when $S$ is a Schubert variety. 
There are only finitely many possibilities for $S$, so we may assume the result for any union of Schubert varieties that is
strictly contained in $S$. If $S$ is not a Schubert variety, write $S=S'\cup S''$ where $S'$, $S''$ are strictly smaller.
Put $X=\overline{{\fb}w{\fb}}/\fb$
and $M=H^0(X,\Ll(\lambda))$.
If $\sk$ is a field of finite characteristic, then $S'\cap S''$ is reduced by Ramanathan \cite[Proposition 1.2.1, Chapter 2]{BK}. By \cite[Corollary 1.6.6]{BK}
it is then also reduced in characteristic zero.
From the Mayer--Vietoris sequence 
$$0\to  H^0(S,\Ll(M))\to H^0(S',\Ll(M))\oplus H^0(S'',\Ll(M))\to H^0(S'\cap S'',\Ll(M))\to\cdots$$
we get
\begin{equation}\label{eq:chiM}
    \chi(S,\Ll(M))=\chi(S',\Ll(M))+\chi(S'',\Ll(M))- \chi(S'\cap S'',\Ll(M)).
\end{equation}
Similarly the Mayer--Vietoris sequence 
$$0\to H^0(\tilde SX ,\Ll(\lambda))\to H^0(\tilde S'X ,\Ll(\lambda))\oplus H^0(\tilde S''X ,\Ll(\lambda))\to H^0(\tilde S'X\cap \tilde S''X ,\Ll(\lambda))\to0,$$ 
gives 
\begin{equation}\label{eq:chilambda}
    \chi(\tilde SX ,\Ll(\lambda))=\chi(\tilde S'X ,\Ll(\lambda)) + \chi (\tilde S''X ,\Ll(\lambda))- \chi(\tilde S'X\cap \tilde S''X ,\Ll(\lambda))
.\end{equation}
By Lemma \ref{lem:capw} we know that $\tilde S'X\cap \tilde S''X =(\tilde S'\cap \tilde S'')X $.
As the right hand sides of equations (\ref{eq:chiM}), (\ref{eq:chilambda})
agree, part \ref{Leray3} follows.

Instead of invoking \cite[Corollary 1.6.6]{BK}, we could choose $N>1$ so that all relevant $H^i(Y,\F)$ are flat over $\Z[1/N]$ and use the Universal coefficient Theorem \ref{th:universal coefficients}
to pass from the case of fields of finite characteristic to the case $\sk=\Z[1/N]$ and from there to the case where $\sk$ is a field of characteristic zero.
\end{proof}

\begin{remark}
    We will not need that 
    $H^i(S,\Ll(H^0({\overline{{\fb}w{\fb}}/{\fb}},\Ll(\lambda)))$ vanishes for $i>0$.
\end{remark}

\

\begin{lemma}\label{lem:kerdecomp}
 Let $S$, $S'$ be  unions of Schubert varieties in $ \bg /{\fb}$ with $S'\subset S$  and let $\lambda$ be dominant.  
  Put $M=\ker({\sH}_w(H^0(S,\Ll(\lambda)))\to {\sH}_w(H^0(S',\Ll(\lambda))))$. Then $M$ has a relative Schubert filtration. The sections of the filtration are the $Q(\mu)$ with $\mu$ an extremal weight of 
${\sH}_w(H^0(S,\Ll(\lambda)))$ but not of ${\sH}_w(H^0(S',\Ll(\lambda)))$. The multiplicity of such $Q(\mu)$ is one, and $\mu\in W\lambda$.
\end{lemma}

\begin{proof}
By Proposition \ref{prop:Leray} we may use Lemma \ref{lem:Schubert module}
to conclude that $M$ has a relative Schubert filtration.
To find out what the sections in the filtration are, it suffices to look at the characters.
By Proposition \ref{prop:multone} both ${\sH}_w(H^0(S,\Ll(\lambda)))$ and ${\sH}_w(H^0(S',\Ll(\lambda)))$ have a character that is a $\Z$-linear combination of characters $\Char(Q(\nu))$ with $\nu\in W\lambda$.
So $M$ also has such a character. But $\Char(Q(\nu))-e^\nu$ only involves weights $\tau$ with $(\tau,\tau)<(\nu,\nu)$. So by inspecting the
extremal weights one gets the coefficients in the expansion of $\Char(M)$
in terms of the $\Char(Q(\nu))$.
\end{proof}

%--------------------------------------------------------------
\subsection{The matrix $\beta_{vw}$}\label{subsec:beta}
%--------------------------------------------------------------
Recall that $\beta_{vw}=\langle [\Oo_{\fX^w}(-\partial \fX^w)],[\tP_v] \rangle$.
We want to pair $[\Oo_{\fX^w}(-\partial \fX^w)]$
with $[\tP_v]=[\Ll(P(-e_v))]$. 
As $0\to\Oo_{\fX^w(-\partial \fX^w)}\to \Oo_{\fX^w}\to \Oo_{\partial \fX^w} \to0$
is exact, it suffices to compute the difference between $\langle [ \Oo_{\fX^w}],[\Ll(P(-e_v))]\rangle$ and $\langle [ \Oo_{\partial \fX^w}],[\Ll(P(-e_v))]\rangle$.


Note that $\fX^w=\overline{{\fb}ww_0x_-}$.

And $P(-e_v)=\Gamma(\overline{{\fb}v^{-1}w_0x_-},\Ll(-w_0ve_v))$.


We hope that $\langle[\fX^w(-\partial \fX^w)],[\tP_v]\rangle$ vanishes unless $vw_0\leq w$.
But the propositions above show that $\langle [ \Oo_{\fX^w}],[\Ll(P(-e_v))]\rangle$  is just the character of the module
$\Gamma(\overline{{\fb(}ww_0\star v^{-1}w_0)x_-},\Ll(-w_0ve_v))$.
Similarly $\langle [ \Oo_{\partial \fX^w}],[\Ll(P(-e_v))]\rangle$ is just the character of the module
$\Gamma(\partial^{w,v},\Ll(-w_0ve_v))$,
where $\partial^{w,v}$ equals 
$(\partial \overline{{\fb}ww_0{\fb}})\;\overline{{\fb}v^{-1}w_0{\fb}}x_-$. Both modules have a relative Schubert filtration and we can get a grip on them by inspecting the extremal weights. (The weights that have the same length as the dominant weight $ve_v$.)



There are three cases.

\begin{itemize}
\item $vw_0=w$. So that is about  the kernel of the surjective map from $\Gamma(\overline{{\fb}w_0x_-},\Ll(-w_0ve_v))$ to $\Gamma(\partial^{w,v},\Ll(-w_0ve_v))$.
One checks it to be $Q(-ve_v)=\sk_{-ve_v}$.
\item $vw_0\nleq w$. Then one sees that 
the kernel of the surjective map from $\Gamma(\overline{{\fb}(ww_0\star v^{-1}w_0)x_-},\Ll(-w_0ve_v))$ to $\Gamma(\partial^{w,v},\Ll(-w_0ve_v))$ vanishes.
\item $vw_0< w$. No claims here.
\end{itemize}

We conclude that the matrix 
 $\beta_{vw}=\langle  [\Oo_{\fX^w}(-\partial \fX^w)],[\tP_v]\rangle$ behaves  as claimed in Theorem \ref{th:triang}.

%--------------------------------------------------------------
\subsection{The matrix $\alpha_{vw}$}\label{subsec:alpha}
%--------------------------------------------------------------

So let us turn to $\alpha_{vw}=\langle [\Oo_{\fX_w}],[\tQ_v]\rangle$.

Now we are dealing with the situation where the Borel subgroup has positive roots, and anti-dominant weights
are to be used as in \cite{WvdKTata}. For instance,  Proposition \ref{prop:Leray} gives
 \begin{proposition}\label{prop:Leray+}
  Let $\lambda$ be an anti-dominant weight. And let $S$ be a ${\fb}^+$ invariant closed subset of $ \bg /{\fb}^+$.
   \begin{enumerate} 
 \item $H^0(\overline{{\fb}^+w{\fb}^+} S,\Ll^+(\lambda))\to H^0({\overline{{\fb}^+w{\fb}^+}/{\fb}^+},\Ll^+(H^0(S,\Ll^+(\lambda)))$ is an isomorphism.
 \label{Leray+1}
  \item  $H^i({\overline{{\fb}^+wx_+}},\Ll^+(H^0(S,\Ll^+(\lambda)))=0$ for $i>0$.
  \label{Leray+3}
 \qed\end{enumerate}
 
 \end{proposition}


By Lemma \ref{lem:Q acyclic} we have 

$\alpha_{vw}=\langle [\Oo_{\fX_w}],[\tQ_v]\rangle=\Gamma(\overline{{\fb}^+wx_+},\Ll^+(\phi_0^* (Q(e_v))))$. Now 
$\phi_0^* (Q(e_v)))$ has $\fb^+$-socle of weight $w_0e_v$ and it is therefore the kernel of the surjection
$\Gamma(\overline{{\fb}^+w_0v^{-1}w_0x_+},\Ll^+(w_0ve_v))\onto\Gamma(\partial(\overline{{\fb}^+w_0v^{-1}w_0x_+}),\Ll^+(w_0ve_v)).$
Combining with Proposition \ref{prop:Leray+}
we see  $\alpha_{vw}$ is the character of the kernel of the surjective map from
$\Gamma(\overline{({\fb}^+w{\fb}^+})(\overline{{\fb}^+w_0v^{-1}w_0x_+}),\Ll^+(w_0ve_v))))$ to $\Gamma(\overline{({\fb}^+w{\fb}^+})\partial(\overline{{\fb}^+w_0v^{-1}w_0x_+}),\Ll^+(w_0ve_v))))$.

\

\begin{lemma}
\label{lem:s to boundary}Let $J=w_0Iw_0$ consist of the simple reflections that fix $w_0ve_v$. 
\begin{enumerate}
\item $w_0v^{-1}w_0$  is a minimal coset representative in $W/W_J$.
\item $\partial(\overline{{\fb}^+w_0v^{-1}w_0x_+})$ contains the union of the 
$\overline{{\fb}^+w_0v^{-1}w_0sx_+}$ with $s$ simple, $s\notin J$.
\end{enumerate}
\end{lemma}
\begin{proof}Let  $\alpha$ be a simple root.
As $e_v$ is a Steinberg weight, $w_0v^{-1}(\alpha)>0$ if and only if $s_\alpha\notin I$.
So $w_0v^{-1}w_0(-w_0(\alpha))<0$ if and only if $s_{-w_0\alpha}\notin J$.
Thus if $\beta$ is a simple root, then $w_0v^{-1}w_0(\beta)<0$ if and only if $s_{\beta}\notin J$.
\end{proof}
\begin{proposition}\label{prop:extremal2}
 Let $u,v\in W$.
 \begin{enumerate}
 \item $ve_v$ is an extremal weight of $\Gamma(\overline{({\fb}^+vw_0{\fb}^+})(\overline{{\fb}^+w_0v^{-1}w_0x_+}),\Ll^+(w_0ve_v))))$, but not of 
$\Gamma((\overline{{\fb}^+vw_0\fb^+})
\partial(\overline{{\fb}^+w_0v^{-1}w_0x_+}),\Ll^+(w_0ve_v))))$.
 \item If $uve_v\neq ve_v$, then $uve_v$ is an extremal weight of $\Gamma(\overline{({\fb}^+vw_0\fb^+})
\partial(\overline{{\fb}^+w_0v^{-1}w_0x_+}),\Ll^+(w_0ve_v))))$.
 \end{enumerate}
 \end{proposition}
\begin{proof}Part (1).
 The extremal weights of 
 $$M:=\Gamma(\overline{({\fb}^+vw_0{\fb}^+})(\overline{{\fb}^+w_0v^{-1}w_0x_+}),\Ll^+(w_0ve_v))))$$ are  elements in the $W$ orbit of $w_0e_v$.
 We have
 $\ell(vw_0)+\ell(w_0v^{-1}w_0)=\ell(w_0)$, so  $(vw_0)\star (w_0v^{-1}w_0)=w_0$, and $w_0w_0ve_v$ is one of the weights of $M$.
 Now suppose $ve_v$ is  an extremal weight of 
$\Gamma((\overline{{\fb}^+vw_0\fb^+})
\partial(\overline{{\fb}^+w_0v^{-1}w_0x_+}),\Ll^+(w_0ve_v))))$.
Then $ve_v$ can be written as $y(w_0ve_v)$ with $y\leq vw_0\star z$ for some $z<w_0v^{-1}w_0$.
Replacing $z$ by a lesser element if necessary, we may assume $y\leq vw_0z$ for some $z<w_0v^{-1}w_0$.
Put $u=w_0y$. Then $ve_v=yw_0ve_v=w_0uw_0ve_v$ with $u\geq w_0vw_0z$. But $uw_0ve_v=w_0ve_v$ implies $u\in W_J$, so 
$u\geq w_0vw_0z$ implies $w_0vw_0z\in W_J$. This contradicts the minimality of $w_0v^{-1}w_0$ in its coset $w_0v^{-1}w_0W_J$.

Part (2). Now consider a weight of the form $uve_v$ with $uve_v\neq ve_v$. 
Then $w_0uw_0w_0ve_v\neq w_0ve_v$ so $w_0uw_0\notin W_J$. There is a simple $s$, $s\notin J$, with $w_0uw_0\geq s$.
Then $uw_0\leq w_0s=(vw_0)\star (w_0v^{-1}w_0s)$ and $\overline{{\fb}^+w_0v^{-1}w_0sx_+}$ is contained in 
$\partial(\overline{\fb^+w_0v^{-1}w_0x_+})$ by Lemma \ref{lem:s to boundary}. So $uw_0w_0ve_v$ is an extremal weight of
$\Gamma((\overline{{\fb}^+vw_0\fb^+})
\partial(\overline{{\fb}^+w_0v^{-1}w_0x_+}),\Ll^+(w_0ve_v))))$.
 \end{proof}

 \begin{proposition}\label{prop:w notless vw_0}
 Let $v,w\in W$ such that $w\nleq vw_0$.
 Every extremal weight of $\Gamma(\overline{({\fb}^+w{\fb}^+})(\overline{{\fb}^+w_0v^{-1}w_0x_+}),\Ll^+(w_0ve_v))))$  is also extremal weight of 
$\Gamma(\overline{({\fb}^+w{\fb}^+})\partial(\overline{{\fb}^+w_0v^{-1}w_0x_+}),\Ll^+(w_0ve_v))))$.

 \end{proposition}
 
 \begin{proof}
 The extremal weights of  $\Gamma(\overline{({\fb}^+w{\fb}^+})(\overline{{\fb}^+w_0v^{-1}w_0x_+}),\Ll^+(w_0ve_v))))$ are of the form
 $uw_0ve_v$ with $u\leq w\star (w_0v^{-1}w_0)$. Suppose $\ell(w)+\ell(w_0v^{-1}w_0)=\ell(ww_0v^{-1}w_0)$.
 Then $w\leq w_0w_0vw_0$, contrary to the assumption. So $\ell(w)+\ell(w_0v^{-1}w_0)<\ell(ww_0v^{-1}w_0)$
 and we may chose $z<w_0v^{-1}w_0$ with $w\star (w_0v^{-1}w_0)=w\star z$. As $u\leq w\star z$, we see that $uw_0ve_v$ is an extremal 
 weight of $\Gamma(\overline{({\fb}^+w{\fb}^+})\partial(\overline{{\fb}^+w_0v^{-1}w_0x_+}),\Ll^+(w_0ve_v))))$
 \end{proof}


One can now deal with the matrix  $\{\alpha_{vw}\}$
in the same manner as for  $\{\beta_{vw}\}$. 

This ends the proof of Theorem \ref{th:triang}.
We next develop its consequences.

%---------------------------------------------------------------------
\section{$\ind$-vanishing for $\fb$-modules}\label{sec:b-cohom_vanish}
%---------------------------------------------------------------------
In this section we draw conclusions about $\fb$-modules by combining  the vanishing in Corollary \ref{cor:PQnabla} with Theorem \ref{th:triang} on triangular transition matrices. 
%--------------------------------------------------------------
\subsection{Cohomological descent from $ \bg/ \fb$ to $\fb$}\label{subsec:cohdesc}
%--------------------------------------------------------------
Let $\sk$ be a field or $\Z$.

Let $\Ff=\Ll ({\sf F})={ \bg}\times ^{\fb}\sf F$ be a $ \bg$-equivariant sheaf on $ \bg/ \fb$ given by a finitely generated $\fb$-module $\sf F$, which is projective over $\sk$. By  \cite[I 3.3, Proposition 5.12]{Jan}  there are isomorphisms 
\begin{equation}\label{eq:cohdescent}
{\rm H}^p({\fb},\sk[{ \bg}]\otimes {\sf F})=\RR^p\ind _{\fb}^{ \bg}(\sf F)={\rm H}^{p}({ \bg/\fb},\Ff),
\end{equation} 
 so it is all about computing sheaf cohomology of a $ \bg$-equivariant coherent sheaf via derived induction.
By \cite[II, Proposition 4.20, B.8]{Jan}, we know that  $\sk[ \bg]$ as a ${ \bg}\times { \bg}$-module via $\rho _l\times \rho _r$ admits a good filtration whose factors are $\nabla _{\lambda}\otimes \nabla _{-w_0\lambda}$ 
with $\lambda \in X({\bt})_+$, each occurring with multiplicity one; that is, $\sk[ \bg]=\varinjlim A_i$ such that $A_i/A_{i-1}=\nabla _{\lambda}\otimes \nabla _{-w_0\lambda}$. The action of ${\fb}$ on $\sk[ \bg]$ is by way of $\rho _r$. 

\begin{lemma}\label{lem:cohdescentlemma}
Assume $H^p(\fb,{\sf F}\otimes\nabla _{\mu})=0$ for $p>0$ and all $\mu$. 
\begin{enumerate}\item Then ${\rm H}^{p}({ \bg/\fb},\Ff)=0$ for all $p> 0$ and ${\rm H}^{0}({ \bg/\fb},\Ff)$ has a good filtration with with associated graded 
$$\bigoplus _{\lambda \in X({\bt})_+}\nabla _{\lambda}\otimes{\rm H}^0({\fb}, \nabla _{-w_0\lambda}\otimes {\sf F}).$$
\item If moreover $H^0(\fb,{\sf F}\otimes\nabla _{\mu})=0$ for all $\mu$, then ${\rm H}^{0}({ \bg/\fb},\Ff)=0$.
\end{enumerate}
\end{lemma}

\begin{proof}
\

(1) By (\ref{eq:cohdescent}), ${\rm H}^{p}({ \bg/\fb},\Ff)={\rm H}^p({\fb},\sk[{ \bg}]\otimes {\sf F})$, 
and using the good filtration on $\sk[ \bg]$ we obtain 
${\rm H}^p({\fb},\sk[{ \bg}]\otimes {\sf F})={\rm H}^p({\fb},\varinjlim A_i\otimes {\sf F})={\rm H}^p({\fb},\varinjlim (A_i\otimes {\sf F}))$, 
because direct limits commute with tensoring with finitely generated projective $\sk$-modules. 

By \cite[I. Lemma 4.17]{Jan}, ${\rm H}^p({\fb},\varinjlim (A_i\otimes {\sf F}))=\varinjlim{\rm H}^p({\fb}, A_i\otimes {\sf F})$. 
Now ${\rm H}^p({\fb}, (A_i/A_{i-1})\otimes {\sf F})={\rm H}^p({\fb},\nabla _{\lambda}\otimes \nabla _{-w_0\lambda}\otimes {\sf F})=\nabla _{\lambda}\otimes {\rm H}^p({\fb},\nabla _{-w_0\lambda}\otimes {\sf F})=0$ for $p>0$, as ${\fb}$ acts on $\sk[ \bg]$ 
by way of $\rho _r$, and ${\rm H}^p({\fb},\nabla _{-w_0\lambda}\otimes {\sf F})=0$ for $p>0$ by the assumption; the long cohomology sequence then gives both statements.
 
 
 
(2) This follows from the preceding item.
\end{proof}

%First observe
\begin{theorem}\label{th:highInd}
Let $\lambda,\mu\in X( \bt)$. Then
\begin{enumerate}
\item $\RR^i\ind_{\fb}^ \bg (P(\lambda)\otimes Q(\mu))=0$ for $i>0$.
\item $\ind_{\fb}^ \bg (P(\lambda)\otimes Q(\mu))$ has a good filtration.
\end{enumerate}
\end{theorem}
\begin{proof}
By Corollary \ref{cor:PQnabla} this follows from (\ref{eq:cohdescent}) and Lemma \ref{lem:cohdescentlemma} above.
\end{proof}

\begin{corollary}\label{cor:P-Q_pairing}
 $\langle[\Ll(P(\lambda))],[\Ll(Q(\mu))]\rangle=\ind_{\fb}^ \bg (P(\lambda)\otimes Q(\mu))$.\qed
\end{corollary}

\

The proof of Theorem \ref{th:highInd} also works over $\Z$. Or one may use the Universal 
coefficient Theorem \ref{th:universal coefficients}. Thus

\begin{theorem}\label{th:highIndZ}
Let $\lambda,\mu\in X( \bt)$.
Then
\begin{enumerate}
\item $\RR^i\ind_{\Bm}^{\Gm} (P(\lambda)_\Z\otimes Q(\mu)_\Z)=0$ for $i>0$.
\item $\ind_{\Bm}^{\Gm} (P(\lambda)_\Z\otimes Q(\mu)_\Z)$  has a good filtration.\qed
\end{enumerate}
\end{theorem}


\begin{theorem}\label{th:indSteinberg_vanishing}
If $\sk=\cC$, then
\begin{enumerate}
\item
$\ind_{\fb}^ \bg (P(-e_v)\otimes Q(e_w))$ vanishes unless $w\leq v$.
\item
$\ind_{\fb}^ \bg (P(-e_v)\otimes Q(e_v))=\sk$.
\end{enumerate}
\end{theorem}
\begin{proof}
$\langle[\tP_v],[\tQ_w]\rangle=\sum_{y, z\in W}\alpha_{wz}\beta_{vy}\langle[\Oo_{X^+_y}],[\Oo_{X^z}(-\partial X^z)]\rangle
=\sum_{y\in W}\alpha_{wy}\beta_{vy}$.
Now apply Theorem \ref{th:triang}.
\end{proof}
We wish to get rid of the restriction $\sk=\cC$.
As Euler characteristics are robust, this provides no difficulty:

The $\ind_{\Bm}^{\Gm} (P(\lambda)_\Z\otimes Q(\mu)_\Z)$ are 
free $\Z$-modules of finite rank,
so the Universal coefficient Theorem \ref{th:universal coefficients}
shows that the  theorem generalizes from $\sk=\cC$
to $\Z$, which in turn implies the case where $\sk$ is a field, again by the Universal coefficient Theorem.


Thus

\begin{theorem}\label{th:indPQ}Let $\sk$ be a field or $\Z$.
Let $v,w\in W$.
\begin{enumerate}
\item $\RR^i\ind_\fb^\bg (P(-e_v)\otimes Q(e_v))=\begin{cases}\sk &\text{if $i=0$}\\
0&\text{else.}\end{cases}$
\item If $w\not\leq v$ then $\RR^i\ind_\fb^\bg (P(-e_v)\otimes Q(e_w))=0$ for all $i$.
\end{enumerate}
\end{theorem}

\begin{corollary}\label{cor:left orthogonal}
Let $M$, $N$ be finite dimensional $ \bg$-modules.
\begin{enumerate}
\item $\Ext^i_{\fb}(M\otimes P(-e_v)^*,N\otimes Q(e_v))=\Ext^i_ \bg (M,N) $ for all $i$.
\item If $w\succ v$ then $\Ext^i_{\fb}(M\otimes P(-e_v)^*,N\otimes Q(e_w))=0$ for all $i$.
\end{enumerate}
\end{corollary}
\begin{proof}
$(1).$ We have a spectral sequence  \cite[I Proposition 4.5]{Jan}
$$\Ext^i_ \bg (M\otimes N^*,\RR^j\ind_{\fb}^ \bg (P(-e_v)\otimes Q(e_w)))\Rightarrow \Ext^{i+j}_{\fb}(M\otimes N^*,P(-e_v)\otimes Q(e_w)).$$
By Theorem \ref{th:highInd} (1), $\RR^j\ind_{\fb}^ \bg (P(-e_v)\otimes Q(e_w))=0$ for $j>0$. Thus, the above spectral sequence degenerates, and with Theorem \ref{th:indPQ}  this gives
$\Ext^{i}_{\fb}(M\otimes N^*,P(-e_v)\otimes Q(e_v))=\Ext^i_ \bg (M\otimes N^*,\RR^0\ind_{\fb}^ \bg (P(-e_v)\otimes Q(e_v)))=
\Ext^i_ \bg (M\otimes N^*,\sk)=\Ext^i_ \bg (M,N)$.



$(2).$ If $w\succ v$ then $w\not\leq v$, and the above spectral sequence still degenerates by Theorem \ref{th:highInd}, and 
Theorem \ref{th:indPQ} gives $\Ext^i_{\fb}(M\otimes P(-e_v)^*,N\otimes Q(e_w))=0$ for all $i$.
\end{proof}

\begin{remark}\label{rem:fr left orthogonal}
Let $\sk=\Z$. Then one should take $M$, $N$ finitely generated and free over $\Z$ in this proof. But see  Corollary \ref{cor:Z left orthogonal} below.
\end{remark}

%--------------------------------------------------------------
\section{Rappels: triangulated categories}\label{sec:rappel_triangulated}
%--------------------------------------------------------------

%--------------------------------------------------------------
\subsection{Triangulated hull }
%--------------------------------------------------------------

\begin{definition}[Triangulated hull]
Given a set $S$ of objects of
a triangulated category $\mathcal D$, its triangulated hull $\hull(S)$ is the smallest strictly full triangulated subcategory that contains $S$.
The triangulated hull of a sequence of subcategories is defined similarly.
\end{definition}

%--------------------------------------------------------------
\subsection{Thick hull }
%--------------------------------------------------------------
\begin{definition}[Thick]A subcategory  of
a triangulated category is called thick if it is closed under direct summands of objects.
\end{definition}

\begin{definition}[Thick hull]
Given a  set $S$ of objects of
a triangulated category $\mathcal D$, its thick hull $\hull(S)_{\oplus}$ is the smallest strictly full triangulated subcategory that contains $S$ which is closed under direct summands of objects.
\end{definition}


%--------------------------------------------------------------
\subsection{Admissible subcategories}\label{subsec:amdsubcat}
%--------------------------------------------------------------

We follow the exposition of  \cite[Section 2.1]{Ef}.
Let $\N\subset \sd$ be a full triangulated subcategory. The right orthogonal to $\N$ is the full subcategory $\N ^{\perp}\subset \sd $ consisting of all objects $X$ such that $\Hom _{\sd}(Y,X)=0$ for any $Y\in \N$. The left orthogonal $^{\perp}\N$ is defined analogously. The orthogonals are also triangulated subcategories.


\begin{definition}\label{def:admissible}
{A full triangulated subcategory $\A$ of $\D$ is called
{\it right admissible} if the inclusion functor $\A\hookrightarrow \D$ has a right adjoint. Similarly, $\A$ is called {\it left
  admissible} if the inclusion functor has a left adjoint. Finally,
$\A$ is {\it admissible} if it is both right and
left admissible.}
\end{definition}


\begin{lemma}\cite[Lemma 3.1]{Bondal}\label{lem:notorious}
Let $\N\subset \D$ be a full triangulated subcategory, and consider its right orthogonal $\N ^{\perp}\subset \D$. Then the following are equivalent:

\begin{itemize}

\item[(1)] The category $\D$ is generated by $\N$ and $\N ^{\perp}$, i.e.
$\hull(\N ^{\perp},\N)=\D$.

\item[(2)] For each object $x\in \D$ there exists an exact triangle $b\rightarrow x\rightarrow c$ with $b\in \N,c\in \N ^{\perp}$.

\item[(3)] The embedding functor $i_{\ast}:\N\hookrightarrow \D$ has a right adjoint $i^!:\D\rightarrow \N$.

\item[(4)] The embedding functor $j_{\ast}:\N ^{\perp}\hookrightarrow \D$ has a left adjoint $j^{\ast}:\D\rightarrow \N ^{\perp}$.

\end{itemize}

\end{lemma}

If the subcategory $\N$ is right (resp., left) admissible, then $\N$ is thick and one can consider the Verdier localization functor with respect to $\N$. We also record the following version of Lemma \ref{lem:notorious}:


\begin{proposition}\cite[Proposition 3.2.8]{Kr}
Let $\N\subset \D$ be a thick subcategory. Then the following are equivalent:

\begin{itemize}

\item[(1)] The inclusion functor $\N\subset \D$ has a right adjoint (i.e., $\N$ is right admissible in $\D$).

\item[(2)] For each object $b\in \D$ there exists an exact triangle $a\rightarrow b\rightarrow c$ with $a\in \N,c\in \N ^{\perp}$.

\item[(3)] The canonical functor $\D \rightarrow \D /\N$ has a right adjoint.

\item[(4)] The composite $\N ^{\perp}\hookrightarrow \D \rightarrow \D/\N$ is a triangle equivalence.

\end{itemize}

\end{proposition}

%--------------------------------------------------------------
\subsection{Semiorthogonal decompositions}\label{subsec:sod}
%--------------------------------------------------------------


If a full triangulated category $\N\subset \D$ is right admissible then by Lemma \ref{lem:notorious}, every object $x\in \D$ fits into a unique 
(up to unique isomorphism) distinguished triangle
\begin{equation}
\dots \longrightarrow  Y\longrightarrow X\longrightarrow Z\longrightarrow Y[1]\rightarrow \dots
\end{equation}
with $Y\in \N$ and $Z\in \N^{\perp}$. One then
says that there is a semiorthogonal decomposition of $\D$ into
the subcategories $(\N^{\perp}, \ \N)$. More generally,
assume given a sequence of full triangulated subcategories $\N _1,\dots,\N_n \subset \D$. Denote $\langle \N _1,\dots,\N _n\rangle:= \hull (\N _1,\dots,\N _n)$, the triangulated subcategory of $\D$ generated by $\N _1,\dots,\N _n$


\begin{definition}\label{def:semdecomposition-filtration}
{A semiorthogonal decomposition of a triangulated category $\D$ is a sequence of full triangulated subcategories $( \N _1,\dots,\N _n)$ in $\D$ such that $\N _i\subsetq \N_j^{\perp}$ for $1\leq i < j\leq n$ and for every object $X\in \D$ there exists a chain of morphisms in $\D$,
\[
\xymatrix@C=.5em{
0_{\ } \ar@{=}[r] & X_n \ar[rrrr] &&&& X_{n-1} \ar[rrrr] \ar[dll] &&&& X_{n-2}
\ar[rr] \ar[dll] && \ldots \ar[rr] && X_{1}
\ar[rrrr] &&&& X_0 \ar[dll] \ar@{=}[r] &  X_{\ } \\
&&& A_{n} \ar@{->}^{[1]}[ull] &&&& A_{n-1} \ar@{->}^{[1]}[ull] &&&&&&&& A_1\ar@{->}^{[1]}[ull] 
}
\]
such that a cone $A_k$ of the morphism $X_k\rightarrow X_{k-1}$ belongs to $\N _k$ for $k=1,\dots ,n$.}
\end{definition}

%--------------------------------------------------------------
\subsection{Mutations}\label{subsec:mutations}
%--------------------------------------------------------------

Let $\D$ be a triangulated category and assume $\D$ admits a semiorthogonal decomposition 
$\D = \langle \A, \B \rangle$ with $\A$ and $\B$ admissible.

\begin{definition}\label{def:left-right_mutation}
The left mutation of $\B$ through $\A$ is defined to be ${\bf L}_{\A}(\B): = \A ^{\perp}$. The 
right mutation of $\A$ through $\B$ is defined to be ${\bf R}_{\B}(\A): = {^{\perp}\B}$.
\end{definition}
One obtains semiorthogonal decompositions $\D = \langle {\bf L}_{\A}(\B), \A\rangle$ and 
$\D = \langle \A , {\bf R}_{\B}(\A)\rangle$.

\begin{lemma}\cite[Lemma 2.7]{Kuzrat}\label{lem:mutations_made_explicit}
There are equivalences ${\bf L}_{\A} :\B\simeq \D/\A\simeq {\bf L}_{\A}(\B)$ and ${\bf R}_{\A} :\A\simeq \D/\B\simeq {\bf R}_{\A}(\B)$.
\end{lemma}

\begin{definition}\label{def:left_and_right_dual_decompositions}
Let $\D = \langle \A _1,\dots,\A_n\rangle$ be a semiorthogonal decomposition with the $\A_i$ admissible. The left dual semiorthogonal decomposition  $\D = \langle \B _n,\dots,\B_1\rangle$ is defined by 
\begin{equation}
\B _i : = {\bf L}_{\A _1}{\bf L}_{\A _2}\dots {\bf L}_{\A _{i-1}}\A _i = {\bf L}_{\langle \A _1,\dots,\A_{i-1}\rangle}\A _i , \quad 1\leq i \leq n .
\end{equation}
The right dual semiorthogonal decomposition $\D = \langle \C _n,\dots,\C_1\rangle$ is defined by 
\begin{equation}
\C _i : = {\bf R}_{\A _n}{\bf R}_{\A _{n-1}}\dots {\bf R}_{\A _{i+1}}\A _i = {\bf R}_{\langle \A _{i+1},\dots,\A_{n}\rangle}\A _i , \quad 1\leq i \leq n .
\end{equation}


\end{definition}

\begin{lemma}\cite[Lemma 2.10]{Kuzrat}\label{lem:mutations_of_completely_orthogonal_subcategories}
Let $\D = \langle \A _1,\dots,\A_n\rangle$ be a semiorthogonal decomposition such that the components $\A _k$ and $\A _{k+1}$ are completely orthogonal, i.e., $\Hom _{\D}(\A _k,\A _{k+1}) = 0$ and $\Hom _{\D}(\A _{k+1},\A _k) = 0$. Then 
\begin{equation}
{\bf L}_{\A_k}\A _{k+1} = \A _{k+1} \qquad \qquad and  \qquad \qquad {\bf R}_{\A_{k+1}}\A _k = \A _k,
\end{equation}
and both the left mutation of $\A _{k+1}$ through ${\A_k}$ and the right mutation of $\A_k$ through $\A _{k+1}$ boil down to a permutation and
\begin{equation}
\D = \langle \A _1,\dots,\A _{k-1},\A _{k+1},\A _k, \A _{k+2},\dots \A_n\rangle
\end{equation}
is the resulting semiorthogonal decomposition of $\D$.

\end{lemma}

Given an object $F \in \D$, we also record explicit descriptions of 
left and right mutations of $F$ through an admissible admissible subcategory $i:\A\hookrightarrow \D$. Let $i^{\ast}$ and $i^{!}$ denote left and right adjoint to functors to $i$, respectively. Then the left mutation ${\bf L}_{\A}(F)$ the (reps. the right mutations ${\bf R}_{\A}(F)$ of $F$ through $\A$) are given by: 
\begin{equation}\label{eq:left-right_mutations_of_an_object}
{\bf L}_{\A}(F): = {\sf Cone}(ii^{!}(F)\rightarrow F), \qquad \qquad {\bf R}_{\A}(F): = {\sf Cone}(F\rightarrow ii^{\ast}(F))[-1].
\end{equation}


%--------------------------------------------------------------
\subsection{Exceptional collections}\label{subsec:FEC}
%--------------------------------------------------------------

Let $\sk$ be a field or $\Z$. Let $\kMod$ denote the category of $\sk$-modules, and $\kmod$ the subcategory of finitely generated $\sk$-modules.
Exceptional collections in $\sf k$-linear triangulated categories are a special case of semiorthogonal decompositions with each component of the decomposition being equivalent to $\Dd ^b({\kmod})$. The above properties of mutations thus specialize to this special case.

\begin{definition}\label{def:exceptcollection}Let $\D$ be a $\sf k$-linear triangulated category.

An object $E \in \D$ of  $\D$ is said to be exceptional if there is an isomorphism of graded $\sf k$-algebras 
\begin{equation}
\Hom _{\D}^{\bullet}(E,E) = {\sf k}.
\end{equation}
A collection of exceptional objects $(E_0,\dots,E_n)$ in $\D$ is called 
exceptional if for $1 \leq i < j \leq n$ one has
\begin{equation}
\Hom _{\D}^{\bullet}(E_j,E_i) = 0.
\end{equation}
The collection $(E_0,\dots,E_n)$ in $\D$ is said to be {\it full} if 
$\hull(E_0,\dots,E_n)^{\perp} = 0.$
\end{definition}
%--------------------------------------------------------------
\subsection{Admissible subcategories from exceptional collections}
%--------------------------------------------------------------
Assume given an exceptional collection $E_0,\dots,E_n$ in  $\D$, and assume  $$\RHom_\D(-,-):\D\times \D\to \Dd^b(\kmod).$$
One proves \cite[Theorem 3.2]{Bondal} that the subcategory $\hull(E_0,\dots,E_n)\subsetq \D$ is admissible in $\D$.
If the collection is full, then ${\D}
= \hull(E_0,\dots,E_n)$.

\

If $\A \subset \D$ is generated by an exceptional object $E$, then by 
 (\ref{eq:left-right_mutations_of_an_object}) the left and right mutations of an object $F\in \D$ through $\A$ are given by the following distinguished triangles:
\begin{equation}
\RHom_\D(E,F)\otimes E\rightarrow F\rightarrow {\bf L}_{\langle E\rangle}(F), \qquad 
{\bf R}_{\langle E\rangle}(F)\rightarrow F\rightarrow \RHom_\D(F,E)^{\ast}\otimes E .
\end{equation}
%--------------------------------------------------------------
\subsection{Dual exceptional collections}\label{subsec:dualFEC}
%--------------------------------------------------------------
Let $\sk$ be a field.
\begin{definition}\label{def:left-right_dual_exceptional_collections}
Let $X$ be a smooth projective variety, and assume given a full exceptional collection $(E_0,\dots,E_n)$ in $\Dd ^{b}(X)$. The right dual exceptional collection $(F_n,\dots, F_0)$ to $(E_0,\dots,E_n)$ is defined as 
\begin{equation}
F_i: = {\bf R}_{\langle E_{i+1},\dots,E_{n} \rangle}E_i, \quad {\rm for} \quad 1\leq i\leq n.
\end{equation}
The left dual exceptional collection $(G_n,\dots, G_0)$ to $(E_0,\dots,E_n)$ is defined as 
\begin{equation}
G_i: = {\bf L}_{\langle E_{1},\dots,E_{i-1} \rangle}E_i, \quad {\rm for} \quad 1\leq i\leq n.
\end{equation}
\end{definition}


\begin{proposition}\cite[Proposition 2.15]{Ef}\label{prop:dual_exc_coll_characterization}
Let $(E_0,\dots,E_n)$ be a full exceptional collection  in a triangulated category $\D$. The left dual exceptional collection $\langle F_n, \dots, F_0\rangle$ is uniquely determined by the following property:
\begin{equation}
\Hom _{\D}(E_i,F _j[l]) = \left\{
   \begin{array}{l}
    {\sf k}, \quad {\rm for} \quad l=0, \  i=j,\\
    0, \quad {\rm otherwise}. \\
   \end{array}
  \right.
\end{equation}
Similarly, the right dual exceptional collection $\langle G _n,\dots, G _0\rangle $ is uniquely determined by the following property:
\begin{equation}
\Hom _{\D}(G_i,E_j[l]) = \left\{
   \begin{array}{l}
    {\sf k}, \quad {\rm for} \quad l=0, \  i=j,\\
    0, \quad {\rm otherwise}. \\
   \end{array}
  \right.
\end{equation}
\end{proposition}


%--------------------------------------------------------------
\section{Generating the categories $\rep({\fb})$ and $\Dd^b (\rep ({\fb}))$}\label{sec:B-generation}
%--------------------------------------------------------------

%--------------------------------------------------------------
\subsection{The categories $\Dd^b (\rep ({\fb}))$ and $\Dd^b (\rep ( \bg))$}%-linear category}
\label{subsec:bounded-triang_cat}
%--------------------------------------------------------------
Let $\sk$ be a field or $\Z$.

We recall that $\Rep(\fb)$ (resp., $\rep(\fb)$) denotes the abelian category 
of rational $\fb$-modules 
(resp., the abelian subcategory of rational $\fb$-modules which are finitely generated over $\sk$). 
We let denote $\Dd(\Rep ({\fb}))$ (resp., $\Dd(\Rep ( \bg))$) the unbounded derived category of $\Rep(\fb)$ (resp., of $\Rep ( \bg)$), 
and $\Dd^b (\rep ({\fb}))$ (resp., $\Dd^b (\rep ( \bg))$) the bounded derived category of the smaller category $\rep(\fb)$ (resp., of $\rep ( \bg)$).

Let $\rep_\fr( \fb)$ denote the full subcategory of $\rep( \fb)$ consisting of the representations that are free over $\sk$.
Define $\rep_\fr( \bg)$
similarly.

We have the exact bifunctors:
$$-\;\otimes_\sk-:\rep_\fr( \fb)\times \rep_\fr( \fb)\to \rep_\fr( \fb)$$ 
and (the internal $\Hom _\sk$ on $\rep_\fr( \fb)$):
$$\Hom_\sk(-,-):\rep_\fr( \fb)\times \rep_\fr( \fb)\to \rep_\fr( \fb).$$ 
Put $M^*=\Hom_\sk(M,\sk)$,
for $M\in\rep_\fr( \fb)$. By \cite[I, Section 2.7, (4)]{Jan}, there is a natural $\fb$-module structure on $M^*$.

\begin{proposition}[Resolution property]\label{prop:resolution}
    Let $ H$ be a flat affine group scheme over a PID $\sk$.
    Then for every finitely generated $ H$-module $N$ there is an exact sequence $$0\to L\to M\to N\to0$$ with the
    $ H$-modules $L$, $M$, finitely generated and free over $\sk$. 
\end{proposition}
\begin{proof}
    This is a special case of \cite[Proposition 3]{Serre}.
\end{proof}

Let $\sk$ be a field or $\Z$
again.

The Proposition implies in standard fashion that every
bounded complex in $\rep( \fb)$ is quasi-isomorphic to a bounded complex with objects in $\rep_\fr( \fb)$.
The above functors can thus be extended as exact bifunctors to the bounded derived category $\Dd ^b(\rep( \fb))$: 
$$
(-)\otimes_\sk (-): \Dd ^b(\rep( \fb))\otimes \Dd ^{b}(\rep( \fb))\rightarrow \Dd ^b(\rep( \fb))
$$
and 
$$\Hom_\sk(-,-):\Dd ^b(\rep( \fb))\times \Dd ^b(\rep( \fb))\to \Dd ^b(\rep( \fb)).$$ 
In other words, we are using that by the dual of \cite[Theorem 10.22, Remark 10.23]{Buehler},
 the bounded derived category of the abelian category $\rep(\fb)$ is equivalent to the bounded 
derived category of its exact subcategory $\rep_\fr(\fb)$.

We often write $\otimes_\sk $ as $\otimes$. We put $M^*=\Hom_\sk(M,\sk)$.
By \cite[I, Section 2.7, (5)]{Jan}, there is an isomorphism $\Hom _\sk(M, N)=M^{\ast}\otimes N$ for $M,N\in \Dd ^b(\rep( \fb))$. There is $(\Hom _\sk,\otimes _\sk)$-adjunction:
$$
\Hom _\sk (-\otimes _\sk-,-)=\Hom _\sk (-,\Hom _\sk(-,-)).
$$
Let $\Dd_{\rep ({\fb})}(\Rep(\fb))$ denote the derived category of complexes of $\Rep(\fb)$ whose cohomology lie in $\rep ({\fb})$.
\begin{lemma}
There is an equivalence of triangulated categories $\Dd^b (\rep ({\fb}))=\Dd ^b_{\rep ({\fb})}(\Rep(\fb))$. 
\end{lemma}

\begin{proof}
This is a particular case of \cite[Appendix, Lemma A.3]{MR}.
\end{proof}

Let $\sk=\Z$.
\begin{corollary}\label{cor:Z left orthogonal}
Let $M, N\in \rep(\Gm)$.
\begin{enumerate}
\item $\Ext^i_ \Bm (M\otimes P(-e_v)^*,N\otimes Q(e_v))=\Ext^i_ \Gm (M,N) $ for all $i$.
\item If $w\not\leq v$ then $\Ext^i_ \Bm (M\otimes P(-e_v)^*,N\otimes Q(e_w))=0$ for all $i$.
\end{enumerate}
\end{corollary}
\begin{proof}We know this already for $M, N\in \rep_\fr(\Gm)$, cf. Remark \ref{rem:fr left orthogonal}.
By Proposition \ref{prop:resolution} the results follow.
\end{proof}
%{
%Let $\sk$ be a field.}
\subsection{The category $\Dd^b (\rep ({\fb}))$ as a $\Dd^b (\rep ( \bg))$-linear category}
\label{subsec:G-linear-triang_cat}

Let $\sk$ be a field or $\Z$.
The restriction functor $\res_{\fb}^{ \bg}: \Dd^+ (\Rep ( \bg))\rightarrow \Dd^+ (\Rep ({\fb}))$ is $t$-exact. 
Its right adjoint is the induction functor $\Rind_{\fb}^{ \bg}$ and the fact that $\res_{\fb}^{ \bg}$ gives a full faithful embedding is a consequence of the Generalized Tensor Identity \cite[I Proposition 4.8]{Jan} and the Kempf vanishing theorem, \cite[Theorem 1.2]{CPSWvdK}, \cite[B.3]{Jan} or Section \ref{subsec:Kempfvan}: 
it implies that $\Rind_{\fb}^{ \bg}\res_{\fb}^{ \bg}={\sf id}_{\Dd^+ (\Rep ( \bg))}$. 

\begin{proposition}
The functor $\Rind_{\fb}^{ \bg}: \Dd ^+(\Rep( \fb))\rightarrow \Dd ^+(\Rep( \bg))$ restricts to a functor $\Rind_{\fb}^{ \bg}: \Dd ^b(\rep( \fb))\rightarrow \Dd ^b(\rep( \bg))$. 
\end{proposition}

\begin{proof}
Set $X= \bg/ \fb$, $Y={\sf pt}$, $f:X\rightarrow Y$, $H= \bg$ in \cite[Proposition A.12]{MR}.
\end{proof}

It follows from the above that $\Dd ^b(\rep ( \bg))$ identifies with a right admissible subcategory of $\Dd^b(\rep ({\fb}))$. 
Both categories are monoidal and $\Dd^b(\rep ({\fb}))$ is a $\Dd ^b(\rep ( \bg))$-linear triangulated category. 
Denote $\Dd ^+_{ \bg}( \bg/ \fb):=\Dd ^+({\QCoh}^{ \bg}( \bg/ \fb))
%=\Dd ^+(\Rep( \fb))
$ and $\Dd ^b_{ \bg}( \bg/ \fb):=\Dd ^b({\sf Coh}^{ \bg}( \bg/ \fb))$.
% =\Dd ^b(\rep( \fb))$.

\begin{proposition}\label{prop:Hom-ind_factorisation}
Let ${\sf F}\in \Dd ^b(\rep( \fb))$ and $\Ff \in \Dd ^b_{ \bg}( \bg/ \fb)$ be the associated complex of equivariant sheaves on $ \bg/ \fb$. 
Then the functor $\RHom _{ \bg/\fb}(\Ff,-): \Dd ^+_{ \bg}( \bg/ \fb)\rightarrow \Dd ^+(\kMod)$ factors canonically through a functor $\Dd ^+_{ \bg}( \bg/ \fb)\rightarrow \Dd ^+(\Rep( \bg))$. 
\end{proposition}

\begin{proof}
Put $X= \bg/ \fb, H= \bg$ in \cite[Appendix, Corollary A.5]{MR}
\end{proof}

%{

The flag variety $ \bg/ \fb$ being proper and smooth, the above functor $\RHom _{ \bg/\fb}(\Ff,-)$ restricts to a functor $\Dd ^b_{ \bg}( \bg/ \fb)\rightarrow \Dd ^b(\kmod)$. 
Let $\Ff$, ${\mathcal G}\in \Dd ^b_{ \bg}( \bg/ \fb))$. Then $\RHom _{ \bg/\fb}(\Ff,\mathcal G)=\RGamma ( \bg/ \fb,(\Ff)^{\vee}\otimes {\mathcal G})$ where $(-)^{\vee}={\mathcal{R}\mathit{Hom}}(-,\Oo _{ \bg/\fb})$ is the Grothendieck duality on $ \bg/ \fb$. Let ${\sf F}\in \Dd^b(\rep (\fb ))$ be such that $\Ll (\sf F)=\Ff$. By \cite[Corollary 2.9]{CPSMackey}, $\RHom _{\bg/\fb}(\Ff,-)=\Rind _{\fb}^{\bg}({\sf F}^{\ast}\otimes (-))$. Thus, the functor $\Dd ^+_{\bg}(\bg/\fb))\rightarrow \Dd ^+(\rep(\bg))$ from Proposition \ref{prop:Hom-ind_factorisation}, which $\RHom _{\bg/\fb}(\Ff,-)$ factors through, is then isomorphic to  $\Rind _{\fb}^{\bg}({\sf F}^{\ast}\otimes (-))$.
%}
\

Let ${\sf Inv}^{ \bg}$ denote the derived functor of invariants $\Dd ^+(\Rep( \bg))\rightarrow \Dd ^+(\kMod)$.

\begin{proposition}\label{prop:derived_B-cohomology}
Let ${\sf F},{\sf G}\in \Dd ^+(\rep( \fb))$ and $\Ff=\Ll ({\sf F})$, ${\mathcal G}=\Ll ({\sf G})\in \Dd ^+_{ \bg}( \bg/ \fb)$ be the associated equivariant complexes of coherent sheaves. There there is a canonical functorial isomorphism
$$
{\sf Inv}^{ \bg}\circ \RHom _{ \bg/\fb}(\Ff,\mathcal G)\xrightarrow{\sim} \RHom _{\Dd ^+(\Rep( \fb))}({\sf F},{\sf G}).
$$
\end{proposition}

\begin{proof}
Put $X= \bg/ \fb$ and $ H= \bg$ in \cite[Appendix, Proposition A.6]{MR}
\end{proof}

The above isomorphism restricts to an isomorphism of functors ${\sf Inv}^{ \bg}\circ \RHom _{ \bg/\fb}(\Ff,\mathcal G)\xrightarrow{\sim} \RHom _{\Dd ^b(\rep( \fb))}({\sf F},{\sf G}).$

%------------------------------------------------------
\subsection{$ \bg$-linear semiorthogonal decompositions}\label{subsec:G-linear-sod}
%------------------------------------------------------

\begin{definition}\label{def:G-linear-cat}
A triangulated category $\D$ is called $ \bg$-linear if $\D$ is equipped with a monoidal action of $\Dd ^b(\rep( \bg))$, i.e. there is a bifunctor $\Dd ^b(\rep( \bg))\times \D\rightarrow \D$. Compare \cite[Section 2.7]{KuzHPD}.
\end{definition}

Our main concern is the category $\Dd ^b(\rep( \fb))$ which becomes a $ \bg$-linear triangulated category under the restriction functor $\res_{\fb}^{ \bg}:\Dd ^b(\rep( \bg))\rightarrow \Dd ^b(\rep( \fb))$.
Later we will also need $\bp$, so let us use $\bp$ instead of $\fb$, having in mind $\bp=\fb$ as an important case.

\begin{definition}\label{def:G-linear-SOD}
A $ \bg$-linear semiorthogonal decomposition of $\Dd^b(\rep ( \bp))$ is a collection $\langle {\sf A}_1,\dots {\sf A}_k\rangle$ of full triangulated $ \bg$-linear subcategories such that the collection $\langle {\sf A}_1,\dots {\sf A}_k\rangle$ is a semiorthogonal decomposition of $\Dd^b(\rep ( \bp))$ in the sense of Definition \ref{def:semdecomposition-filtration}.
\end{definition}


The following proposition is just a variation of \cite[Lemma 2.7]{Kuzbasech}:

\begin{proposition}\label{prop:G-linear_sod}
A pair of $\Dd ^b(\rep ( \bg))$-linear subcategories ${\sf A},{\sf B}\subsetq \Dd ^{b}(\rep( \bp))$ is semiorthogonal (i.e.\  ${\sf A}\subsetq{\sf B}^\perp$) if and only if 
the equality $\Rind_{ \bp}^{ \bg}(N^{\ast}\otimes _k M)=0$ holds for all $M\in {\sf A}$, $N\in {\sf B}$.   
\end{proposition}

\begin{proof}

($\Leftarrow$). Assume given $M\in {\sf A}$, $N\in {\sf B}$ and assume $\Rind_{ \bp}^{ \bg}(N^{\ast}\otimes _k M)=0$. 
Then 
\begin{eqnarray*}&\RHom _{\Dd ^{b}(\rep( \bp))}(N,M)=\RHom _{\Dd ^{b}(\rep( \bp))}(\sk,N^*\otimes M)=\\
&\RHom _{\Dd ^{b}(\rep(\bg))}(\sk,\Rind_{ \bp}^{ \bg}(N^{\ast}\otimes _k M))=0.
\end{eqnarray*}

($\Rightarrow$). Let $M\in {\sf A}$, $N\in {\sf B}$ and 
$\RHom _{\Dd ^{b}(\rep( \bp))}(N,M)=0$ for all such $M$, $N$.
Let be $L$ an arbitrary object of $\Dd ^{b}(\rep( \bg))$. Then 
$$
\RHom _{\Dd ^{b}(\rep( \bg))}(L,\Rind_{ \bp}^{ \bg}(N^{\ast}\otimes _{\sk}M))=\RHom _{\Dd ^{b}(\rep( \bp))}(\res _{ \bp}^{ \bg}L,N^{\ast}\otimes _{\sk}M)
$$
and 
$$
\RHom _{\Dd ^{b}(\rep( \bp))}(\res _{ \bp}^{ \bg}L,N^{\ast}\otimes _{\sk}M)
=\RHom _{\Dd ^{b}(\rep( \bp))}(N\otimes _{\sk}\res _{ \bp}^{ \bg}L,M)=0;
$$
the latter equality holds since the subcategory $\sf B$ is $ \bg$-linear and hence stable under tensoring with objects of $\Dd ^{b}(\rep( \bg))$. 
It follows that $\RHom _{\Dd ^{b}(\rep( \bg))}(L,\Rind_{ \bp}^{ \bg}(N^{\ast}\otimes _{\sk}M))=0$ for an arbitrary object $L\in \Dd ^{b}(\rep( \bg))$; hence, $\Rind_{ \bp}^{ \bg}(N^{\ast}\otimes _{\sk}M)=0$.
\end{proof}
Let us simplify notation and write $\RHom _{ \Dd ^b(\rep( \bp))}$ as $\RHom _{ \bp}$,  and $\RHom _{ \Dd ^b(\rep( \bg))}$ as $\RHom _{ \bg}$.
\begin{definition}
A functor $\Phi: \Dd (\rep( \bp))\rightarrow \Dd (\rep( \bp))$ is called $ \bg$-linear if for all $M \in \Dd (\rep( \bp))$, $N \in \Dd ^b(\rep( \bg))$ there are given bifunctorial isomorphisms 
$$\Phi (\res_{ \bp}^{ \bg}(N)\otimes M)=\res_{ \bp}^{ \bg}(N)\otimes \Phi (M).$$
\end{definition}


\begin{proposition}\label{prop:G-linear_exceptional}
\
\begin{enumerate}
\item Assume given a fully faithful $ \bg$-linear functor $\Phi: \Dd ^b(\rep( \bg))\rightarrow \Dd ^b(\rep( \bp))$. Then $\Phi$ is isomorphic to the functor $\Phi _E(-)=(-)\otimes E$ where $E\in \Dd ^b(\rep( \bp))$ such that $\RHom _{ \bp}(E,E)=\sk$.

\item Each object $E\in \Dd ^b(\rep( \bp))$ such that $\Rind_{ \bp}^{ \bg}(E^{\ast}\otimes E)=\sk$ gives a fully faithful $ \bg$-linear functor $\Phi: \Dd ^b(\rep( \bg))\rightarrow \Dd ^b(\rep( \bp))$. 

\item Under the assumptions of (1), the object $E$ satisfies $\Rind_{ \bp}^{ \bg}(E^{\ast}\otimes E)=\sk$.
\end{enumerate}

\end{proposition}
\begin{proof}
(1) Take $E=\Phi(\sk)$.

(2) For $M,N\in \Dd ^b(\rep( \bg))$ we have
\begin{eqnarray*}& \RHom _{ \bp}(M\otimes E,N\otimes E)=\RHom_\bp(M\otimes N^\ast,E^\ast\otimes E)=\\
&\RHom_\bg(M\otimes N^\ast,\Rind_\bp^\bg( E^\ast\otimes E))=\RHom_\bg(M\otimes N^\ast,\sk)=
\RHom_\bg(M,N).
\end{eqnarray*}

(3) Under the assumptions of (1)
\begin{eqnarray*}&\RHom _{ \bg}( M, \sk)=\RHom_\bp(M\otimes E, E)=\\&\RHom_\bp(M, E^\ast\otimes E)=\RHom_\bg(M,\Rind_\bp^\bg( E^\ast\otimes E)).
\end{eqnarray*}
for $M\in \Dd^b(\rep(\bg))$. By the Yoneda Lemma it follows that $\Rind_{ \bp}^{ \bg}(E^{\ast}\otimes E)=\sk$.
\end{proof}
\begin{remark}\label{rem:ind-to-B-exceptional}
See Section \ref{subsec:cohdesc} for the precise relation between the two vanishing conditions ${\rm R}^{i}\ind_{\fb}^{ \bg}(E^{\ast}\otimes E)=0$ for $i>0$ and $\RHom ^{i} _{\fb}(E,E)=0$ for $i>0$.
\end{remark}


%--------------------------------------------------------------
\subsection{Generating $\Dd ({\QCoh}(\bg/\bp))$}\label{subsec:generatingG/P}
%--------------------------------------------------------------

\begin{definition}
Let $\D$ be a compactly generated triangulated category. A set $S$ 
of compact objects of $\D$ is called a generating set if $\Hom _{\D}(S,X)=0$ implies $X=0$ and $S$ is closed under the shift functor, i.e. $S =S[1]$.
\end{definition}

\begin{proposition}\label{prop:ample_line_bundle_gen_set}
Let $X$ be a quasi--compact, separated scheme, and $\Ll$ be an ample line bundle on $X$. Then the set $\langle \Ll ^{\otimes m}[n]\rangle, m,n\in \mathbb Z$ is a generating set for $\Dd ({\QCoh}(X))$.
\end{proposition}

\begin{proof}
See \cite[Examples 1.10 and 1.11]{Neem}.
\end{proof}

\begin{corollary}\label{cor:line_bundles-generation}
The set of line bundles $\langle \Ll _{\lambda}[n],\lambda \in X(\bt)_+, n\in \mathbb Z\rangle$ is a generating set for $\Dd({\QCoh}( \bg/ \bp))$.
\end{corollary}

\begin{corollary}\label{cor:generatingGmodB}
The smallest thick full triangulated subcategory of $\Dd^b( \bg/ \bp)$ containing $\langle \Ll _{\lambda}[n],\lambda \in X(\bt)_+, n\in \mathbb Z\rangle$ is $\Dd^b( \bg/ \bp)$.
\end{corollary}

\begin{proof}
Follows from \cite[Corollary 2.3]{Neem}.
\end{proof}


The goal of the rest of this section---and the two sections after it---is to construct a collection of objects $X_p\in \Dd ^b(\rep( \fb))$, $p\in W$, each satisfying the condition in (2) of Proposition \ref{prop:G-linear_exceptional}; by part (1) of the Proposition
each of those objects will also be $\fb$-exceptional. Furthermore, the collection of objects $X_p$ will form a $ \bg$-linear semiorthogonal decomposition of $\Dd^b(\rep ({\fb}))$; the ultimate statement is Theorem \ref{th:semi-orthogonal}.


%--------------------------------------------------------------
\subsection{Generating $\Dd^b (\rep ( \bg))$}\label{subsec:generatingrepG}
%--------------------------------------------------------------

\begin{proposition}\label{prop:stronggeneratingrepG}Let $\sk$ be a field or $\Z$.
The smallest strictly full triangulated subcategory of $\Dd^b (\rep ( \bg))$ that contains the set of modules $\nabla _{\lambda}$,
$\lambda \in X(\bt)_+$, is the whole $\Dd^b (\rep ( \bg))$.
\end{proposition}

\begin{proof}First we show that $\rep(\bg)$ lies in the subcategory.
We choose a real valued additive injective height function $\hgt$ on the weight lattice which is positive on positive roots.
Consider a nonzero $M\in\rep( \bg)$. Say $\mu$ is its highest weight with respect to $\hgt$. 
Thus $\hgt(\mu)\geq0$. Assume that all representations with a smaller highest weight are in the subcategory. We have a map $\Delta_\mu\to \nabla_\mu$ whose kernel and cokernel have
lower weights, so  $\Delta_\mu$ lies in the subcategory.
Let $M_{\mu,\mathrm {triv}}$ denote the weight space $M_{\mu}$ provided with a trivial $ \bg$ action.
By \cite[Proposition 21]{FvdK} there is a natural map 
$\Delta_\mu\otimes_\sk M_{\mu,\mathrm {triv}}\to M$ whose kernel and cokernel have lower weights.  
Now notice that $\Delta_\mu\otimes_\sk M_{\mu,\mathrm {triv}}$ lies in the subcategory. So $M$ does too.
So $\rep(\bg)$ lies in the subcategory. 
Now use that every object of $\Dd^b(\rep(\bg))$  is quasi-isomorphic to a bounded complex in $\rep( \bg)$ and that bounded complexes are repeated cones of pure objects. (cf.\ ``Stupid  truncations'' \cite[2.5]{Kuzbasech}.) 
\end{proof}

\begin{proposition}\label{prop:G linear hull}
Let $S$ be a set of objects of $\Dd^b(\rep(\fb))$. 

Then $\hull(\{\nabla_\nu\otimes M\mid \nu\in X( \bt)_+,M\in S\})$ is $\bg$-linear.
\end{proposition}
\begin{proof}
By Proposition \ref{prop:stronggeneratingrepG} we have $\hull(\{\nabla_\nu\mid \nu\in X( \bt)_+\})= \Dd^b (\rep ( \bg))$. 
So $\hull(\{\nabla_\nu\otimes M\mid \nu\in X( \bt)_+,M\in S\})$ equals $\hull(\{\Dd^b (\rep ( \bg))\otimes M\mid M\in S\})$.
\end{proof}
%--------------------------------------------------------------
\subsection{Generating $\Dd^b (\rep ({\fb}))$}\label{subsec:D^brepB-generation}
%--------------------------------------------------------------

Let $\sk$ be a field or $\Z$. 

\begin{theorem}\label{th:derived_generation}
Given a $p\in W$, the triangulated hull in $\Dd(\Rep(\fb))$ of the two categories
\begin{eqnarray*}
&\hull(\{\nabla_\lambda\otimes Q(e_v)\}_{v\succ p,\lambda\in X( \bt)_+}),\\
&\hull(\{\nabla_\lambda\otimes P(-e_v)^*\}_{v\preceq p,\lambda\in X( \bt)_+}) 
\end{eqnarray*}
is $\Dd^b (\rep ({\fb}))$.
\end{theorem}

\begin{proof}
This will take the rest of the section. We will categorify the theorem of Steinberg that says that
the $[\sk_{e_v}]$ generate $R(\fb)$ as an $R(\bg)$-module and we will apply the same reasoning with a few $\sk_{e_v}$'s replaced with $Q(e_v)$ or $P(-e_v)^*$. Our arguments are similar to the proof of \cite[Theorem 2]{Ana}.
\end{proof}

%--------------------------------------------------------------
\subsection{Generating $\rep({\fb})$}\label{subsec:repB-generation}
%--------------------------------------------------------------

  We say that a full subcategory of an abelian category has the 2 out of 3 property if,
  whenever $0\to N_1\to N_2 \to N_3 \to 0$ is exact and two of the $N_i$ are in the subcategory, then so is the third.
  
  \subsection{The set $\{M_v\}_{v\in W}$}
  
   Let us be given a set $\{M_v\}_{v\in W}$ of objects of $\rep_\fr({\fb})$  with the following properties.
  The multiplicity of the weight $e_v$ in $M_v$ is one. Every weight $\lambda$ of $M_v$ satisfies $(\lambda,\lambda)\leq (e_v,e_v)$.
  If $\lambda $ is a weight of $M_v$ with $(\lambda,\lambda)= (e_v,e_v)$, then  $\lambda$ is a weight of $P(-e_v)^*$. So all weights $\lambda$
  of $M_v$ precede $e_v$ in the antipodal excellent order, notation $\lambda\leq_ae_v$.
Examples of possible choices of $M_v$ are $\sk_{e_v}$, $Q(e_v)$, $P(-e_v)^*$.
  We use $e^\lambda$ as another notation for $\sk_\lambda$. 
  
  \begin{theorem}[Generation]\label{th:generation}
  The smallest strictly full additive subcategory  that 
  \begin{itemize}
  \item contains the $M_v$,  
  \item has the
  2 out of 3 property and 
  \item contains with every $e^\lambda$ and every fundamental representation $\nabla(\omega_i)$ also
   $e^\lambda\otimes\nabla(\omega_i)$,
  \end{itemize}
   is the category $\rep({\fb})$ of finite dimensional ${\fb}$-modules.
  \end{theorem}
  
We first prove a lemma and two propositions.
   \begin{lemma}\label{lem:aed}
    Let $\lambda\in X(\bt)$, $\alpha\in \Pi$.
    Then $s_\alpha\lambda\leq_e\lambda$ if and only if $s_\alpha\lambda\geq_d\lambda$.
    
    Recall that $s_\alpha\lambda\leq_e\lambda$ if and only if $s_\alpha\lambda\geq_a\lambda$ by Lemma \ref{lem:ea}.
\end{lemma}
\begin{proof}
    Write $\lambda=w\nu$ with $\nu$
    dominant and $w$ minimal. We have $s_\alpha w>w$ or $s_\alpha w<w$. If $s_\alpha w>w$, then $\lambda$ is a weight of $P(s_\alpha\lambda)$ and the socle of $P(s_\alpha\lambda)$, which has weight $s_\alpha\lambda$ by \cite[Lemma 2.2.9]{WvdKTata}, must lie in the $\fb$-submodule of $P(s_\alpha\lambda)$
    generated by the weight space of weight $\lambda$. This implies
$s_\alpha\lambda\leq_d\lambda$. The other case is similar.
\end{proof}
   \begin{proposition}\label{prop:touch}
  Let $\lambda$ be a weight in the $W$ orbit of the dominant weight $\varpi$. 
  Assume $\alpha_i^\vee(\varpi)\leq1$ for all
  simple $\alpha_i$. 
  Take $w$ minimal so that $\lambda=w\varpi$. Let there be an $\alpha_j$  such that $\alpha_j^\vee(\varpi)=1$ and 
  $w\leq w s_j$. Let $\mu\in w(\varpi-\omega_j)+W\omega_j$. Then $(\mu,\mu)\leq (\varpi,\varpi)$ and if $(\mu,\mu)=(\varpi,\varpi)$,
  then there is $v\geq w$ with $\mu=v\varpi$.
  So $\mu\leq_a\lambda$.
  \end{proposition}
  \begin{proof}
  Put $\tau=\varpi-\omega_j$. Let $W_I$, $W_J$, $W_K$ be the stabilizers in $W$ of $\tau$, $\omega_j$, $\varpi$ respectively.
  Choose $z$ minimal so that $w^{-1}\mu=\tau+z\omega_j$. Assume $(\mu,\mu)\geq(\varpi,\varpi)$.
 % The longest weights of $\nabla(\tau)\otimes\nabla(\omega_j)$ are all in $W\varpi$, so $\mu\in W\varpi$ or $(\mu,\mu)<(\varpi,\varpi)$.
%  Say $\mu\in W\varpi$.
%  As $\tau+z\omega_j$ has the same length as $\tau+\omega_j$, we must have $(\tau,\omega_j-z\omega_j)=0$.
  Choose a reduced expression $s_ks_{k-1}\cdots s_1$ for $z$ and put $z_0=\id$, $z_i=s_iz_{i-1}$. 
  As $z_{i+1}\omega_j>_ez_i\omega_j$, we have by Lemma \ref{lem:aed} that
the path $z_0\omega_j,\cdots, z_k\omega_j$ from $\omega_j$ to $z\omega_j$ is strictly descending for $\leq_d$.
  Along the path $(\tau,z_i\omega_j)$ 
  can only go down, but $(\tau,u\omega_j)<(\tau,v\omega_j)$ implies 
  $(\tau+u\omega_j,\tau+u\omega_j)<(\tau+v\omega_j,\tau+v\omega_j)$.
  So the path must consist of steps in directions perpendicular
  to $\tau$. So $z\in W_I$.
  Now $w$ is minimal in its coset $wW_K$, and moreover we have $w\leq w s_j$. That makes it minimal in $wW_I$ also.
  This shows $wz\geq w$. Now note that $\mu=w(\tau+z\omega_j)=w(z\tau+z\omega_j)=wz\varpi$.
  \end{proof}
  
  \begin{proposition}\label{prop:fartouch}
  Let $\lambda$ be a weight in the $W$ orbit of the dominant weight $\varpi$. 
  Assume $\alpha_j^\vee(\varpi)>1$ for some
  simple $\alpha_j$. Take $w$ minimal so that $\lambda=w\varpi$.  Let $\mu\in w(\varpi-\omega_j)+W\omega_j$. 
  Then $(\mu,\mu)\leq (\varpi,\varpi)$ and if $(\mu,\mu)=(\varpi,\varpi)$,
  then there is $v\geq w$ with $\mu=v\varpi$.  So $\mu\leq_a\lambda$.
  \end{proposition}
  \begin{proof}
  Put $\tau=\varpi-\omega_j$. Let $W_I$, $W_J$  be the stabilizers in $W$ of $\tau$, $\omega_j$ respectively.
  Observe that $W_I$ is also the stabilizer $W_K$ of $\varpi$.
  Choose $z$ minimal so that $w^{-1}\mu=\tau+z\omega_j$.
  %The longest weights of $\nabla(\tau)\otimes\nabla(\omega_j)$ are all in $W\varpi$, so $\mu\in W\varpi$ or $(\mu,\mu)<(\varpi,\varpi)$.
 % Say $\mu\in W\varpi$.
Assume $(\mu,\mu)\geq(\varpi,\varpi)$. 
%As $\tau+z\omega_j$ has the same length as $\tau+\omega_j$, we must have $(\tau,\omega_j-z\omega_j)=0$.
  Along the path $z_0\omega_j,\cdots, z_k\omega_j$ from $\omega_j$ to $z\varpi_j$, given by an irreducible expression for $z$, the inner product $(\tau,z_i\omega_j)$ 
  can only go down, so the path must consist of steps in directions perpendicular
  to $\tau$. So $z\in W_I=W_K$.
  Now $w$ is minimal in its coset $wW_K=wW_I$.  This shows $wz\geq w$. Now note that $\mu=w(\tau+z\omega_j)=w(z\tau+z\omega_j)=wz\varpi$.
  \end{proof}       

  
  \begin{proof}[\rm \bf  Proof of Theorem \ref{th:generation}]
  Let $\varpi$ be dominant and let $\lambda$ be a weight in its $W$ orbit. Choose $w$ minimal so that 
  $\lambda=w\varpi$.
  Assume that for all $v\geq w$ with $v\varpi\neq \lambda$ the representation $e^{v\varpi}$ 
   is in the subcategory.
Assume also that for all weights $\mu$ with $(\mu,\mu)<(\lambda,\lambda)$ 
  the one dimensional representation $e^\mu$ of weight $\mu$ is in the subcategory. 
  In other words, assume that $e^\mu$  is in the subcategory for all $\mu$ with $\mu<_a\lambda$.
  
  We claim that then    $e^\lambda$ is in the subcategory.
  The theorem easily follows from the claim by induction along $<_a$.
  
  So let us prove the claim.
  
  There are several cases.
  
  If $\lambda$ is a Steinberg weight $e_v$, then one uses the given properties of $M_v$.
  
  So we may assume it is not a Steinberg weight.
   
  There are two cases.
  
  The first case is that $\alpha_i^\vee(\varpi)\leq 1$ for all $i$.
  As $\lambda$ is not a Steinberg weight there must be a $j$ with $\alpha_j^\vee(\varpi)=1$ but $w(\alpha_j)$ positive.
 Then $\ell(ws_j)>\ell(w)$ and we are in the situation of Proposition \ref{prop:touch}. 
 Let $\tau=\varpi-\omega_j$. Then $\tau$ is strictly shorter than $\lambda$, so $e^{w\tau}$ is in our subcategory, and therefore
  $N=e^{w\tau}\otimes \nabla(\omega_j)$ is in the subcategory. The weights $\mu$ of $N$  lie in the convex hull of $w\tau+W\omega_j$ and are either shorter than $\lambda$
  or they are of the form $v\varpi$ with $v\geq w$. That is $\mu\leq_a\lambda$.
  So $e^\mu$ is in the subcategory for all weights of $N$ different from $\lambda$ and $\lambda$ has multiplicity one in $N$. It follows that $e^\lambda$ is in the subcategory.
  
  The second case is that some $\alpha_j^\vee(\varpi)> 1$. Put $\tau=\varpi-\omega_j$. Again $\tau$ is strictly shorter than $\varpi$.
  By Proposition \ref{prop:fartouch}
  the weights  $\mu$ of $N=e^{w\tau}\otimes \nabla(\omega_j)$ are either shorter than $\lambda$
  or they are of the form $v\varpi$ with $v\geq w$.
 So $e^\mu$ is in the subcategory for all weights of $N$ different from $\lambda$ and $\lambda$ has multiplicity one in $N$. It follows that $e^\lambda$ is in the subcategory.
  \end{proof}

  \begin{corollary}\label{cor:if not Steinberg then tensor}
If $\sigma$ is not a Steinberg weight, then  
 there is a fundamental weight $\omega_j$ and a weight $\mu$ with $(\mu,\mu)<(\sigma,\sigma)$,
 so that the $\fb$-module $L=\nabla_{\omega_j}\otimes e^\mu$  has
$\sigma$ as a weight  of  multiplicity one, and so that all weights  $\nu$ of $L$ satisfy 
$\nu\leq_a\sigma$.
\end{corollary}
\begin{proof}
See the proof of Theorem \ref{th:generation}.
\end{proof}

\begin{corollary}\label{cor:Steinberg_alternative}Let $\sk$ be an algebraically closed field. Let $W$ be partitioned
arbitrarily into three subsets $W_1$, $W_2$, $W_3$. One gets a basis of $R({\fb})$ over $R( \bg )$ by taking
   $\{\sk_{e_v}\}$ for ${v\in W_1}$, $Q(e_v)$ for ${v\in W_2}$,  $P(-e_v)^*$ for ${v\in W_3}$. 
 \end{corollary}

 \begin{remark}
{
The results of this section give an alternative proof of Steinberg's theorem \cite{St} that the $\sk_{e_w}$ generate
$R({\fb})$ over $R( \bg )$. See also Corollary \ref{cor:ordered_generate} or \cite[Theorem 2]{Ana}. Our proof  assumes only that $\bg$ is split and $\sk$ is a field or $\Z$, while
Steinberg takes the field algebraically closed. But by \cite[\S 3]{Serre} this makes no difference.
 }
 \end{remark}
    
\begin{proof}[\rm \bf  Proof of Theorem \ref{th:derived_generation}]
By Theorem \ref{th:generation}, for any $p\in W$ the set
 $\{\nabla_\lambda\otimes Q(e_v)\}_{v\succ p,\lambda\in X( \bt)_+}\cup\{\nabla_\lambda\otimes P(-e_v)^*\}_{v\preceq p,\lambda\in X( \bt)_+}$ generates $\rep({\fb})$ as an abelian category in a specific manner, corresponding in the derived
 category with taking cones, shifts, or using the $\bg$-linear structure of Proposition \ref{prop:G linear hull}. Therefore the hull contains all pure objects (objects concentrated in one degree).
 Considering the canonical truncation of an object of $\Dd^b (\rep ({\fb}))$ and resolving each cohomology via pure objects,
 %the above generators of $\rep({\fb})$, 
 we get the statement. One could also use stupid truncations, cf.\ \cite[2.5]{Kuzbasech}.
\end{proof}


%--------------------------------------------------------------
\section{Construction of the objects $X_p$ and $Y_p$}\label{sec:X_p&Y_p}
%--------------------------------------------------------------
By Theorem \ref{th:derived_generation}, we have
\begin{equation}\label{eq:trianghullP-Q}
\hull( \{\nabla_\lambda\otimes Q(e_v)\}_{v\succ  p,\lambda\in X( \bt)_+}\cup\{\nabla_\lambda\otimes P(-e_v)^*\}_{v\preceq  p,\lambda\in X( \bt)_+})%_{\oplus} 
=\Dd ^b(\rep( \fb)).
\end{equation}
for all $p\in W$.

%--------------------------------------------------------------
\subsection{Cut at $p\in W$}\label{subsec:cut_at_p}
%--------------------------------------------------------------



Introduce the following notation: 
\begin{itemize}
\item ${\sf Q}_{\succeq p}:=\hull(\{\nabla_\lambda\otimes Q(e_v)\}_{v\succeq  p,\lambda\in X( \bt)_+})$




\item ${\sf Q}_{\succ p}:=\hull(\{\nabla_\lambda\otimes Q(e_v)\}_{v\succ  p,\lambda\in X( \bt)_+})$





\item ${\sf P}_{\preceq p}:=\hull(\{\nabla_\lambda\otimes P(-e_v)^*\}_{v\preceq  p,\lambda\in X( \bt)_+})$



\item ${\sf P}_{\prec p}:=\hull(\{\nabla_\lambda\otimes P(-e_v)^*\}_{v\prec  p,\lambda\in X( \bt)_+})$

\end{itemize}
\begin{remark}
    We will need to find exceptional objects to conclude that ${\sf Q}_{\succeq p}$, ${\sf P}_{\preceq p}$, ${\sf Q}_{\succeq p}\cap{\sf P}_{\preceq p}$ are admissible subcategories. That is what this section and the next one are about.
\end{remark} 
\begin{remark}
    Do not confuse ${\sf Q}_{\succ p}$
    with $\hull(\{\nabla_\lambda\otimes Q(e_v)\}_{e_v>_a  e_p,\lambda\in X( \bt)_+})$.
    The latter may seem more natural, as the $Q(\mu)$ belong with $>_a$, not $\succ$.
    The Steinberg weights
    show two faces, and this is essential. On the one hand they are indexed by the Weyl group and are thus
    ordered by the Bruhat order, on the other hand they can be characterised  in terms of $>_a$, see Remark \ref{rem:Steinberg from anti}.
\end{remark}

\ 

Denote $\D:=\hull({\sf Q}_{\succ p} \cup {\sf P}_{\preceq p})$. Then (\ref{eq:trianghullP-Q}) gives that $\D$ is just $\Dd ^b(\rep( \fb))$.
\begin{equation}\label{eq:Q>pPleqp-generation}
\D=\hull({\sf Q}_{\succ p} \cup {\sf P}_{\preceq p}).
\end{equation}
By Corollary \ref{cor:left orthogonal} and Corollary \ref{cor:Z left orthogonal}, for all $p\in W$:
\begin{equation}
{\sf Q}_{\succ p}\subsetq{\sf P}_{\preceq p}^{\perp}.
\end{equation}
Note that we are taking the right orthogonal in $\D$.

By Lemma \ref{lem:notorious}  we have a left adjoint of the inclusion of ${\sf P}_{\preceq p}^\perp$ into $\D$.
Let $j^*$ denote this left adjoint of the embedding functor ${\sf P}_{\preceq p}^\perp\hookrightarrow \D$. Then
we have 
\begin{equation}\label{eq:image is perp}
j^*(\D)={\sf P}_{\preceq p}^\perp,
\end{equation}
\begin{equation}
j^*({\sf P}_{\preceq p})=0,
\end{equation}
\begin{equation}
j^*({\sf Q}_{\succ p})= {\sf Q}_{\succ p}.
\end{equation}
Now consider the full subcategory of $\D$ whose objects are the $X$ with $j^*(X)\in {\sf Q}_{\succ p}$.
It is triangulated and contains both ${\sf P}_{\preceq p}$ and  ${\sf Q}_{\succ p}$. So it contains 
$\hull({\sf Q}_{\succ p} \cup {\sf P}_{\preceq p})$ and must be all of $\D$.
In other words, $j^*(\D)= {\sf Q}_{\succ p}$. Taken together with equation (\ref{eq:image is perp}) this shows
\begin{equation}\label{eq:P perp is Q}
{\sf P}_{\preceq p}^\perp= {\sf Q}_{\succ p}.
\end{equation}
Thus, by Definition \ref{def:semdecomposition-filtration}, $\D$ has a semiorthogonal decomposition:
\begin{equation}\label{eq:p-q-SOD-1}
\D=\langle \;{\sf Q}_{\succ p},{\sf P}_{\preceq p}\; \rangle .
\end{equation}
Then the inclusion ${\sf Q}_{\succ p}\hookrightarrow \D$ has a left adjoint and the inclusion ${\sf P}_{\preceq p}\hookrightarrow \D$ has a right adjoint. So, ${\sf Q}_{\succ p}$ is left admissible in $\D$ and ${\sf P}_{\preceq p}$ is right admissible in $\D$. 
Dually to (\ref{eq:P perp is Q}) we also have 
\begin{equation}
{\sf P}_{\preceq p}={}^\perp {\sf Q}_{\succ p}.
\end{equation}
Moreover: 
\begin{proposition}\label{prop:p-q-SOD-2}
\begin{equation}\label{eq:p-q-SOD-2}
\D=\langle {\sf Q}_{\succeq p},{\sf P}_{\prec p}\rangle 
\end{equation}
\end{proposition}

\begin{proof}
    If $p\neq\id$, then $p$ has a predecessor $p'$ for the total order $\prec$ on $W$ and ${\sf Q}_{\succeq p}={\sf Q}_{\succ p'}$,
    ${\sf P}_{\prec p}={\sf P}_{\preceq p'}$.
    And if $p=\id\in W$, then actually $Q(e_p)=P(-e_p)^*=\sk$.
\end{proof}
%--------------------------------------------------------------
\subsection{Defining $X_p$ and $Y_p$}\label{subsec:X_p-Y_p}
%--------------------------------------------------------------


Let $X_p$ be the image of $P(-e_p)^*$ under the left adjoint of the inclusion of ${\sf Q}_{\succeq p}$ into $\D$.
More precisely, denote $i^q_{\succeq p}:{\sf Q}_{\succeq p}\hookrightarrow \D$ the embedding functor and let ${i^q}_{\succeq p}^{\ \ast}$ be its left adjoint. Then
\begin{equation}\label{eq:def_of_X_p}
X_p:={i^q}_{\succeq p}^{\ \ast}(P(-e_p)^*).
\end{equation}
By definition, $X_p\in {\sf Q}_{\succeq p}$. 
We have an exact triangle 
\begin{equation}\label{eq:def_triang_X_p}
\dots \rightarrow P(-e_p)^*\rightarrow X_p\rightarrow {\sf cone}(P(-e_p)^*\rightarrow X_p)\rightarrow \dots 
\end{equation}
Note that ${\sf cone}(P(-e_p)^*\rightarrow X_p)\in{\sf P}_{\prec p}=\hull(\{\nabla_\lambda\otimes P(-e_v)^*\}_{v\prec  p,\lambda\in X( \bt)_+})$. Therefore, 
$$X_p\in\hull(  P(-e_p)^*\cup {\sf P}_{\prec p})\subsetq\hull(\{\nabla_\lambda\otimes P(-e_v)^*\}_{v\preceq  p,\lambda\in X( \bt)_+})={\sf P}_{\preceq p}.$$
So, 
\begin{equation}\label{eq:X in Q and P}
X_p\in {\sf Q}_{\succeq p}\cap {\sf P}_{\preceq p}.    
\end{equation}
Now let $Y_p$ be the image of $Q(e_p)$ under the 
right adjoint of the inclusion of
${\sf P}_{\preceq p}=\{\nabla_\lambda\otimes P(-e_v)^*\}_{v\preceq p,\lambda\in X( \bt)_+}$ into $\D$. More precisely, denote $i^p_{\preceq  p}:{\sf P}_{\preceq p}\hookrightarrow \D$ the embedding functor and let ${i^p}_{\preceq p}^{\ !}$ be its right adjoint. Then
\begin{equation}\label{eq:def_of_Y_p}
Y_p:={i^p}_{\preceq p}^{\ !}(Q(e_p)).
\end{equation}
By definition, $Y_p\in {\sf P}_{\preceq p}$. We have an exact triangle 
\begin{equation}\label{eq:def_triang_Y_p}
\dots \rightarrow Y_p\rightarrow Q(e_p)\rightarrow {\sf cone}(Y_p\rightarrow Q(e_p))\rightarrow \dots 
\end{equation}
Note that ${\sf cone}(Y_p\rightarrow Q(e_p))\in {\sf P}_{\preceq p}^{\perp}=
{\sf Q}_{\succ p}$. Therefore, 
%:=\langle\{\nabla_\lambda\otimes Q(e_v)\}_{v\succ  p,\lambda\in X( \bt)_+}\rangle$. 
$$Y_p\in \hull( Q(e_p)\cup {\sf Q}_{\succ p})\subsetq\hull(\{\nabla_\lambda\otimes Q(e_v)\}_{v\succeq  p,\lambda\in X( \bt)_+})={\sf Q}_{\succeq p}.$$
So, \begin{equation}\label{eq:Y in Q and P}Y_p\in {\sf Q}_{\succeq p}\cap {\sf P}_{\preceq p} .
\end{equation}
\begin{lemma}\label{lem:samehull}Let $p\in W$.
\begin{itemize}
    \item $\hull(\{\nabla_\lambda\otimes X_v\}_{v\preceq p})=\hull(\{\nabla_\lambda\otimes P(-e_v)^*\}_{v\preceq  p,\lambda\in X( \bt)_+}).$


    
    \item $\hull(\{\nabla_\lambda\otimes Y_v\}_{v\succeq  p,\lambda\in X( \bt)_+})=\hull(\{\nabla_\lambda\otimes Q(e_v)\}_{v\succeq  p,\lambda\in X( \bt)_+}).$
\end{itemize}
In particular, $\hull(\{\nabla_\lambda\otimes X_v\}_{v\in W}) =\D$.
\end{lemma}


\begin{proof}
As $\cone(P(-e_p)^*\to X_p)$ lies in $\hull(\{\nabla_\lambda\otimes P(-e_v)^*\}_{v\prec  p,\lambda\in X( \bt)_+})$, we may argue by induction on the size of $\{\;v\in W\mid v\preceq p\;\}$. 

Similarly,
$\cone(Y_v\to Q(e_v))$ lies in $\hull(\{\nabla_\lambda\otimes Q(e_v)\}_{v\succ  p,\lambda\in X( \bt)_+})$ and we may use induction on the size of 
$\{\;v\in W\mid v\succeq p\;\}$.
\end{proof}



%--------------------------------------------------------------
\subsection{Computing morphisms between $X_p$ and $Y_p$}\label{subsec:X_p-Y_p-morphisms}
%--------------------------------------------------------------


The objects $X_v,Y_v$ are the key ingredient for constructing semiorthogonal sequences in $\Dd^b (\rep(\fb))$. Combined with the results of Section \ref{sec:X_p&Y_p-isom} below, the objects $X_v$'s will give the sought-for exceptional collection in $\D=\Dd^b (\rep(\fb))$. The goal of this section is to compute $\Hom_\D^{\bullet}(X_v,Y_v)$. This computation is achieved by the results of Section \ref{sec:b-cohom_vanish}.

\begin{lemma}\label{lem:X_p-Y_p-isom}
    \label{lem:XYmix}Let $\D=\Dd^b (\rep ({\fb}))$.
    \begin{enumerate}
        \item $\Hom_\D( X_v,Y_v[i])=\begin{cases}\sk &\text{if $i=0$},\\
0&\text{else.}\end{cases}$

\

\item If $w\succ  v$ then $\Hom_\D(Y_v, X_w[i])=0$ for all $i$.
    \end{enumerate}
\end{lemma}

\begin{proof}
$(1).$ Setting $M=N=\sk$ in Corollary \ref{cor:left orthogonal} or Corollary \ref{cor:Z left orthogonal}, we see with Proposition \ref{prop:cohinduced}
that
$\RHom_{\fb}(P(-e_v)^*, Q(e_v))={\rm H}^{\ast}( \bg,\sk)=\sk$. 
Applying $\Hom_\D(-, Q(e_v))$ to the triangle (\ref{eq:def_triang_X_p}), we obtain 
\begin{eqnarray}
& \dots \rightarrow \Hom _{\D}(\cone(P(-e_v)^*\to X_v),Q(e_v))\rightarrow \Hom _{\D}(X_p,Q(e_v))\rightarrow \\ \nonumber 
& \Hom _{\D}(P(-e_v)^*,Q(e_v))\rightarrow \dots 
\end{eqnarray}
Now $\cone(P(-e_v)^*\to X_v)\in {\sf P}_{\prec v}$ and $Q(e_v)\in {\sf P}_{\prec v}^\perp$. 
Thus, $$\Hom_\D(\cone(P(-e_v)^*\to X_v), Q(e_v)[i])=0$$ for all $i\in \mathbb Z$. 
We conclude that $\Hom _{\D}(X_v,Q(e_v))[i]=\Hom _{\D}(P(-e_v)^*,Q(e_v))[i]$ for all $i\in \mathbb Z$, and by Theorem \ref{th:indPQ}
the latter group is trivial if $i\neq 0$ and is isomorphic to $\sf k$ for $i=0$. Thus, 
\begin{equation}\label{eq:X_v-to-Q(e_v)}
\Hom _{\D}(X_v,Q(e_v))[i]=\begin{cases}\sk &\text{if $i=0$}\\
0&\text{else.}\end{cases}
\end{equation}

Now applying $\Hom_\D(X_v,-)$ to the triangle (\ref{eq:def_triang_Y_p}), we obtain 
\begin{eqnarray}
& \dots \rightarrow \Hom _{\D}(X_v,Y_v)\rightarrow \Hom _{\D}(X_v,Q(e_v))\rightarrow \\ \nonumber 
& \Hom _{\D}(X_v,{\sf cone}(Y_v\rightarrow Q(e_v))\rightarrow \dots 
\end{eqnarray}
and ${\sf cone}(Y_v\rightarrow Q(e_v)\in{\sf P}_{\preceq v}^{\perp}$, 
while $X_v\in {\sf Q}_{\succeq v}\cap {\sf P}_{\preceq v}\subset {\sf P}_{\preceq v}$. 
We conclude that $\Hom _{\D}(X_v,{\sf cone}(Y_v\rightarrow Q(e_v)[i])=0$ for all $i\in \mathbb Z$. 
From the previous exact triangle we obtain isomorphisms $\Hom _{\D}(X_v,Y_v[i])=\Hom _{\D}(X_v,Q(e_v)[i])$ for all $i\in \mathbb Z$ and from (\ref{eq:X_v-to-Q(e_v)}) the latter group is isomorphic to $\sk$ for $i=0$ and is trivial otherwise. 



$(2).$ Is an easier consequence of Corollary \ref{cor:left orthogonal}  or Corollary \ref{cor:Z left orthogonal}.
\end{proof}

As is clear from Lemma \ref{lem:X_p-Y_p-isom}, it is desirable to know  that $X_v$ is isomorphic to $Y_v$ for all $v\in W$.

%--------------------------------------------------------------
\section{Isomorphism of $X_p$ with $Y_p$}\label{sec:X_p&Y_p-isom}
%--------------------------------------------------------------
In this section we show that $X_p$
is isomorphic with $Y_p$, using a refinement of Theorem \ref{th:generation}. This is where the Steinberg weights tie everything together.

\

Let $\lambda \in X( \bt)$. Choose a dominant weight $\lambda ^+\in X({\bt})_+\cap W\lambda$ where $W\lambda$ is the $W$-orbit through $\lambda$ and $w\in W$ such that $\lambda ^+=w\lambda$ with $w$ being minimal with that property. 

Then put $$\pre_\lambda=\{v\mid (e_v,e_v)<(\lambda,\lambda)\}\cup\{ v\mid e_v=v^{-1}\lambda ^+\neq \lambda,~v>w\}.$$
That is:
$$\pre_\lambda=\{v\in W\mid e_v<_a\lambda\},$$
where $<_a$ is the antipodal excellent order \ref{def:antipodal}.


We have the following technical variation on Theorem \ref{th:generation}.

\begin{theorem}%[Generation of triangulated hull]
\label{th:cone generation}
  Let $\lambda,\mu\in X(\bt)$, $\mu\neq\lambda$.
  Assume that $\mu$
  is a weight of $P(-\lambda)^*$ or
  that $(\mu,\mu)<(\lambda,\lambda)$.
  
  Then
  $e^\mu$ is an object of the triangulated hull of
$\{\nabla_\nu\otimes M_v\mid \nu\in X( \bt)_+,~v\in\pre_\lambda\}.$
  \end{theorem}
  
\begin{proof}
We adapt the proof of Theorem \ref{th:generation}. First let us rephrase what we need to show. We need to show that if $\mu<_a\lambda$, then $e^\mu$ is in the triangulated hull of 

$\{\nabla_\nu\otimes M_v\mid \nu\in X( \bt)_+,~v\in W, ~e_v<_a\lambda\}.$
By  Proposition \ref{prop:G linear hull} this hull is $\bg$-linear.
We want to show by induction along $<_a$ that $e^\mu$ is in the  hull.
 So assume $e^\sigma$ is in the hull for $\sigma<_a\mu$.

There are two cases: $\mu$ is a Steinberg weight or it is not. If it is a Steinberg weight, then it is an $e_v$ with $e_v<_a\lambda$. And all weights $\sigma$ of $M_v$ that are distinct from $\mu=e_v$ satisfy $\sigma<_a\mu$.
If $\mu$ is not a Steinberg weight, then by Corollary \ref{cor:if not Steinberg then tensor} there is a module $N$  in  $\{\nabla_\nu\otimes e^\tau\mid \nu\in X( \bt)_+,~(\tau,\tau)<(\mu,\mu)\}$ 
with $\mu$ a weight  of $N$ of multiplicity one, and all weights  $\sigma$ of $N$ satisfying 
$\sigma\leq_a\mu$. But $(\tau,\tau)<(\mu,\mu)$ implies $\tau<_a\mu$.
\end{proof}
For completeness we mention
\begin{corollary}\label{cor:ordered_generate}
    Let $\lambda\in X(\bt)$. Then $e^\lambda$ lies in the $R(\bg)$-submodule of $R(\fb)$
generated by the $[M_v]$ with $e_v\leq_a\lambda$.
\end{corollary}
\begin{proof}
    `Same proof' by induction along $\leq_a$.
\end{proof}
\begin{remark}\label{rem:Steinberg from anti}
    Thus $\lambda$ is a Steinberg weight if and only if $e^\lambda$ does not lie in the $R(\bg)$-submodule of $R(\fb)$
generated by the $e^\mu$ with $\mu<_a\lambda$.
 If $\lambda$ is a Steinberg weight, then $\lambda=e_w$ where $w\in W$ is the element of minimal length making $w\lambda$ dominant.
\end{remark}

\

Let $f_p$ be the natural map from $P(-e_p)^*$ to $Q(e_p)$. Notice that $f_p$ factors as a surjection $P(-e_p)^*\to\sk_{e_p}$, followed by an injection $\sk_{e_p}\to Q(e_p)$.
We have an exact triangle 
$$\dots \rightarrow \kernel(f_p)[1]\to \cone(f_p)\to \coker(f_p)\rightarrow \dots$$ 
Take $M_v=Q(e_v)$ for $v\succ p$ and $M_v=P(-e_v)^*$ for $v\prec p$. We do not need to specify $M_p$, because $p\notin \pre_{e_p}$.



\begin{corollary}\label{cor:cone-f_p}
Let $p\in W$. Then $\cone(f_p)$ belongs to the hull of the union of the
following three sets:
\begin{eqnarray*}
\{\;\nabla_\lambda\otimes Q(e_v)&\mid& \lambda \mbox{ is dominant,  } v\succ  p\mbox{ and }(e_v,e_v)<(e_p,e_p)\;\}\cup\\
\{\;\nabla_\lambda\otimes Q(e_v)&\mid& \lambda \mbox{ is dominant, }  v\succ p\mbox{ and }v^{-1} pe_p=e_v\neq e_p\;\}\cup\\
\{\;\nabla_\lambda\otimes P(-e_v)^*&\mid& \lambda \mbox{ is dominant, } v\prec p \mbox{ and }(e_v,e_v)<(e_p,e_p)\;\}.
\end{eqnarray*}
In particular, $\cone(f_p)$ belongs to the hull of $ {\sf P}_{\prec p}\cup {\sf Q}_{\succ p}$.
\end{corollary}

\begin{proof}
Apply Theorem \ref{th:cone generation} with $\lambda=e_p$ to 
$\kernel(f_p)$ and $\coker(f_p)$.
\end{proof}
\begin{remark}
    Observe that the proof involves the partial order $<_a$, but the conclusion, that $\cone(f_p)$ belongs to the hull of $ {\sf P}_{\prec p}\cup {\sf Q}_{\succ p}$, refers to $\prec$. This conclusion just
    uses that only $M_v$ with $v\neq p$ are used. It is at the Steinberg weights that the Bruhat order meets the antipodal excellent order. Both orders are important. Compare Proposition \ref{prop:p-q-SOD-2} and Corollary \ref{cor:ordered_generate}.
\end{remark}
We are now in a position to prove the isomorphism $X_p=Y_p$. 
Note that ${i^p}_{\preceq p}^{\ !}Q(e_p)=Y_p$ and ${i^p}_{\preceq p}^{\ !}(X_p)=X_p$ as ${i^p}_{\preceq p}^{\ !}={\sf id}_{{\sf P}_{\preceq p}}$ on ${\sf P}_{\preceq p}$. 
Moreover, the functor ${i^p}_{\preceq p}^{\ !}$ annihilates ${\sf Q}_{\succ p}$
as ${\sf Q}_{\succ p}={\sf P}_{\preceq p}^{\perp}$. That is:
\begin{eqnarray*}
{i^p}_{\preceq p}^{\ !}({\sf Q}_{\succ p})=0,
\end{eqnarray*}
and 
\begin{eqnarray*}
{i^p}_{\preceq p}^{\ !}({\sf P}_{\prec p})\subset {\sf P}_{\prec p}.
\end{eqnarray*}
As $\cone(f_p)$ belongs to the hull of $ {\sf P}_{\prec p}\cup {\sf Q}_{\succ p}$, we conclude that \begin{eqnarray}{i^p}_{\preceq p}^{\ !}(\cone(f_p))\in {\sf P}_{\prec p}.\end{eqnarray}
Consider the exact triangle 
\begin{equation}
\dots \rightarrow P(-e_v)^*\rightarrow Q(e_v)\rightarrow \cone(f_p)\rightarrow \dots 
\end{equation}
Applying to the triangle the functor ${i^p}_{\preceq p}^{\ !}$ and remembering 
that ${i^p}_{\preceq p}^{\ !}(Q(e_v))=Y_p$, we obtain 
\begin{equation}
\dots \rightarrow P(-e_v)^*\rightarrow Y_p\rightarrow {i^p}_{\preceq p}^{\ !}(\cone(f_p))\rightarrow \dots 
\end{equation}
with ${i^p}_{\preceq p}^{\ !}(\cone(f_p))\in {\sf P}_{\prec p}$.
Thus
\begin{equation}\label{eq:cone in P} \cone(P(-e_v)^*\rightarrow Y_p)\in  {\sf P}_{\prec p}.
\end{equation} 
Now apply to the triangle 
\begin{equation}
\dots \rightarrow P(-e_v)^*\rightarrow Y_p\rightarrow \cone(P(-e_v)^*\rightarrow Y_p)\rightarrow \dots 
\end{equation}
the projection functor ${i^q}_{\succeq p}^{\ \ast}$ onto ${\sf Q}_{\succeq p}$.
We have ${i^q}_{\succeq p}^{\ \ast}(P(-e_v)^*)=X_p$ and ${i^q}_{\succeq p}^{\ \ast}(Y_p)=Y_p$, as $Y_p\in{\sf Q}_{\succeq p}\cap {\sf P}_{\preceq p}\subset {\sf Q}_{\succeq p}$ by (\ref{eq:Y in Q and P}). Further ${i^q}_{\succeq p}^{\ \ast}(\cone(P(-e_v)^*\rightarrow Y_p))=0$ because of (\ref{eq:cone in P}), as ${i^q}_{\succeq p}^{\ \ast}$ annihilates ${\sf P}_{\prec p}$. Thus, $X_p=Y_p$ and by Lemma \ref{lem:X_p-Y_p-isom}, (1), the object $X_p$ is exceptional in $\D=
 \Dd^b (\rep ({\fb}))$.

\begin{remark}
    Actually we need a little more than exceptionality. See part 1 of Theorem \ref{th:X_p-G-exceptional-coll}.
\end{remark}
%--------------------------------------------------------------
\section{Semi-orthogonal decomposition of $\Dd ^b(\rep(\fb))$ as a 
$\bg$-linear category}\label{sec:SODofB-modoverG-mod}
%--------------------------------------------------------------

 Let $\sk$ be a field.
\begin{theorem}\label{th:X_p-G-exceptional-coll}
Let $\D=\Dd^b (\rep ({\fb}))$.
Let $M$, $N$ be finite dimensional $ \bg $-modules. Let $v,w\in W$.
\begin{enumerate}
\item $\Hom_\D(M\otimes X_v,N\otimes X_v[i])=\Ext^i_ \bg (M,N) $ for all $i$.
\item If $w\succ  v$ then $\Hom_\D(M\otimes X_v,N\otimes X_w[i])=0$ for all $i$.
\end{enumerate}
\end{theorem}
\begin{proof}
$(1).$ Recall that for $p\in W$, the object $X_p$ is both the projection of $P(-e_v)^*$ onto the category ${\sf Q}_{\succeq p}$ and the projection of $Q(e_p)$ onto ${\sf P}_{\preceq p}$. 
Tensoring the triangles (\ref{eq:def_triang_X_p}) and (\ref{eq:def_triang_Y_p}) with modules $M$ and $N$, respectively, we get:
\begin{equation}\label{eq:otimesM}
\dots \rightarrow M\otimes P(-e_p)^*\rightarrow M\otimes X_p\rightarrow M\otimes {\sf cone}(P(-e_p)^*\rightarrow X_p)\rightarrow \dots 
\end{equation}
and 
\begin{equation}\label{eq:otimesN}
\dots \rightarrow N\otimes X_p\rightarrow N\otimes Q(e_p)\rightarrow N\otimes {\sf cone}(X_p\rightarrow Q(e_p))\rightarrow \dots 
\end{equation}
We have that ${\sf cone}(P(-e_p)^*\rightarrow X_p)\in {\sf P}_{\prec p}$ and ${\sf cone}(X_p\rightarrow Q(e_p))\in {\sf Q}_{\succ p}$, and both subcategories ${\sf P}_{\prec p}$ and ${\sf Q}_{\succ p}$ are $ \bg$-linear, thus also
$M\otimes {\sf cone}(P(-e_p)^*\rightarrow X_p)\in {\sf P}_{\prec p}$ and $N\otimes {\sf cone}(X_p\rightarrow Q(e_p))\in {\sf Q}_{\succ p}$. 

We have $\Hom _{\D}(M\otimes {\sf cone}(P(-e_p)^*\rightarrow X_p),N\otimes X_p[i]))=0$ for all $i$ since $X_p\in {\sf Q}_{\succeq p}\cap {\sf P}_{\preceq p}\subset {\sf Q}_{\succeq p}$ and ${\sf Q}_{\succeq p}={\sf P}_{\prec p}^{\perp}$ by 
 (\ref{eq:p-q-SOD-2}). Applying $\Hom _{\D}(-,N\otimes X_p)$ to the triangle (\ref{eq:otimesM}), we then obtain $\Hom _{\D}(M\otimes X_p,N\otimes X_p[i])=\Hom _{\D}(M\otimes P(-e_p)^*,N\otimes X_p[i])$ for all $i$.
 
Next, we have $\Hom _{\D}(M\otimes X_p,N\otimes {\sf cone}(X_p\rightarrow Q(e_p)[i]))=0$ for all $i$ since $X_p\in {\sf Q}_{\succeq p}\cap {\sf P}_{\preceq p}\subset {\sf P}_{\preceq p}$ and ${\sf Q}_{\succ p}={\sf P}_{\preceq p}^{\perp}$ by 
 (\ref{eq:p-q-SOD-1}). Applying $\Hom _{\D}(M\otimes X_p,-)$ to the triangle (\ref{eq:otimesN}), we then obtain $\Hom _{\D}(M\otimes P(-e_p)^*,N\otimes X_p[i])=\Hom _{\D}(M\otimes P(-e_p)^*,N\otimes Q(e_p)[i])$ for all $i$. The latter group is isomorphic to $\Ext ^i_{\fb}(M\otimes P(-e_p)^*,N\otimes Q(e_p))$. Now the statement in (1) follows by Corollary \ref{cor:left orthogonal}, (1).



$(2).$ The second part follows similarly using Corollary \ref{cor:left orthogonal}, (2).
\end{proof}
\begin{remark}\label{rem:flat representative}
If $\sk=\Z$, one may use the resolution property Proposition \ref{prop:resolution} to choose a  representative of the isomorphism class of $X_v$ so
    that the representative
    is a bounded complex of modules which are flat over $\Z$. With such choice 
     the proof of Theorem \ref{th:X_p-G-exceptional-coll} goes through for $M,N\in\rep_\fr(\bg)$ and by Proposition \ref{prop:resolution} 
    the result then follows for $M,N\in\rep(\bg)$.
\end{remark}

In fact we have
\begin{theorem}Let $\sk$ be a field or $\Z$.
\label{th:Z-X_p-G-exceptional-coll}
Let $\D=\Dd^b (\rep ({\fb}))$.
Let $M, N\in \Dd^b(\rep(\bg))$ and $v,w\in W$.
\begin{enumerate}
\item $\RHom_\D(M\otimes X_v,N\otimes X_v)=\RHom_ {\Dd^b(\rep(\bg))} (M,N) $,
\item If $w\succ  v$ then $\RHom_\D(M\otimes X_v,N\otimes X_w)=0$.
\end{enumerate}
\end{theorem}
\begin{proof}

Part (1).   It suffices to treat the case $M=\sk$
because 
$\RHom_\D(L\otimes X_v,N\otimes X_v)=\RHom_\D( X_v,L^*\otimes N\otimes X_v)$ and
$\RHom_  {\Dd^b(\rep(\bg))}  (L,N) =\RHom_  {\Dd^b(\rep(\bg))} (\sk,L^*\otimes N)$.  By Theorem \ref{th:X_p-G-exceptional-coll} or Remark
\ref{rem:flat representative} we  know the result when $N\in \rep(\bg)$ and $\RHom_ \bg (\sk,N) $ is concentrated in at most one degree. 
For instance, this is the case when $N=\nabla_\nu$ for some dominant $\nu$.
So it holds for  $N$ in the hull $\Dd^b(\rep(\bg))$ of the $\nabla_\nu,\nu\in X(\bt)_+$. 

Part (2) is easier.
\end{proof}

\begin{corollary}
\label{cor:Rind_B^G(X_p*otimesX_p)}

\

\begin{enumerate}
\item $\Rind_{\fb}^{ \bg}(X_v^{\ast}\otimes X_v)=\sk$ .
\item If $w\succ  v$ then $\Rind _{\fb}^{ \bg}(X_v^{\ast}\otimes X_w)=0$ .
\end{enumerate}
\end{corollary}
\begin{proof}Part (1)
\begin{eqnarray*}&\RHom _{ \Dd^b(\rep(\bg))}( M, \sk)=\RHom_\D(M\otimes X_v, X_v)=\\&\RHom_\D(M, X_v^\ast\otimes X_v)=
\RHom_{\Dd^b(\rep(\bg))}(M,\Rind_\fb^\bg( X_v^\ast\otimes X_v)),
\end{eqnarray*}
for $M\in \Dd^b(\rep(\bg))$. By the Yoneda Lemma it follows that $\Rind_{\fb}^{ \bg}(X_v^{\ast}\otimes X_v)=\sk$.

Part (2) is easier.
\end{proof}
Let $\sk$ be a field or $\Z$.
\begin{proposition}\label{prop:sfX_p-admissibility}
    Let $p\in W$.
    The 
strictly full subcategory of
$\D=\Dd ^b(\rep(\fb))$ generated by $\{M\otimes X_p \mid M\in \Dd ^b(\rep( \bg))\}$
is an admissible subcategory.
\end{proposition}

\begin{proof}
Recall that by Proposition \ref{prop:resolution} we may extend the exact bifunctors
 $$-\;\otimes_\sk-:\rep_\fr( \fb)\times \rep_\fr( \fb)\to \rep_\fr( \fb)$$ and $$\Hom_\sk(-,-):\rep_\fr( \fb)\times \rep_\fr( \fb)\to \rep_\fr( \fb)$$ 
 to the bounded derived category of $\rep( \fb)$. We put $M^*=\Hom_\sk(M,\sk)$.
Then the right adjoint of the inclusion of the subcategory is $$\Rind_ \fb^ \bg(-\otimes_\sk X_p^*)\otimes_\sk X_p$$ 
and the left adjoint is $$(\Rind_ \fb^ \bg((-)^*\otimes_\sk X_p))^*\otimes_\sk X_p.$$
\end{proof}


Let $\sk$ be a field or $\Z$. In the latter case $\bg$ means $\Gm$.
%-----------------------------------------------------------------
\begin{theorem}\label{th:semi-orthogonal}
%-----------------------------------------------------------------
For $v\in W$, denote $\sf X _v$ to be the 
strictly full subcategory of
$\D=\Dd ^b(\rep(\fb))$ generated by $\{M\otimes X_v \mid M\in \Dd ^b(\rep( \bg))\}$.
Then the category $\D$ has a $ \bg$-linear semiorthogonal decomposition 
\begin{equation}
\D= \langle \sf X _v\rangle _{v\in W}
\end{equation}
with respect to the order $\prec$ on the Weyl group $W$. Each subcategory $\sf X _v$ is equivalent to $\Dd ^b(\rep( \bg))$.
\end{theorem}
\begin{proof}
Denote $\Phi _p:\Dd^b (\rep( \bg ))\rightarrow \D$ the $ \bg$-linear functor $M\rightarrow M\otimes _{\sk} X_p$. 
Now Corollary \ref {cor:Rind_B^G(X_p*otimesX_p)}
 and Proposition 
\ref{prop:G-linear_exceptional} give that $\Phi _p$ is fully faithful. 




Let now $w\succ v$. 
By  Corollary \ref {cor:Rind_B^G(X_p*otimesX_p)} and \cite[I Proposition 3.6]{Jan} we have
 $\Rind _{\fb}^{ \bg}(X_v^{\ast}\otimes X_w\otimes M)=\Rind _{\fb}^{ \bg}(X_v^{\ast}\otimes X_w)\otimes M=0$ for $M\in \Dd ^b(\rep( \bg))$.
By Proposition \ref{prop:sfX_p-admissibility}, each subcategory ${\sf X}_v, v\in W$ is admissible in $\Dd ^b(\rep( \fb))$. 
Proposition \ref{prop:G-linear_sod} then gives that the sequence $\langle \sf X _v\rangle _{v\in W}$ of admissible subcategories of $\D$ is $ \bg$-semiorthogonal with respect to the order $\prec$ on $W$. 
Lemma \ref{lem:samehull} states that the triangulated hull of $\langle \sf X _v\rangle _{v\in W}$ coincides with $\D$. 
Thus, it is a $ \bg$-linear semi-orthogonal decomposition of $\D$.
\end{proof}




%--------------------------------------------------------------
\section{Full exceptional collections in $\Dd ^b( \bg/ \fb)$}\label{sec:FEConGmodB}
%--------------------------------------------------------------
We recall first the sheafification functor $\Ll$. 
Let $ \bg$ be a flat $\sf k$-group scheme acting freely (from the right) on a flat $\sf k$-scheme $X$ such that $X/ \bg$ is a scheme. 
Associated to each $ \bg$-module $M$ is a sheaf $\Ll (M)=\Ll _{X/ \bg}(M)$ on $X/ \bg$, \cite[I, Section 5.8]{Jan}. 
The functor $\Ll : {\Rep}({ \bg})\rightarrow {\rm Sh}(X/ \bg)$ is exact and lands in the subcategory ${\QCoh}(X/ \bg)$ of quasi--coherent 
$\Oo _{X/ \bg}$-modules of ${\rm Sh}(X/ \bg)$, \cite[I, Proposition 5.9]{Jan}. 
Since $\Ll$ is exact, denote also $\Ll:\Dd ({\Rep}({ \bg}))\rightarrow \Dd _{{\QCoh}(X/ \bg)}{\rm Sh}(X/ \bg)\footnote{The derived category of complexes on ${\rm Sh}(X/ \bg)$ with quasi-coherent cohomology.}$ its derived functor.

Now set $ \bg:= \fb, X= \bg$ (so $X= \bg/ \fb$) in the definition of the functor $\Ll$. It restricts to the functor between the bounded derived categories $\Dd ^b({\rep}({\fb}))\rightarrow \Dd ^b( \bg/ \fb)$. 
We put $\tX_v:=\Ll(X_v)$ for $v\in W$. 

\


Let $\sk$ be a field or $\Z$. In the latter case $\bg$ means $\Gm$. Recall that $\Dd^b ( \bg/\fb):=\Dd^b (\Coh( \bg/\fb))$.
Theorem \ref{th:semi-orthogonal} implies the following:

\begin{theorem}\label{th:coh-semi-orthogonal}Let $\D=\Dd^b (\bg/\fb)$.
Let $v,w\in W$.\\
\begin{enumerate}
\item $\Hom_\D(\tX_v,\tX_v[i])=\begin{cases}\sk &\text{if $i=0$},\\
0&\text{else.}\end{cases}$\\

\

\item If $w\succ  v$ then $\Hom_\D(\tX_v,\tX_w[i])=0$ for all $i$.\\
\item The triangulated hull of $\{\tX_v\mid v \in W\}$ is $\D$.
\end{enumerate}
In other words, the collection of objects $(\tX_v)_{v \in W}$ is a full exceptional collection in $ \D$.
\end{theorem}


\begin{proof}
$(1)$. By Proposition \ref{prop:derived_B-cohomology}, the functor 
$\RHom _{ \bg/\fb}(\Ll (-),\Ll (-))$ canonically factors through $\Rind _{\fb}^{ \bg}((-)^{\ast}\otimes (-))$
and $\RHom_\D(\tX_v,\tX_v)=\RHom _{ \bg/\fb}(\Ll (X_v),\Ll (X_v))={\sf For}\circ \Rind_{\fb}^{ \bg}(X_v^{\ast}\otimes X_v)$ where $\sf For$ is the forgetful functor $\Dd ^b(\rep( \bg))\rightarrow \Dd ^b( {\kmod})$. This also holds when $\sk=\Z$.
Now by Corollary \ref{cor:Rind_B^G(X_p*otimesX_p)} (1), $\Rind_{\fb}^{ \bg}(X_v^{\ast}\otimes X_v)=\sk$. 



$(2)$. This is the same argument as in (1), using Corollary \ref{cor:Rind_B^G(X_p*otimesX_p)} (2).





$(3).$  We 
will prove that 
the 
triangulated hull $\langle \tX_v, v\in W\rangle$ of $\{\tX_v, v\in W\}$ in $\D=\Dd^b (\Coh( \bg/\fb))$
is the whole $\Dd^b (\Coh( \bg/\fb))$. Observe that by the previous items (1) and (2), the collection of $\tX_v, v\in W$ is exceptional. 
Thus,  $\langle \tX_v, v\in W\rangle\subsetq \Dd^b (\Coh( \bg/\fb))$ is an admissible (hence thick) subcategory of $\Dd^b (\Coh( \bg/\fb))$, see Section \ref{sec:rappel_triangulated}. We are bound to show that $\langle\tX_v, v\in W\rangle^{\perp}=0$. Let us show that $\langle \tX_v, v\in W\rangle$ contains the category $\Dd ^b_{ \bg}( \bg/ \fb):=\Dd ^b({\sf Coh}^{ \bg}( \bg/ \fb))$
of all $ \bg$-equivariant objects of $\Dd^b (\Coh( \bg/\fb))$. 

By Lemma \ref{lem:samehull},  $\hull(\{\nabla_\lambda\otimes X_v\}_{v\in W})$ coincides with the hull of $\rep( \fb)$ in $\Dd^b (\rep ({\fb}))$. 
In particular, the abelian category $\rep( \fb)\subset \Dd^b (\rep ({\fb}))$ is contained in  $\hull(\{\nabla_{\lambda}\otimes X_v\}_{v\in W})$. 
By definition, $\tX_v=\Ll (X_v)$; applying the functor $\Ll$ to the above inclusion, we see that the category of $ \bg$-equivariant coherent sheaves ${\sf Coh}^{ \bg}( \bg/ \fb)=\Ll (\rep( \fb))$ is contained in the image under 
$\Ll$ of  $\hull(\{\nabla_{\lambda}\otimes X_v\}_{v\in W})$. 
Recall that the functor $\Ll$ is monoidal; thus, $\Ll (\nabla_{\lambda}\otimes X_v)=\Ll (\nabla_{\lambda})\otimes\tX_v$. 
Now $\Ll (\nabla_{\lambda})$ is a ($ \bg$-equivariant) trivial vector bundle on $ \bg/ \fb$, so the image of 
$\hull(\{\nabla_{\lambda}\otimes X_v\}_{v\in W})$ under $\Ll$ coincides with the hull $\langle \tX_v, v\in W\rangle$ of $\{ \tX_v, v\in W\}$. 
We conclude that $\langle \tX_v, v\in W\rangle$ contains the category $\Dd ^b_{ \bg}( \bg/ \fb)$.

Now since $\Dd ^b_{ \bg}( \bg/ \fb)\subsetq \langle \tX_v, v\in W\rangle$, the  category $\langle \tX_v, v\in W\rangle$
contains all line bundles $\Ll _{\lambda},\lambda \in X(\bt)_+$, the group $ \bg$ being simply connected. Now
$\langle \tX_v, v\in W\rangle^\perp=0$ by Corollary \ref{cor:line_bundles-generation}.
 By Lemma \ref{lem:notorious} we get $\Dd^b (\rep ({\fb}))=\langle \tX_v, v\in W\rangle$.
\end{proof}

\begin{remark}\label{rem:basechangeG-to-point}
{\rm Theorem \ref{th:coh-semi-orthogonal} is in fact a simple instance of base change for semiorthogonal decompositions, \cite{Kuzbasech}, which in this particular case can be stated in elementary terms. In more detail, consider a flat morphism $\pi\colon X\rightarrow S$ between smooth projective varieties, a closed point $s\in S$ and the base change diagram along the embedding $i_s\colon s\hookrightarrow S$:

\begin{figure}[H]
$$\xymatrix @C5pc @R4pc {
X_s\ar[r]^i\ar@<-0.1ex>[d]_{\pi _s} & X \ar@<-0.4ex>[d]^{\pi} \\
       s\ar[r]^{i_s} & S}
$$
%\caption{}	
\end{figure}

Assume given an $S$-linear semiortogonal decomposition $\langle \D _1,\dots ,\D _n\rangle$ of $\Dd ^b(X)$, such that each admissible subcategory $\D _i$ is equivalent to $\Dd ^b(S)$. Then each embedding functor $\Phi _k:\D _k\rightarrow \Dd ^b(X), k=1,\dots,n$ is given by $\Phi _k(-)=(-)\otimes E_k$ where $E_k,k=1,\dots n$ is a collection of objects of $\Dd ^b(X)$ with the following two properties: 1) $R\pi _{\ast}R{\mathcal Hom}(E_i,E_j)=0$ for $i>j$ and 2) $R\pi _{\ast}{\mathcal Hom}(E_k,E_k)=\Oo _S$ for all $k$ (cf. Propositions \ref{prop:G-linear_sod} and \ref{prop:G-linear_exceptional}). Restricting the objects $E _k\in \Dd ^b(X)$ to the fiber $X_s$, one obtains a collection of objects $\Ee _k:=i^{\ast}E_k, k=1,\dots n$. Now the claim is that the collection $\Ee _k, k=1,\dots n$ is exceptional in $\Dd ^b(X_s)$: for instance, to see that $\RHom _{X_s}(\Ee _i,\Ee _j)=0$ for $i>j$, by \cite[Lemma 2.32]{KuzHPD} base change holds for the above Cartesian square, thus \
\begin{eqnarray}
&\RHom _{X_s}(\Ee _i,\Ee _j)=R{\pi _s}_{\ast}(\Ee _i^{\ast}\otimes ^{\mathbb L}\Ee _j)=R{\pi _s}_{\ast}i^{\ast}(E_i^{\ast}\otimes ^{\mathbb L}E _j)= \\
& i^{\ast}_sR\pi _{\ast}(E_i^{\ast}\otimes ^{\mathbb L}E _j)=i^{\ast}_sR\pi _{\ast}R{\mathcal Hom}(E_i,E_j)=0
\end{eqnarray}

by 1) above. Similarly for $\RHom _{X_s}(\Ee _i,\Ee _i)={\sf k}$ for $i=1,\dots n$. Finally, the objects $\Ee _i,i=1,\dots n$ generate $\Dd ^b(X_s)$: for if there was a non-trivial object $\Ff \in \langle \Ee _1,\dots ,\Ee _n\rangle  ^{\perp}$, its pushforward ${i_s}_{\ast}\Ff \in \Dd ^b(X)$ would be a non-trivial object in the right orthogonal to the semiortogonal decomposition $\langle \D _1,\dots ,\D _n\rangle$ of $\Dd ^b(X)$, a contradiction.
}
\end{remark}

%--------------------------------------------------------------
\subsection{Variations} 
%--------------------------------------------------------------

The $\tX_p$ may depend on the choice of the total order $\prec$
on $W$. Therefore it is not clear that $\tX_p$ is perpendicular to
$\tX_q$ when $\ell(p)=\ell(q)$.
In fact computer assisted computations in $\sK_ {\bt}( G /{\fb})$ indicate 
that this fails already for type ${\bf B}_3$ with $\ell(p)=\ell(q)=3$ and for type ${\Cc}_3$ with $\ell(p)=\ell(q)=2$.

%----------------------------------------------------------------------

\section{Generalised flag varieties}\label{sec:caseP}

Let $\bp$ be a parabolic subgroup containing $\fb$. (If $\sk=\Z$, then $\bp$ means $\mathbb P$.)
We seek a full exceptional collection on $\bg/\bp$.
Let $W_\bp$ be the Weyl group of $\bp$,
generated inside $W$ by the $s_\alpha$ with $\alpha$ simple and
$\bp_\alpha\subset \bp$.
Let $W^\bp$ be the set of minimal coset representatives of $W/W_\bp$, cf.\  Lemma \ref{lem:minimal-coset-representative}. 
According to Steinberg \cite{St} we may use as generators of  the $R(\bg)$-module $R(\bp)$, over $\cC$, 
the classes of the
irreducible $\bp$-modules with highest weight $e_v$ where $v$
runs over $W^\bp$, and $e_v$ is still defined as in \ref{subsec:Steinberg}. Inspired by that,
we restrict our total order $\prec$ from $W$ to $W^\bp$.
Note our convention that the notation $e_v$ keeps the meaning it had when $\bp=\fb$. But  only the $e_v$ with $v\in W^\bp$ will be relevant.
\begin{remark}The map $W\to W^\bp$ which sends $w\in W$ to the minimal representative of the coset $wW_\bp$ is a poset map.
If one has chosen a total order on $W^\bp$ refining the (restriction to $W^\bp$) of the Bruhat order, then this
chosen order can be extended to
a total order on $W$ that refines the Bruhat order on $W$. We already had to choose $\prec$ on $W$ in section \ref{subsec:totalorderonW}, so it makes sense to keep that order $\prec$ and restrict it to $W^\bp$.
\end{remark}

For $v\in W^\bp$ we will find ${\hat X}_v\in \Dd^b (\rep ( \bp ))$  such that
 
\begin{theorem}Let $\sk$ be a field or $\Z$.
\label{th:Z-X^P_p-G-exceptional-coll}
Let $\D=\Dd^b (\rep ( \bp ))$.
Let $M, N\in \Dd^b(\rep(\bg))$ and $v,w\in W^\bp$.
\begin{enumerate}
\item $\RHom_\D(M\otimes {\hat X}_v,N\otimes {\hat X}_v)=\RHom_ {\Dd^b(\rep(\bg))} (M,N) $,
\item If $w\succ  v$ then $\RHom_\D(M\otimes {\hat X}_v,N\otimes {\hat X}_w)=0$.
\end{enumerate}
\end{theorem}

\begin{theorem}\label{th:Psemi-orthogonal}
%-----------------------------------------------------------------
For $v\in W^\bp$, denote ${\hat{\sf X}} _v$ to be the 
strictly full subcategory of
$\D=\Dd ^b(\rep(\bp))$ generated by $\{M\otimes {\hat X}_v \mid M\in \Dd ^b(\rep( \bg))\}$.
Then the category $\D$ has a $ \bg$-linear semiorthogonal decomposition 
\begin{equation}
\D= \langle {\hat{\sf X}} _v\rangle _{v\in W^\bp}
\end{equation}
with respect to the order $\prec$ on  $W^\bp$. Each subcategory ${\hat{\sf X}} _v$ is equivalent to $\Dd ^b(\rep( \bg))$.
\end{theorem}

\begin{definition}
 We put $\hat\tX_v:=\Ll_{\bg/\bp}({\hat X}_v)$ for $v\in W^\bp$. 
\end{definition}


\begin{theorem}\label{th:Pcoh-semi-orthogonal}Let $\D=\Dd^b (\Coh( \bg / \bp))$.
Let $v,w\in W^\bp$.\\
\begin{enumerate}
\item $\Hom_\D(\hat\tX_v,\hat\tX_v[i])=\begin{cases}\sk &\text{if $i=0$},\\
0&\text{else.}\end{cases}$\\



\item If $w\succ  v$ then $\Hom_\D(\hat\tX_v,\hat\tX_w[i])=0$ for all $i$.\\
\item The triangulated hull of $\{\hat\tX_v\mid v \in W^\bp\}$ is $\D$.
\end{enumerate}
In other words, the collection of objects $(\hat\tX_v)_{v \in W^\bp}$ is a full exceptional collection in $ \D$.
\end{theorem}

\

To prove these Theorems we will need to find replacements--often with hatted notation--of the key ingredients used in 
the case $\bp=\fb$.
This will take the rest of the section. 



Recall that by \cite[II Corollary 4.7]{Jan} we may view $\rep(\bg)$ as a subcategory of $\rep(\bp)$
and $\Dd^b(\rep(\bg))$ as a subcategory of
$\Dd^b(\rep(\bp))$, so that we may suppress 
$\res^\bg_\bp$ in the notation. Similarly we may suppress
$\res^\bp_\fb$. And when we say that a certain $\fb$-module $M$ is a $\bp$-module, this will mean that $M=\res^\bp_\fb\ind_\fb^\bp(M)$.

\begin{remark}
The strategy for constructing $\hat\tX_v$ is the same as for $\tX_v$, but
the construction does not immediately imply the precise relation between $\tX_v$, $\hat\tX_v$ and $\bg/\fb\to \bg/\bp$.
More specifically, one expects that $\tX_v$ is always
the pull back of $\hat\tX_v$ for $v\in W^\bp$. This is indeed the case:
\end{remark}

\begin{theorem}\label{th:P and B}
Let $v\in  W^\bp$ and let $\pi:\bg/\fb\to \bg/\bp$ be the natural map.
Then ${\hat{\sf X}} _v$ equals ${{\sf X}} _v$ in $\Dd^b(\rep(\fb))$ and $\tX_v=\pi^*\hat\tX_v$.
\end{theorem}

Recall that 
$\Pi=\{\alpha _1,\dots ,\alpha _r\}$ is the set of simple roots and 
$(\omega_i,\alpha_j^\vee)=\delta_{ij}$. 
If $\alpha=\alpha_i$, we also write $\omega_i$ as $\omega_\alpha$.


\begin{notation}
Let $\Pi^\bp$ be the set of simple roots $\alpha$ for which $\bp_\alpha\subsetq \bp$.
Let $\Pi^{\notin \bp}$ be the set of simple roots outside $\Pi^\bp$.
We say that $\lambda$ is $\bp$-dominant if $(\alpha^\vee,\lambda)\geq0$ for all $\alpha\in \Pi^\bp$.
The set of $\bp$-dominant weights is denoted $X^\bp_+$.
Every $W_\bp$-orbit of weights intersects $X^\bp_+$ in a unique element.
Let $\bg_\bp$ be the semisimple subgroup of $\bp$ with $\Pi^\bp$ as simple roots. (So $\bt\bg_\bp$ is a Levi subgroup $\bl_\bp$ of $\bp$.)
Note that $\bg_\bp$ is simply connected.
\end{notation}


\begin{lemma}\label{lem:PisP}Let $v\in W^\bp$.
Then $P(-e_v)$ is a $\bp$-module and therefore $P(-e_v)^*$ is a $\bp$-module.
Conversely, if $e_w$ is  $\bp$-dominant, then $w\in W^\bp$.
\end{lemma}
\begin{proof}
We want to show that $P(-e_v)=\res^\bp_\fb\ind_\fb^\bp(P(-e_v))$.
Let $\bp=\bu_{\bp}\bl_{\bp}$ be a Levi decomposition of $\bp$ where $\bu_{\bp}$ is the unipotent radical of $\bp$ and $\bl_{\bp}$ is the Levi component. Let $w_0^{\bp}\in W_\bp$ be the longest element and consider a reduced decomposition $w_0^{\bp}=s_{1}\cdots s_{l}$ of $w_0^\bp$. The  functors ${\sf H}_{w_0^{\bp}}$ and $\res^\bp_\fb\ind_\fb^\bp$  are identical, because 
    $\overline{\fb w_0^{\bp}\fb}/\fb=\bp/\fb$. So we now want to prove that ${\sf H}_{w_0^{\bp}}(P(-e_v))=P(-e_v)$. 
By Lemma \ref{lem:Hwstar}
we have ${\sf H}_{w_0^{\bp}}=\sH_{s_{1}}\circ\cdots\circ\sH_{s_{l}}$.
So it suffices to show that $\sH_s(P(-e_v))=P(-e_v)$ for $s\in W_\bp$ simple.
 Now $v$ is a minimal coset representative in $W/W_\bp$, so $vs>v$ for $s\in W_\bp$.
Then $s v^{-1}>v^{-1}$, so $s v^{-1}w_0<v^{-1}w_0$, and $s\star v^{-1}w_0=v^{-1}w_0$.
Thus $\sH_s(P(-e_v))=
{\sH}_{s}{\sH}_{v^{-1}w_0}(-w_0ve_v)={\sH}_{v^{-1}w_0}(-w_0ve_v)=P(-e_v)$ (Cf. Lemma \ref{lem:Hwstar}). 

Conversely, let  $e_w$ be $\bp$-dominant and  $\alpha\in \Pi^\bp$.
 Now $we_w$ is dominant and $w^{-1}$ is minimal amongst the $z$ with $zwe_w=e_w$. Thus if $s_\alpha e_w= e_w$ then $s_\alpha w^{-1}>w^{-1}$.
If $s_\alpha e_w\neq e_w$ then $P(e_w)$ is no $\bp_\alpha$-module, because its $\fb$-socle has weight $e_w$. Thus $\ind_\fb^{\bp_\alpha}P(e_w)={\sf H}_{s_\alpha}{\sf H}_{ w^{-1}}P(we_w)$ with $s_\alpha\star  w^{-1}\neq  w^{-1}$
and
$s_\alpha w^{-1}>w^{-1}$. 
  We have $w s_\alpha>w$ in either case, so $w\in W^\bp$.

\end{proof}

\begin{lemma}\label{lem:indpreserve}Let $\lambda\in X^\bp_+$.
Let $N\in\rep(\fb)$ be such that all weights $\mu$ of $N$ satisfy $\mu\leq_a\lambda$.
Then all weights $\mu$ of $\ind_\fb^\bp(N)$ satisfy $\mu\leq_a\lambda$ and the natural map $\ind_\fb^\bp(N)\to N$ 
induces an isomorphism $\ind_\fb^\bp(N)_\lambda\to N_\lambda$.
\end{lemma}

\begin{proof}
Recall there is a a reduced decomposition of $w_0^{\bp}=s_{1}\cdots s_{l}$ of $w_0^\bp$ and that $\res^\bp_\fb\ind_\fb^\bp=\sH_{s_{1}}\circ\cdots\circ\sH_{s_{l}}$. 
 Therefore we may assume $\bp$ is a minimal parabolic, say
$\bp=\bp_\alpha$. Thus $\bp/\fb$ is a projective line and we may use \cite[II Proposition 5.2]{Jan}.
If $\mu$ is a weight and $(\alpha^\vee,\mu)\geq0$,
then $s_\alpha \mu\leq_a \mu$, $\RR^1\ind_\fb^\bp(\mu)=0$, 
$(\ind_\fb^P(\mu))_\mu=\sk_\mu$ and all weights $\nu$ of $\ind_\fb^P(\mu)$ lie on the line segment joining $\mu $ with 
$s_\alpha \mu$. In particular, they satisfy $\nu\leq_a\mu$. And if $(\alpha^\vee,\mu)<0$,
then $\ind_\fb^P(\mu)=0$ and all weights $\nu$ of
$\RR^1\ind_\fb^P(\mu)$ lie strictly between the endpoints of the line segment. Thus $\nu<_a\mu$.
If $N$ has just weight $\lambda$, then it is  well known that $\ind_\fb^\bp(N)_\lambda\to N_\lambda$
is an isomorphism \cite[Propositions 17, 20]{FvdK}. 
Now use induction on the number of weights of $N$.  
\end{proof}



 %--------------------------------------------------------------
\subsection{Generating $\rep( \bp )$}\label{subsec:repP-generation}
%--------------------------------------------------------------

  We say that a full subcategory of an abelian category has the 2 out of 3 property if,
  whenever $0\to N_1\to N_2 \to N_3 \to 0$ is exact and two of the $N_i$ are in the subcategory, then so is the third.
  
  \subsection{The set $\{\hat M_v\}_{v\in W^\bp}$}
  For $\bp$-dominant $\lambda$ we use $\hat\nabla_\lambda$ as another notation for $\ind_\fb^\bp\sk_\lambda$. 
   Let us be given a set $\{\hat M_v\}_{v\in W^\bp}$ of objects of $\rep_\fr( \bp )$  with the following properties.
   For each ${v\in W^\bp}$  there is a diagram of $\bp$-modules
   $$\hat M_v\stackrel{g}\longrightarrow \hat N_v\stackrel{f}\longleftarrow \hat\nabla_{e_v} $$
   such that every weight $\mu$ of $\kernel(f)$, $\kernel(g)$, $\coker(f)$, $\coker(g)$ satisfies $\mu<_a e_v$.
  
  
  Examples of possible choices of $\hat M_v$ are $\hat\nabla_{{e_v}}$ with $f=g=\id$, $\ind_\fb^\bp Q(e_v)$ with $g=\id$, $P(-e_v)^*$
  with $f=\id$. 
  
  The multiplicity of the weight $e_v$ is one in $\hat M_v$, $\hat N_v$, $\hat\nabla_{e_v}$. 
  Every weight $\lambda$ of $\hat M_v$ satisfies $\lambda\leq_a e_v$.

\

  The following Theorem is similar to \cite[Theorem 2]{Ana}, which is proved directly in the context of $\bg$ and $\bp$. Our argument differs in that it refers back to the proof for the $\bp=\fb$ case.
  \begin{theorem}[Generation]\label{th:Pgeneration}
  The smallest strictly full additive subcategory of $\rep(\bp)$ that 
  \begin{itemize}
  \item contains the $\hat M_v$,  
  \item has the
  2 out of 3 property and 
  \item contains with every $\hat\nabla_\lambda$ also
   $\hat\nabla_\lambda\otimes\nabla_{\omega_i}$ for every fundamental representation $\nabla_{\omega_i}$,
  \end{itemize}
   is the category $\rep( \bp )$ of finite dimensional $ \bp $-modules.

  \end{theorem}  
  
  More specifically
  
  \begin{theorem}[Generation of initial interval]\label{th:intervalgeneration}
  Let $\tau\in X^\bp_+$. Let $V\in\rep(\bp)$ such that every weight $\mu$ of $V$ satisfies $\mu<_a\tau$.
Then  $V$ is an object of
  the smallest strictly full additive subcategory $\tS$ of $\rep(\bp)$ 
  satisfying 
  \begin{itemize}
  \item $\tS$ contains $\hat M_v$ for $v\in W^\bp$ with $e_v<_a\tau$,  
  \item  $\tS$ has the
  2 out of 3 property,
 \item  If  $\tS$ contains $\hat\nabla_\lambda$, then $\tS$ also contains
   $\hat\nabla_\lambda\otimes\nabla_{\omega_i}$, for every fundamental representation $\nabla_{\omega_i}$. 
  \end{itemize}
  \end{theorem}  

\begin{proof}
We argue by induction along the well-ordered partial order $\leq_a$. So we assume the Theorem when $\tau$ is replaced by 
a $\sigma\in X_+^\bp$ with $\sigma<_a\tau$. Let $\sigma\in X_+^\bp$ with $\sigma<_a\tau$.
Thus $V\in \tS$ when all weights $\xi$
of $V$ satisfy $\xi<_a\sigma$.

Step 1. First we  wish to show that $\tS$ contains at least one $\bp$-module $N$ with $N_\sigma=\sk_\sigma$ and with weights $\xi$
that satisfy $\xi\leq_a\sigma$.
If $\sigma$ is a Steinberg weight $e_v$, then $v\in W^\bp$ by Lemma \ref{lem:PisP}
and we simply take $N=\hat M_v$.
If $\sigma$ is not a Steinberg weight,  then by Corollary \ref{cor:if not Steinberg then tensor}
 there is a fundamental weight $\omega_j$ and a weight $\mu$ with $(\mu,\mu)<(\sigma,\sigma)$,
 so that the $\fb$-module $L=\nabla_{\omega_j}\otimes e^\mu$  has
$\sigma$ as a weight  of  multiplicity one, and so that all weights  $\nu$ of $L$ satisfy 
$\nu\leq_a\sigma$. Now take $N=\ind_\fb^\bp L=\nabla_{\omega_j}\otimes\ind_\fb^\bp e^\mu$. 
By Lemma \ref{lem:indpreserve} every weight $\nu$ of $N$ satisfies $\nu\leq_a\sigma$ and 
$N_\sigma=\sk_\sigma$. In particular, $\ind_\fb^\bp e^\mu$ is nonzero, so  $\mu\in X^\bp_+$ and $\ind_\fb^\bp e^\mu=\hat\nabla_\mu$.
And every weight $\nu$ of $\hat\nabla_\mu$ satisfies $(\nu,\nu)<(\sigma,\sigma)$, hence $\nu<_a\sigma$, so $\hat\nabla_\mu$  lies in $\tS$ and  
$N=\nabla_{\omega_j}\otimes\hat\nabla_\mu$ lies in $\tS$.

Step 2. Next we wish to show that $\hat\nabla_\sigma$ lies in $\tS$. To this end we look for more $\bp$-modules $N$ 
in $\tS$ with $N_\sigma=\sk_\sigma$ 
and such that the weights $\xi$ of $N$
satisfy $\xi\leq_a\sigma$.
Start with the $N$ from Step 1.
Let $N^1$ be the span of the weight spaces $N_\mu$ with $(\mu,\omega_\alpha)<(\sigma,\omega_\alpha)$ for at least one
$\alpha\in \Pi^{\notin \bp}$. Then $N^1$ is a $\bp$-submodule that lies in $\tS$. So we may replace $N$ with
$N/N^1$ and further assume $N^1=0$. 
Let $N^2$ be the span of the weight spaces $N_\mu$ with $(\mu,\omega_\alpha)\leq(\sigma,\omega_\alpha)$ for all
$\alpha\in \Pi^{\notin \bp}$. Then $N^2$ is a $\bp$-submodule and $N/N^2$ lies in $\tS$. So we may replace $N$ with
$N^2$ and further assume $N^1=0$, $N=N^2$. 
Let $N^3$ be the $\bp$-submodule of $N$ generated by $N_\sigma$.
Then $N/N^3$ is in $\tS$, so we may replace $N$ with $N^3$. Now the unipotent radical of $\bp$ acts trivially
on $N$ and the projection $p$ of $N$ onto its weight space $N_\sigma$ is $\bt$-equivariant and $\bg_\bp$-equivariant,
hence $\bp$-equivariant. By Lemma \ref{lem:indpreserve} the map $N\to \hat\nabla_\sigma$ induced by $p$ has kernel and cokernel in $\tS$. So $\hat\nabla_\sigma$ lies  in $\tS$.


Step 3.
Finally we want to show that $\tS$ contains every  $\bp$-module $V$ all whose weights $\mu$ satisfy $\mu<_a\tau$.
Consider such a $V$.
We may and shall assume  that $\tS$ contains every  $\bp$-module
 whose set of weights is a proper subset of the set of weights of $V$.
Say $V$ is nonzero. Choose an extremal weight $\sigma$ of $V$ that is $\bp$-dominant.
Let $V^1$ be the span of the weight spaces $V_\mu$ with $(\mu,\omega_\alpha)<(\sigma,\omega_\alpha)$ for at least one
$\alpha\in \Pi^{\notin \bp}$. This is a $\bp$-submodule, and if $V^1$ is nonzero, then $V$ is in $\tS$ because both
$V^1$ and $V/V^1$ are. 
Let $V^2$ be the span of the weight spaces $V_\mu$ with $(\mu,\omega_\alpha)\leq(\sigma,\omega_\alpha)$ for all
$\alpha\in \Pi^{\notin \bp}$. This is a $\bp$-submodule, and if $V^2\not= V$, then $V$ is in $\tS$ because both
$V^2$ and $V/V^2$ are. 
So we further assume $V^1=0$ and $V=V^2$. Then  the unipotent radical of $\bp$ acts trivially
on $V$. Let $V^3$ be the $\bp$-submodule generated by $V_\sigma$.
If $V^3\not= V$, then $V$ is in $\tS$ because both
$V^3$ and $V/V^3$ are. So we may assume $V=V^3$. The weights $\xi$ of $V$ now
satisfy $\xi\leq_a\sigma$.
The projection $p$ of $V$ onto its weight space $V_\sigma$ is $\bt$-equivariant and $\bg_\bp$-equivariant,
hence $\bp$-equivariant. We have  $\ind_\fb^\bp(V_\sigma)=\hat\nabla_\sigma\otimes (V_{\sigma})_\triv$, 
where $(V_{\sigma})_\triv$ is $V_{\sigma}$
with trivial $\bp$-action, cf.\ \cite[Proposition 17]{FvdK}. Consider the map $f:V\to \ind_\fb^\bp(V_\sigma)=\hat\nabla_\sigma\otimes (V_{\sigma})_\triv$, corresponding with $p$. Both $\kernel(f)$ and $\coker(f)$ lie in $\tS$. 
As $\hat\nabla_\sigma$ lies in $\tS$, so does $\hat\nabla_\sigma\otimes (V_{\sigma})_\triv$.
So $V$ lies in $\tS$.
\end{proof}


\begin{remark}
    In  Theorem \ref{th:Pgeneration}
    and
    Theorem \ref{th:intervalgeneration} we deal with $\rep(\bp)$,
    even over $\Z$,
    but the Corollary is about $R(\bp)$ as in the paper \cite{Ana}. Of course it is easier to work in $R(\bp)$, but basically the arguments are the same and we get the same understanding of Steinberg weights and a similar constructive decomposition of elements of $R(\bp)$ over $R(\bg)$ as in the proofs of Ananyevskiy.
\end{remark}


\begin{definition}
For $v\in W^\bp$ put $\hat Q(e_v)=\ind_\fb^\bp Q(e_v)$. And then for $p\in W^\bp$
\

\begin{itemize}

\item $\hat{\sf Q}_{\succeq p}:=\hull(\{\nabla_\lambda\otimes \hat Q(e_v)\}_{v\succeq p,  v\in W^\bp,\lambda\in X( \bt)_+})$




\item $\hat{\sf Q}_{\succ p}:=\hull(\{\nabla_\lambda\otimes \hat Q(e_v)\}_{v\succ p,  v\in W^\bp,\lambda\in X( \bt)_+})$





\item $\hat{\sf P}_{\preceq p}:=\hull(\{\nabla_\lambda\otimes P(-e_v)^*\}_{v\preceq p,  v\in W^\bp,\lambda\in X( \bt)_+})$



\item $\hat{\sf P}_{\prec p}:=\hull(\{\nabla_\lambda\otimes P(-e_v)^*\}_{v\prec p,  v\in W^\bp,\lambda\in X( \bt)_+})$

\end{itemize}
 \end{definition}

We get from Theorem \ref{th:Pgeneration} the following replacement of Theorem \ref{th:derived_generation} 

\begin{theorem}\label{th:Pderived_generation}
Given a $p\in W^\bp$, the triangulated hull in $\Dd(\Rep(\bp))$ of the two categories $\hat{\sf P}_{\preceq p}$ and 
$\hat{\sf Q}_{\succ p}$ 
is $\Dd^b (\rep ( P ))$. 
\end{theorem} 

 
 We also have the following replacement of Theorem \ref{th:indPQ}
 
 \begin{theorem}\label{th:PindPQ}Let $\sk$ be a field or $\Z$.
Let $v,w\in W^\bp$.
\begin{enumerate}
\item $\RR^i\ind_\bp^\bg (P(-e_v)\otimes \hat Q(e_v))=\begin{cases}\sk &\text{if $i=0$,}\\
0&\text{else.}\end{cases}$
\item If $w\not\leq v$ then $\RR^i\ind_\bp^\bg (P(-e_v)\otimes \hat Q(e_w))=0$ for all $i$.
\end{enumerate}
\end{theorem}

\begin{proof}By Theorem \ref{th:universal coefficients}, it suffices to treat the case where $\sk$ is a field.
By \cite[Theorem 1.9.(a)(ii)]{WvdK} the module $Q(e_w)$ is $\ind_\fb^\bp$ acyclic. Lemma \ref{lem:PisP}, together with the Generalized Tensor Identity \cite[I Proposition 4.8]{Jan}, give
$$\RR^j\ind_\fb^\bp (P(-e_v)\otimes Q(e_w))=P(-e_v)\otimes\RR^j\ind_\fb^\bp ( Q(e_w))=\begin{cases}P(-e_v)\otimes\hat  Q(e_w)&\text{if $j=0$,}\\
0&\text{else.}\end{cases}$$
So
$\RR^i\ind_\fb^\bg (P(-e_v)\otimes  Q(e_w))=\RR^i\ind_\bp^\bg (P(-e_v)\otimes \hat Q(e_w))$.
But the left hand side is known from Theorem \ref{th:indPQ}.
\end{proof}
\begin{proof}[\rm \bf Proof of Theorems \ref{th:Z-X^P_p-G-exceptional-coll},  \ref{th:Psemi-orthogonal}, \ref{th:Pcoh-semi-orthogonal}]
    

With these replacements in hand, we now proceed as in the case $\bp=\fb$:
For $p\in W^\bp$, let $\hat X_p$ be the image of $P(-e_p)^*$ under the left adjoint of the inclusion of $\hat{\sf Q}_{\succeq p}$ 
into $\D=\Dd ^b(\rep( \bp))$. 
And let $\hat Y_p$ be the image of $\hat Q(e_p)$ under the 
right adjoint of the inclusion of
$\hat{\sf P}_{\preceq p}$ into $\D$. 

Let $f_p$ be the natural map from $P(-e_p)^*$ to $Q(e_p)$. Recall that $f_p$ factors as a surjection $P(-e_p)^*\to\sk_{e_p}$, followed by an injection $\sk_{e_p}\to Q(e_p)$. Let $\hat f_p:P(-e_p)^*\to \hat Q(e_p)$ be induced by $f_p$. Then $\hat f_p$ induces an isomorphism of $\bt$-modules 
$(P(-e_p)^*)_{e_p}\to (\hat Q(e_p))_{e_p}$ and the weights $\mu$ of $\kernel (\hat f_p)$, $\coker(\hat f_p)$ satisfy $\mu<_a e_p$.
We have an exact triangle 
$$\dots \rightarrow \ker(\hat f_p)[1]\to \cone(\hat f_p)\to \coker(\hat f_p)\rightarrow \dots$$ in $\Dd ^b(\rep( \bp))$.
Take $\hat M_v=Q(e_v)$ for $p\prec v\in W^\bp$ and $\hat M_v=P(-e_v)^*$ for  $p\succ v\in W^\bp$. It does not matter what we choose for $\hat M_p$ itself.

Using Theorem \ref{th:intervalgeneration} instead of Theorem \ref{th:cone generation}, we find that 
$\cone(\hat f_p)$ lies in the triangulated hull of $\{\hat M_v\otimes \nabla_\lambda \mid v\in W^\bp,\ e_v <_a e_p,\ \lambda\in X(\bt)_+\}$. 
So $\cone(\hat f_p)$ lies in the triangulated hull of $\hat{\sf P}_{\prec p}\cup\hat{\sf Q}_{\succ p}$.
One uses this to show that $\hat X_p$ and $\hat Y_p$ are isomorphic.
Mutatis mutandis the old constructions and proofs go through and 
Theorems \ref{th:Z-X^P_p-G-exceptional-coll}, \ref{th:Psemi-orthogonal}, \ref{th:Pcoh-semi-orthogonal} follow.
\end{proof}
Remains to prove Theorem \ref{th:P and B}.
\begin{lemma}\label{lem:inducing Q}
Let $s$ be a simple refection and $\bp_s$ the corresponding minimal parabolic. 
\begin{itemize}
\item If $s\lambda<_d\lambda$ then $0\to Q(s\lambda)\to{\sH_s}(Q(\lambda))\to Q(\lambda)\to 0$ is exact,
\item If $s\lambda=\lambda$ then ${\sH_s}(Q(\lambda))=Q(\lambda)$,
\item If $s\lambda>_d\lambda$ then ${\sH_s}(Q(\lambda))=0$.
\end{itemize}
\end{lemma}
\begin{proof}
First assume $s\lambda\leq_d\lambda$. Then $\lambda$ is $\bp_s$-dominant. The extremal weights of ${\sH}_s(Q(\lambda))$ are the extremal
weights  $\lambda$, $s\lambda$ of ${\sH}_s(\sk_\lambda)$, because the weights $\nu$ of $Q(\lambda)$ different from $\lambda$ are too short to contribute, cf. proof of Lemma \ref{lem:indpreserve}.
By definition $Q(\lambda)$ is the kernel of a surjective map
 $P(\lambda)\rightarrow {\rm H}^0(\partial X_w,\Ll (\lambda ^+))$, and we get by  Lemma \ref{lem:kerdecomp} a relative Schubert filtration on 
 ${\sH}_s(Q(\lambda))$, with sections described by the extremal weights. Moreover,  the submodule ${\sH}_s(\sk_\lambda)$ of
${\sH}_s(Q(\lambda))$ has the $\fb$-socle $\sk_{s\lambda}$, so we know in which order the sections of the filtration occur. The first two parts follow.




Now assume $s\lambda>_d\lambda$.
 By \cite[Theorem 1.9, (ii)]{WvdK}, $Q(\lambda)$ is acyclic for ${\sH}_w,w\in W$. Applying ${\sH}_s$ to the exact sequence $0\to Q(\lambda)\to{\sH}_s(Q(s\lambda))\to Q(s\lambda)\to 0$, and using the idempotency ${\sH}_s{\sH}_s={\sH}_s$ of ${\sH}_s$'s, we conclude that ${\sH}_s(Q(\lambda))=0$.
\end{proof}


\begin{lemma}\label{lem:ind of Q(e_w)}
Let $s$ be a simple refection and $\bp_s$ the corresponding minimal parabolic. 
Let $w\in W$.
\begin{itemize}
\item If $se_w<_de_w$ then $ws\succ w$ and $0\to Q(e_{ws})\to{\sH_s}(Q(e_w))\to Q(e_w)\to 0$ is exact,
\item If $se_w=e_w$ then ${\sH_s}(Q(e_w))=Q(e_w)$,
\item If $se_w>_de_w$ then ${\sH_s}(Q(e_w))=0$.
\end{itemize}
\end{lemma}
\begin{proof}Recall that $w^{-1}$ is a minimal coset representative of the stabilizer in $W$ of the dominant weight $we_w$.
    If $se_w<_de_w$, then $se_w>_ee_w$ by Lemma \ref{lem:aed}, so $\ell(sw^{-1})=\ell(w^{-1})+1$ and $sw^{-1}$ is also a minimal coset representative of the stabilizer in $W$ of the dominant weight $we_w$.
    So $sw^{-1}\alpha>0$ for simple roots $\alpha$ perpendicular to  $we_w$. It follows that $se_w$
    is the Steinberg weight $e_{ws}$. The other points are clear from Lemma \ref{lem:inducing Q}.
\end{proof}
\begin{lemma}\label{lem:indP of Q(e_w)}
Let $w\in W^\bp$.
Then $\hat Q(e_w)\to Q(e_w)$ is surjective and its kernel lies in ${\sf Q}_{\succ p}$.
\end{lemma}
\begin{proof}
Observe that $\hat Q(e_w)$ is nonzero by Theorem \ref{th:PindPQ}. Recall there is a a reduced decomposition  $w_0^{\bp}=s_{1}\cdots s_{l}$ of $w_0^\bp$ and that $\res^\bp_\fb\ind_\fb^\bp=\sH_{s_{1}}\circ\cdots\circ\sH_{s_{l}}$.  Therefore we may assume $\bp$ is a minimal parabolic and apply Lemma \ref{lem:ind of Q(e_w)}, taking into account that all modules $Q(\lambda)$ are 
$\sH_{s}$ acyclic for simple $s$.
\end{proof}

\begin{lemma}\label{lem:inclusion of hat}
For $w\in W^\bp$ one has $\hat{\sf Q}_{\succ p}\subset {\sf Q}_{\succ p}$.
\end{lemma}
\begin{proof}
By descending induction along $\succ$, using Lemma \ref{lem:indP of Q(e_w)}.
\end{proof}

\begin{proof}[\rm\bf Proof of Theorem \ref{th:P and B}]
Take $p\in W^\bp$. By equation (\ref{eq:def_of_Y_p}) we have 
\begin{equation}
    Y_p={i^p}_{\preceq p}^{\ !}(Q(e_p)).
\end{equation}
 By Lemma \ref{lem:indP of Q(e_w)} we have 
\begin{equation}
{i^p}_{\preceq p}^{\ !}(Q(e_p))={i^p}_{\preceq p}^{\ !}(\hat Q(e_p)).\end{equation}
As ${\sf cone}(\hat Y_p\rightarrow \hat Q(e_p))\in \hat{\sf P}_{\preceq p}^{\perp}=
\hat{\sf Q}_{\succ p}\subset  {\sf Q}_{\succ p}$, we also have \begin{equation}
{i^p}_{\preceq p}^{\ !}(\hat Q(e_p))={i^p}_{\preceq p}^{\ !}(\hat Y_p).    
\end{equation} 
Further $\hat Y_p=\hat X_p\in \hat{\sf P}_{\preceq p}\subset {\sf P}_{\preceq p}$, so that 
\begin{equation}
{i^p}_{\preceq p}^{\ !}(\hat Y_p)=\hat Y_p.    
\end{equation}
Taken together, 
$Y_p={i^p}_{\preceq p}^{\ !}(Q(e_p))={i^p}_{\preceq p}^{\ !}(\hat Q(e_p))={i^p}_{\preceq p}^{\ !}(\hat Y_p)=\hat Y_p$.
\end{proof}

\subsection{A pairing}
Let $X=\bg/\fb$, $Y=\bg/\bp$ and let $\pi:X\to Y$ be the  projection. It induces a
map\\ $\pi^*:K_\bt(Y)\to K_\bt(X)$ and therefore, by subsection \ref{subsec:alphabet}, a pairing 
$$K_\bt(Y)\times K_\bt(X)\to R(\bt), \quad ([\F],[\G])\mapsto \langle \pi^*[\F],[\G]\rangle.$$
%Let $w_0^P$ be the longest element of $W_P$.
For $v,w\in W^\bp$ with $vw_0\not\leq ww_0$
the pairing sends $([\Ll_Y(P(-e_v))],[\Oo_{\fX^{ww_0}}(-\partial \fX^{ww_0})])$ to $\beta_{v,ww_0}=0$.
And it sends $([\Oo_{\overline{\fb^+v\bp}/\bp}],[\Oo_{\fX^{ww_0}}(-\partial \fX^{ww_0})])$ to $\delta_{vw}$ for $v,w\in W^\bp$.
If $\sk=\cC$ it follows  that $[\Ll_Y(P(-e_v))]$ is in the $R(\bt)$-span of the Schubert classes $[\Oo_{\overline{\fb^+w\bp}/\bp}]$ with $w\in W^\bp$, $w\leq v$.

%$B^+wP/B=B^+wB/B=B^+wx_-=B^+ww_0x+$.

%----------------------------------------------------------------------
\section{Explicit calculations in rank 2}\label{sec:rank_2}
%----------------------------------------------------------------------
This section works out in detail the case of flag varieties of rank two groups. Unless 
specified otherwise, we work over $\mathbb Z$. The ensuing full exceptional collections consist of pure objects. These collections have the property that their duals in the sense of the definition of Section \ref{subsec:dualFEC} are given by by applying the Grothendieck duality to the initial ones followed by tensoring with the square root of the canonical sheaf $\Ll (2\rho)$.

\

We first record several fundamental vanishing theorems that will be used throughout this section.

%----------------------------------------------------------------
\subsection{Borel-Weil-Bott theorem}\label{th:BWB_van}
%----------------------------------------------------------------

Define $$\overline C _{\mathbb Z}=\{\lambda\in X({\bt}) \ | \ 0\leq \langle \lambda + \rho,\beta ^{\vee}\rangle\} \ {\rm for \ all} \ \beta \in {\rm R}_+$$ if ${\rm char}(k)=0$ and $$\overline C _{\mathbb Z}=\{\lambda\in X({\bt}) \ | \ 0\leq \langle \lambda + \rho,\beta ^{\vee}\rangle\}\leq p \ {\rm for \ all} \ \beta \in {\rm R}_+$$ if ${\rm char}(k)=p>0$.

\begin{theorem}\cite[Part II, Corollary 5.5]{Jan}\label{th:Bott-Demazure_th}

\begin{itemize}[(a)]

\item[(a)] If $\lambda \in \overline C _{\mathbb Z}$ with $\lambda \notin X({\bt})_+$, then 

$$
{\rm H}^i({ \bg}/{\fb},\Ll _{w\cdot\lambda})=0
$$

for all $w\in W$. Here $\cdot$ is the dot--action of the Weyl group $W$ on 
$X({\bt})$.

\

\item[(b)] If $\lambda \in  \overline C _{\mathbb Z}\cap  X({\bt})_+$, then 
for all $w\in W$ and $i\in \mathbb N$

$$
{\rm H}^i({ \bg}/{\fb},\Ll _{w\cdot \lambda})={\rm H}^0({ \bg}/{\fb},\Ll _{\lambda}),
$$

\

if $i=l(w)$ and otherwise ${\rm H}^i({ \bg}/{\fb},\Ll _{w\cdot \lambda})=0$. Here $l(w)$ is the length function.

\end{itemize}

\end{theorem}

%----------------------------------------------------------------
\subsection{Kempf vanishing theorem}\label{subsec:Kempfvan}
%----------------------------------------------------------------

\begin{theorem}\cite[Part II, Proposition 4.5]{Jan}\label{th:Kempf_van} and \cite[Section 2.6]{Donkin2}
If $\lambda \in X({\bt})_+$, then 

$$
{\rm H}^i({ G}/{\fb},\Ll _{\lambda})=0
$$
for all $i>0$.

\end{theorem}

By \cite[Corollary 2.6, (ii)]{AndChev}, $
{\rm H}^0({\mathbb G}/{\mathbb B},{\mathbb L} _{\lambda})$ are free $\mathbb Z$-modules for $\lambda \in X({\Tm})_+$. Combined with 
the Universal coefficient Theorem \ref{th:universal coefficients} and 
Theorem \ref{th:Kempf_van}, this gives ${\rm H}^i({\mathbb G}/{\mathbb B},{\mathbb L} _{\lambda})=0$ for $i>0$.

%----------------------------------------------------------------
\subsection{Kodaira vanishing theorem}
%----------------------------------------------------------------

\begin{theorem}\cite[Theorem 1.2.9]{BK}\label{th:Kodaira_van}
Let $k$ be a field of characteristic $p>0$ and $X$ be a Frobenius split $k$-variety. Assume given an ample line bundle $\Ll$ on $X$. Then $H^i(X,\Ll \otimes \omega _X)=0$ for $i>0$.
\end{theorem}

Generalized flag varieties $ \bg/\bp$ are Frobenius split, \cite[Theorem 2.2.5]{BK}, thus Kodaira vanishing holds for these.

\begin{corollary}\label{cor:excseqlinebun}
Let $X$ be a Frobenius split Fano $k$-variety such that ${\rm Pic}(X)=\mathbb Z$ with $\Ll\in {\rm Pic}(X)$ being the positive generator. Let $n$ be the index of $X$, i.e.\ $\omega _{X}^{-1}=\Ll ^{\otimes n}$. Then the collection of line bundles $\Oo _X,\Ll,\dots, \Ll ^{\otimes (n-1)}$ is an exceptional collection on $X$.
\end{corollary}

\begin{proof}
Setting $\Ll=\omega _{X}^{-1}$ in Theorem \ref{th:Kodaira_van}, we obtain $H^i(X,\Oo _X)=0$ for $i>0$. It follows that for any line bundle $\mathcal M$ on $X$ one has $\Ext _X^i({\mathcal M},{\mathcal M})=H^i(X,\Oo _X)=0$ for $i>0$, hence the collection $\Oo _X,\Ll,\dots, \Ll ^{\otimes (n-1)}$ consists of exceptional bundles. The same theorem implies that $\Ext _X^k(\Ll ^{\otimes j},\Ll ^{\otimes i})=H^k(X,\Ll ^{\otimes (i-j)}))=0$ for $i<j$ and $k>0$ since $-n+1\leq i-j<0$. Finally, $\Hom _X(\Ll ^{\otimes j},\Ll ^{\otimes i})=H^0(X,\Ll ^{\otimes (i-j)})=0$ since the bundle $\Ll ^{\otimes (j-i)}$ is ample and hence has a non-trivial section.
\end{proof}

%---------------------------------------------------------------
\subsection{Demazure functors}\label{subsec:Demfun}
%---------------------------------------------------------------

By \cite[Section 5.17-5.18]{Jan}, the Joseph-Demazure functors ${\sH}_w$ from Section \ref{subsec:J-D} can be related to functors which are defined  geometrically. Let $R$ be the root system of a simple simply connected group $ \bg$ of rank $r$ and $\Pi=\{\alpha _1,\dots ,\alpha _r\}\subset R$ the simple roots. 
Let $s_{\alpha _i}, i=1,\dots,r$ be the associated simple reflections of the Weyl group $W$. Given a simple root $\alpha$, let $\pi _{\alpha}: { \bg}/{\fb}\rightarrow { \bg}/{\bp}_{\alpha}$ denote the projection. Associated to $\alpha$ is an endofunctor ${\rm R}^0{\sf D}_{\alpha}:=\pi _{\alpha}^{\ast}{\rm R}^0{\pi _{\alpha}}_{\ast}$ of ${\QCoh}^{ \bg}({ \bg}/{\fb})$ (the functor $\pi _{\alpha}^{\ast}$ is exact, the morphism 
$\pi _{\alpha}$ being smooth). Let $w$ be an element of the Weyl group. For a reduced expression of $w = s_{\alpha _1}s_{\alpha _2}\dots s_{\alpha _n}$ consider the functor ${\rm R}^0{\sf D}_w ={\rm R}^0{\sf D}_{\alpha _1}\circ \dots \circ {\rm R}^0{\sf D}_{\alpha _n} = \pi _{\alpha _1}^{\ast}{\rm R}^0{\pi _{\alpha _1}}_{\ast}\dots \pi _{\alpha _n}^{\ast}{\rm R}^0{\pi _{\alpha _n}}_{\ast}$. The functor ${\rm R}^0{\sf D}_w $ is also a left exact functor.

Similarly to the case of ${\sH}_w$, let ${\sf D}_w={\sf D}_{\alpha _1}\circ \dots \circ {\sf D}_{\alpha _n} = \pi _{\alpha _1}^{\ast}{\pi _{\alpha _1}}_{\ast}\dots \pi _{\alpha _n}^{\ast}{\pi _{\alpha _n}}_{\ast}$ denote the right derived functor of ${\rm R}^0{\sf D}_w$; it is an endofunctor of $\Dd ^{+}({\QCoh}^{ \bg}({ \bg}/{\fb}))$\footnote{Note that the functor ${\sf D}_{w}$ is also defined on the derived category of non-equivariant coherent sheaves on $ \bg/\fb$ and eventually will be considered on that category.}. Since the morphisms $\pi _{\alpha _i}$'s are smooth and proper, the functors ${\sf D}_w$'s restrict to endofunctors of the bounded derived category $\Dd ^b({\sf Coh}^{ \bg}({ \bg}/{\fb}))$\footnote{The same footnote as above applies.}.

Now set ${ \bg}:={\fb}, X={ \bg}$ (res.,  ${ \bg}:={\bp}_{\alpha}, X={ \bg}$) in the definition of the sheafification functor $\Ll$ and denote $\Ll _{\fb}$ (resp., $\Ll _{{\bp}_{\alpha}}$) the corresponding functors ${\Rep}({\fb})\rightarrow {\QCoh}^{ \bg}({ \bg}/{\fb})$ (resp., ${\Rep}({\bp}_{\alpha})\rightarrow {\QCoh}^{ \bg}({ \bg}/{\bp}_{\alpha}$)).

\begin{proposition}\label{prop:Joseph-to-Demazure}
Let $M\in \Dd ^+({\Rep}( \fb))$ (resp., $M\in \Dd ^b({\rep}( \fb))$). Then there is an isomorphism $\Ll _{\fb}({\rm R}{\sH}_w(M))={\sf D}_w\Ll _{\fb}(M)$.
\end{proposition}

\begin{proof}
See \cite[Lemma 5]{Math}.
\end{proof}

%---------------------------------------------------------------
\subsection{Demazure resolutions}\label{subsec:Demres}
%---------------------------------------------------------------

Let $w\in W$ and $s_{\alpha _i}s_{\alpha _{i_2}}\dots s_{\alpha _{i_k}}$ be a reduced expression of $w$. Associated to $s_{\alpha _i}s_{\alpha _{i_2}}\dots s_{\alpha _{i_k}}$ is the Bott-Samelson-Demazure-Hansen resolution $Z_w$ of the Schubert variety $X_w$ associated to $w$, \cite[Section 2.2]{BK} and \cite[Section 1.2]{WvdKTata}. Both $X_w$ and $Z_w$ are acted upon by $\fb$ and one can form $ \bg$-varieties by setting ${\mathcal X}_w= \bg\times ^{\fb}X_w\subset X\times X$ and $\Zz _w= \bg\times ^{\fb}Z_w$. Then $\Zz _w$ is a resolution of singularities of ${\mathcal X}_w$. Note that the variety $\Zz _w$ depends on the choice of a reduced expression of $w$.

Given two smooth varieties $X$ and $Y$ over $\sk$ and an object 
$\mathcal P \in \Dd ^b(X\times Y)$, denote $\Phi _{\mathcal P}(-): = {\pi _Y}_{\ast}(\pi _{X}^{\ast}(-)\otimes \mathcal P)$ the 
integral transform between $\Dd ^b(X)$ and $\Dd ^b(Y)$, where $\pi _{X}, \pi _{Y}$ are the projections of $X\times Y$ onto corresponding factors.

By the construction, the variety $\Zz _w$ is a closed subvariety of $X^{l(w)}$. Denote $p_w$ and $q_w$ the two projections $Z_w$ onto the first and the last factor of $X^{l(w)}$, respectively.


\begin{theorem}\label{th:Demazure_kernel}\cite[Corollary 2.2.7]{WvdKTata}

Let $w\in W$ and consider a reduced decomposition of $w$ as above. There is an isomorphism of endofunctors of $\Dd ^b(X)$: 

\begin{equation}
{\sd }_w=\Phi _{\Oo _{\tX _w}}={q_w}_{\ast}p_w^{\ast}
\end{equation}
where $\Phi _{\Oo _{\tX _w}}$ is the integral functor associated to $\Oo _{\tX _w}$.
\end{theorem}

\begin{proof}
Follows from the flat base change and 
rationality of singularities of Schubert varieties, \cite{And}.
\end{proof}


To unburden the notation, let $X$ denote ${ \bg}/{\fb}$. Occasionally though when there is a need, we will stress the notation ${ \bg}/{\fb}$. For a simple root $\alpha _i$ let ${\bp}_{\alpha _i}$ denote the minimal parabolic generated by ${\fb}$ and $\bu_{\alpha _i}$, and denote $Y_i:={ \bg}/{\bp}_{\alpha _i}$ with $\pi _{\alpha _i}: X\rightarrow Y_i$. Given a simple root $\alpha_ i$, we denote $\sd _i$ the Demazure operator associated to $\alpha_ i$, i.e.\ $\sd _i:=\sd _{\alpha _i}$.
Throughout this section, we employ the notation of Section \ref{subsec:triangulated}. The following theorem is the starting point for unraveling the category $\Dd ^b({\rep}({\fb}))$ as a module category over  $\Dd ^b({\rep}( \bg))$:

\begin{theorem}[Demazure, \cite{Demres}]\label{th:Demw0}
Let $w_0$ be the longest element of $W$.
There is an isomorphism of functors ${\sd}_{w_0}=\pi ^{\ast}\pi _{\ast}=\sd _{\Oo _X}$, where $\pi:X\rightarrow {\rm pt}$ is the projection to a point.
\end{theorem}

\begin{proof}
Setting $w=w_0$ in Theorem \ref{th:Demazure_kernel} gives the result. 
\end{proof}

The functors ${\sd }_i$ are idempotents: ${\sd }_i\circ {\sd }_i={\sd }_i$. Indeed, by the projection formula one has ${\sd }_i\circ {\sd }_i={\sd }_i\otimes {\sd }_i\Oo _X={\sd }_i$. This leads to the appearance of the $0$--Hecke monoid. The following statement is a very simple consequence of that idempotency property that will be used throughout:

\begin{proposition}\label{prop:Demazure_reduction}
Let $w\in W$ and $s_{\alpha _i}s_{\alpha _{i_2}}\dots s_{\alpha _{i_k}}$ be a reduced expression of $w$ starting from $s_{\alpha _i}$. Let $\Ee, \Ff$ be two objects of $\Dd ^b(X)$ and ${\sf D}_w$ be the Demazure functor associated to $w$. Then 
\begin{equation}\label{eq:Demazure_reduction}
\RHom _X({\sf D}_i\Ee,{\sf D}_w\Ff)=\RHom _X({\sf D}_i\Ee,{\sf D}_{s_{\alpha _i}w}\Ff).
\end{equation}
\end{proposition}

\begin{proof}
We have ${\sf D}_w={\sf D}_i{\sf D}_{s_{\alpha _i}w}$, hence $\RHom _X({\sf D}_i\Ee,{\sf D}_w\Ff)=\RHom _X({\sf D}_i\Ee,{\sf D}_i{\sf D}_{s_{\alpha _i}w}\Ff)$. Now the functor ${\pi _i}^{\ast}$ is fully faithful, hence the latter group is isomorphic to $\RHom _{Y_i}({\pi _i}_{\ast}\Ee,{\pi _i}_{\ast}{\sf D}_{s_{\alpha _i}w}\Ff)$. By adjunction $({\pi _i}^{\ast},{\pi _i}_{\ast})$ we see that the latter group is isomorphic to $\RHom _X({\sf D}_i\Ee,{\sf D}_{s_{\alpha _i}w}\Ff)$.
\end{proof}

Let $\omega _i$ be a fundamental weight. The space of global sections $\Gamma (X,\Ll ({\omega _i}))$ is isomorphic to $\nabla _{\omega _i}$ (the induced module). Let $\omega _X$ (resp., $\omega _{Y_i}$) denote the canonical line bundle on $X$ (resp., on $Y_i$). Then ${\pi _{\alpha _i}}_{\ast}\omega _{X}=\omega _{Y_i}[-1]$. The relative canonical sheaf $\omega _{X/Y_i}$ is isomorphic to $\Ll (-\alpha _i)$ \cite[Lemme 6]{Dem}, and hence ${\pi _i}_{\ast}\Ll ({-\alpha _i})=\Oo _{Y_i}[-1]$. More generally, given an object $\Ff \in \Dd ^b(X)$, Grothendieck duality for the morphism ${\pi _{\alpha _i}}$ gives an isomorphism ${\pi _{\alpha _i}}_{\ast}(\Ff \otimes \Ll (-\alpha _i))[1]={\pi _{\alpha _i}}_{\ast}(\Ff ^{\vee})^{\vee}$ where the inner $(-)^{\vee}$ is the duality $R{\mathcal Hom}_X(\Ff ,\Oo _X)$ on $X$ and the outer $(-)^{\vee}$ is the duality $R{\mathcal Hom}_{Y_i}(\Ff ,\Oo _{Y_i})$ on $Y_i$. Since the pullback ${\pi _{\alpha _i}^{\ast}}$ commutes with $(-)^{\vee}$, we obtain 
an isomorphism ${\sd }_i(\Ff \otimes \Ll (-\alpha _i))[1]={\sd }_i(\Ff ^{\vee})^{\vee}$.

Given a line bundle $\Ll (\chi), \chi \in X({\bt})$ and a minimal parabolic ${\bp}_{\alpha}$ generated by $\fb$ and $\bu_{\alpha}$,
the degree of the restriction of $\Ll (\chi)$ to fibres of the projection $\pi _{\alpha}:{ \bg}/{\fb}\rightarrow { \bg}/{\bp}_{\alpha}$ is equal to $\langle \chi, \alpha ^{\vee}\rangle$, \cite[Lemme 2]{Demres}.

Given a fundamental weight $\omega _i$, there is a short exact sequence \cite[Section 3]{Dem}:
\begin{equation}\label{eq:Demazureses}
0\rightarrow \Ll ({\omega _i-\alpha _i})\rightarrow {\sd}_i\Ll ({\omega _i})\rightarrow \Ll ({\omega _i})\rightarrow 0.
\end{equation}

\begin{proposition}\label{prop:rightadjD_i}
The right adjoint functor ${\sd }_i^{\vee}(-)$ to ${\sd }_i$ is isomorphic to ${\sd }_i(-)\otimes \Ll (-{\alpha _i})[1]$.
\end{proposition}

\begin{proof}
There are adjunction pairs of functors $({\pi _i^{\ast}},{\pi _i}_{\ast})$ and $({\pi _i}_{\ast},{\pi _i^!})$ is an adjoint pair. Now ${\sd }_i={\pi _i^{\ast}}{\pi _i}_{\ast}$, hence its right adjoint is isomorphic to ${\sd }_i^{\vee}={\pi _i^!}{\pi _i}_{\ast}$. The Grothendieck duality for the morphism $\pi _i:X\rightarrow Y_i$ gives an isomorphism ${\pi _i^!}(-)={\pi _i^{\ast}}(-)\otimes \omega _{X/Y_i}[1]$ where $\omega _{X/Y_i}$ is the relative canonical sheaf for the projection $\pi _i$ and the shift $[1]$ appears since the relative dimension $X$ over $Y_i$ is equal to $1$. Finally, recalling the isomorphism $\omega _{X/Y_i}=\Ll (-\alpha _i)$ we obtain the result.
\end{proof}

%---------------------------------------------------------------
\subsection{The bundle $\Ll (-\rho)$} 
%---------------------------------------------------------------

Recall that the canonical sheaf $\omega _X$ is isomorphic to $\Ll (-2\rho)$ where $\rho$ is the sum of fundamental weights. The following proposition will turn to be useful in checking semiorthogonality properties of exceptional sequences.

\begin{proposition}\label{prop:Demazure_-rho-orthogonal}
Let $\Ee \in \Dd ^b(X)$ and $w\in W, l(w)\geq 1$ with the associated Demazure operator ${\sd }_w$. Then $\RHom _X({\sd}_w\Ee,\Ll (-\rho))=0$.
\end{proposition}

\begin{proof}
Let $s_{\alpha _i}s_{\alpha _{i_2}}\dots s_{\alpha _{i_k}}$ be a reduced expression of $w$. Then ${\sd }_w={\sd }_{i_1}{\sd }_{i_2}\dots {\sd }_{i_k}={\sd}_{i_1}{\sd}_{s_{\alpha _1}w}$. By adjunction $({\pi _i}^{\ast},{\pi _i}_{\ast})$, we obtain

\begin{equation}
\RHom _X({\sd}_w\Ee,\Ll (-\rho))=\RHom _{Y_1}({\pi _{i_1}}_{\ast}{\sd}_{s_{\alpha _1}w}\Ee,{\pi _i}_{\ast}\Ll (-\rho)).
\end{equation}

Now $\RHom _{Y_1}({\pi _{i_1}}_{\ast}{\sd}_{s_{\alpha _1}w}\Ee,{\pi _i}_{\ast}\Ll (-\rho))=0$  since ${\pi _i}_{\ast}\Ll (-\rho)=0$ as $\langle -\rho,\alpha _i^{\vee}\rangle =-1$.
\end{proof}

%---------------------------------------------------------------
\subsection{An exact triangle}
%---------------------------------------------------------------
Let $S$ be a smooth projective scheme and $\Ee$ be a vector bundle of rank $r$ on $S$. Consider the associated projective bundle $\pi :   \Pp (\Ee)\rightarrow S$. Denote $\Oo _{\pi}(-1)$ the line bundle on $  \Pp (\Ee)$ of relative degree $-1$, such that $\pi _{\ast}\Oo _{\pi}(1)=\Ee ^{\ast}$. 
One has \cite[Corollary 8.36]{Huyb}:

\begin{theorem}\label{th:Orlovth}
The category $\Dd ^b(  \Pp (\Ee))$ has a semiorthogonal decomposition: 

\begin{equation}
\Dd ^b(  \Pp (\Ee)) = \langle \pi ^{\ast}\Dd ^b(S)\otimes \Oo _{\pi}(-r+1),\dots ,  \pi ^{\ast}\Dd ^b(S)\otimes \Oo _{\pi}(-1),\pi ^{\ast}\Dd ^b(S)\rangle .
\end{equation}

\end{theorem}

The following proposition generalizes the sequence (\ref{eq:Demazureses}) for an arbitrary object of $\Dd ^b(X)$:

\begin{proposition}\label{prop:alpha_i-triangle}
For any object $\Ee \in \Dd ^b(X)$ there is an exact triangle 
\begin{equation}\label{eq:alpha_i-triangle}
\dots \longrightarrow {\sd }_i(\Ee \otimes \Ll (-\omega _i))\otimes \Ll (\omega _i-\alpha_i)\longrightarrow {\sd }_i\Ee \rightarrow \Ee \stackrel{[1]}\longrightarrow \dots 
\end{equation}
\end{proposition}

\begin{proof}
The triangle is obtained by Theorem \ref{th:Orlovth}. Note that the cone 
of the canonical adjunction morphism ${\sd }_i\Ee \rightarrow \Ee$ is of the form $\pi _i^{\ast}(?)\otimes \Ll (\omega _i-\alpha_i)$. To calculate the unknown object $?$ in the cone, tensor the adjunction map with $\Ll (-\omega _i)$ and apply ${\pi _i}_{\ast}$. One obtains ${\pi _i}_{\ast}{\rm Cone}({\sd }_i\Ee \rightarrow \Ee)=?[-1]={\pi _i}_{\ast}(\Ee\otimes \Ll (-\omega _i))$ giving the triangle (\ref{eq:alpha_i-triangle}).
\end{proof}

Note that the object $\Ee$ in the above proposition is not required to be $\bg$-equivariant. However, if $\Ff\in \Dd ^b(X)$ is a $\bg$-equivariant object of $\Dd ^b(X)$, i.e. $\Ff \in \Dd ^b_{\bg}(X)=\Dd ^b(\rep(\fb))$ (and therefore $\Ff=\Ll(\sf F)$ for an object ${\sf F}\in \Dd ^b(\rep(\fb))$, then the maps in the triangle (\ref{eq:alpha_i-triangle}) are $\bg$-equivariant. More specifically, in that case triangle (\ref{eq:alpha_i-triangle}) comes from an exact triangle in $\Dd ^b(\rep(\fb))$ upon applying the functor $\Ll$. Precisely, (\ref{eq:alpha_i-triangle}) is obtained from the following triangle below 
by applying Proposition \ref{prop:Joseph-to-Demazure} to it:
\begin{proposition}\label{prop:repB-alpha_i-triangle}

Let ${\sf F}$ be an object $\Dd ^b(\rep(\fb))$ and $\alpha _i\in S$ a simple root. Let ${\sH}_i$ denote the Joseph-Demazure functor ${\sH}_{s_{\alpha _i}}$. Then there is an exact triangle in $\Dd ^b(\rep(\fb))$
\begin{equation}\label{eq:repB-alpha_i-triangle}
\dots \longrightarrow {\rm R}{\sH}_i({\sf F}\otimes {\sf k}_{-\omega _i})\otimes {\sf k}_{\omega _i-\alpha_i}\longrightarrow {\rm R}{\sH}_i{\sf F} \rightarrow {\sf F} \stackrel{[1]}\longrightarrow \dots 
\end{equation}
\end{proposition}

\begin{remark}\label{rem:Q_alternate}
{Recall the formula for the functor $Q(-)$ (Section \ref{subsec:Schubfilt} and \cite[Definition 7.2.1]{WvdKTata}). Given a weight $\lambda \in X({\bf T})$, let $\lambda^+$ be the dominant weight in the Weyl group orbit of $\lambda$ and let $w\in W$ be minimal so that $\lambda=w\lambda^+$ (cf. Definition \ref{def:dualJosephmod}). Choose a reduced expression for $w=s_{\alpha _1}s_{\alpha _2}\dots s_{\alpha _k}$. For a ${\bf B}$-module $M$ and a simple root $\alpha \in S$ define the functor $F_{\alpha}(M):={\res}_{\bf B}^{{\bf P}_{\alpha}}({\mathbb Z}_{-\rho}\otimes {\ind}_{\bf B}^{{\bf P}_{\alpha}}({\mathbb Z}_{s_{\alpha}(\rho)}\otimes _{\mathbb Z}M))$\footnote{Recall that $\fb$ is generated by negative roots, so the signs of $\rho$ and of $s_{\alpha _i}(\rho)$ in each factor are opposite to those from \cite[Definition 7.2.1]{WvdKTata}.}. Set 
\begin{equation}\label{eq:Q_integral}
Q(\lambda)_{\mathbb Z}:=F_{\alpha _1}\dots F_{\alpha _k}{\mathbb Z}_{\lambda^+}.
\end{equation}

Then \cite[Proposition 7.2.2]{WvdKTata} implies that $Q(\lambda)$ defined above is the integral form of relative Schubert modules of Definition \ref{def:defineQ}.
}
\end{remark}

\

We can reconcile the exact triangle (\ref{eq:repB-alpha_i-triangle}) from Proposition \ref{prop:repB-alpha_i-triangle} with the above description of the functor $Q$. For a simple root $\alpha _i$, the functor $F_{\alpha _i}(-)={\res}_{\bf B}^{{\bf P}_{\alpha _i}}({\mathbb Z}_{-\rho}\otimes {\ind}_{\bf B}^{{\bf P}_{\alpha _1}}({\mathbb Z}_{s_{\alpha _i}(\rho)}\otimes _{\mathbb Z}(-))$
is left exact; its right derived functor is given by the cone 
${\rm R}{\sH}_i(-\otimes {\sf k}_{-\omega _i})\otimes {\sf k}_{\omega _i-\alpha_i}$
in (\ref{eq:repB-alpha_i-triangle}) of Proposition \ref{prop:repB-alpha_i-triangle}: to see it, note that $\langle s_{\alpha _i}(\rho),\alpha _i^{\vee}\rangle =\langle -\omega _i,\alpha _i^{\vee}\rangle =-1$ and $\langle -
\rho,\alpha _i^{\vee}\rangle = \langle \omega _i-\alpha _i,\alpha _i^{\vee}\rangle =-1$. So, the triangle (\ref{eq:repB-alpha_i-triangle}) may as well be written as follows:
\begin{equation}\label{eq:Q-alpha_i-triangle}
\dots \longrightarrow {\rm R}{\sH}_i({\sf F}\otimes {\sf k}_{s_{\alpha _i}(\rho)})\otimes {\sf k}_{-\rho}\longrightarrow {\rm R}{\sH}_i{\sf F} \rightarrow {\sf F} \stackrel{[1]}\longrightarrow \dots 
\end{equation}

Given a reduced expression $s_{\alpha _1}\dots s_{\alpha _k}$ of an element $w\in W$, there is an isomorphism of functors ${\rm R}{\sH}_w={\rm R}{\sH}_1\dots {\rm R}{\sH}_k$. Iterating the cones in exact triangles (\ref{eq:repB-alpha_i-triangle}) for $i=1,\dots,k$, we obtain a right derived functor of $Q$.

%----------------------------------------------------------------
\subsection{More on Demazure functors}
%----------------------------------------------------------------
We collect here various auxiliary statements about the effect of Demazure functors on objects of interest.

Let $ \bg$ be a simple simply connected algebraic group and $\omega _1,\dots,\omega _r$ be the fundamental weights of ${ \bg}$. Given a fundamental weight $\omega _i$, define the vector bundle $\Psi _1^{\omega _i}$ as the kernel of the canonical surjective morphism $\nabla _{\omega _i}\otimes \Ll (0)\rightarrow \Ll (\omega _i)$.  Thus, there is the short exact sequence:

\begin{equation}\label{eq:Psi_omega_i}
0\rightarrow \Psi _1^{\omega _i}\rightarrow \nabla _{\omega _i}\otimes \Ll (0)\rightarrow \Ll (\omega _i)\rightarrow 0,
\end{equation}
and $\Psi _1^{\omega _i}$ is the left mutation of $\Ll (\omega _i)$ through $\Ll (0)$. For referential purposes, we also record the dual sequence 

\begin{equation}\label{eq:dualPsi_omega_i}
0\rightarrow \Ll (-\omega _i)\rightarrow (\nabla _{\omega _i})^{\ast}\otimes \Ll (0)\rightarrow (\Psi _1^{\omega _i})^{\ast}\rightarrow 0
\end{equation}

We can now state a few simple propositions describing the action of elementary Demazure operators ${\sf D}_k$.

\begin{proposition}\label{prop:pi_k_*Psi _l}


\begin{equation}\label{eq:D_kPsi _lac}
{\sf D}_k(\Psi _1^{\omega _l})^{\ast}=\left\{
   \begin{array}{l}
    (\nabla _{\omega _l})^{\ast}\otimes \Ll (0), \quad k=l,\\
    \\
    (\Psi _1^{\omega _l})^{\ast}, \quad k\neq l. \\
   \end{array}
  \right.
\end{equation}
\end{proposition}

\begin{proof}
Consider the sequence (\ref{eq:dualPsi_omega_i}) for $i=l$ and apply to it the functor ${\pi _{\alpha _k}}_{\ast}$. Now observe that $\langle -\omega _l,\alpha _k^{\vee}\rangle =-\delta _{kl}$,
hence ${\sf D}_k\Ll (-\omega _k)=0$ and ${\sf D}_k\Ll (-\omega _l)=\Ll (-\omega _k)$ if $l\neq k$. 

The middle term of (\ref{eq:dualPsi_omega_i}) is the trivial bundle, hence remains unaffected by ${\sf D}_k$, and the statement follows.
\end{proof}

\begin{proposition}\label{prop:pi_k_*Psi _l(-l)a}
Let the Dynkin type of $ \bg$ be ${\sf A}_n$. Then 

\begin{equation}\label{eq:pi_k_*Psi _l(-l)b}
%{\pi _{\alpha _k}}_{\ast}
{\sf D}_k((\Psi _1^{\omega _l})^{\ast}\otimes \Ll (-\omega _l))=\left\{
   \begin{array}{l}
    \Ll (-\omega _{k-1}-\omega _{k+1}), \quad k=l,\\
    \\
    (\Psi _1^{\omega _l})^{\ast}\otimes \Ll (-\omega _l), \quad k\neq l. \\
   \end{array}
  \right.
\end{equation}

The weight $\omega _{k-1}$ (resp., $\omega _{k+1}$) is set to be $0$ if $k=1$ (resp., $k=n$).
\end{proposition}

\begin{proof}
Tensoring the sequence (\ref{eq:dualPsi_omega_i}) for $i=l$ with $\Ll (-\omega _l)$, we obtain 

\begin{equation}\label{eq:dualPsi_omega_l(-l)}
0\rightarrow \Ll (-2\omega _l)\rightarrow (\nabla _{\omega _l})^{\ast}\otimes \Ll (-\omega _l)\rightarrow (\Psi _1^{\omega _l})^{\ast}\otimes \Ll (-\omega _l)\rightarrow 0.
\end{equation}

Applying the functor ${\pi _{\alpha _k}}_{\ast}$ to it and using once again the equality $\langle -\omega _l,\alpha _k^{\vee}\rangle =-\delta _{kl}$, we see that the sequence is unaffected by ${\pi _{\alpha _k}}_{\ast}$ if $k\neq l$ which asserts the second line of (\ref{eq:pi_k_*Psi _l(-l)b}). If $k=l$ then ${\pi _{\alpha _l}}_{\ast}\Ll (-\omega _l)=0$ and we need to compute ${\pi _{\alpha _l}}_{\ast}\Ll (-2\omega _l)$. Now $-2\omega _l=
-\alpha _l-\omega _{l-1}-\omega _{l+1}$ (the Cartan matrix of ${\sf A}_n$) where $\omega _{l-1}$ (resp., $\omega _{l+1}$) is set to be $0$ if $l=1$ (resp., $l=n$). Using the isomorphism ${\pi _{\alpha _l}}_{\ast}\Ll (-\alpha _l)=\Oo _{Y_l}[-1]$ and projection formula, we obtain an isomorphism ${\sf D}_l\Ll (-2\omega _l)=\Ll (-\omega _{l-1}-\omega _{l+1})[-1]$. Now the exact triangle obtained by applying ${\pi _{\alpha _l}}_{\ast}$ to (\ref{eq:dualPsi_omega_l(-l)}) gives the statement.
\end{proof}

\begin{remark}\label{rem:prop_pi_k_*Psi _l(-l)_othertypes}
{The above Proposition \ref{prop:pi_k_*Psi _l(-l)a} can easily be generalized to all simple groups, but the formulae for ${\sf D}_k((\Psi _1^{\omega _k})^{\ast}\otimes \Ll (-\omega _k))$ will then depend on the fundamental weight $\omega _k$. We record the examples of ${\fb}_2$ and of ${ \bf G}_2$ that we will need later. Let $\alpha, \beta$ be the simple roots with $\alpha$ being the short root. Then if $ \bg$ is of type ${\fb}_2$, we have:

\begin{equation}\label{eq:pi_k_*Psi _l(-l)c}
%{\pi _{\alpha _k}}_{\ast}
{\sf D}_k((\Psi _1^{\omega _l})^{\ast}\otimes \Ll (-\omega _l))=\left\{
   \begin{array}{l}
    \Ll (-\omega _{\beta}), \quad k=l, \quad \alpha _k=\alpha,\\
    \\
    \Ll (-2\omega _{\alpha}), \quad k=l, \quad \alpha _k=\beta,\\
    \\
    (\Psi _1^{\omega _l})^{\ast}\otimes \Ll (-\omega _l), \quad k\neq l. \\
   \end{array}
  \right.
\end{equation}

If $ \bg$ is of type ${\bf G}_2$, we have 

\begin{equation}\label{eq:pi_k_*Psi _l(-l)d}
%{\pi _{\alpha _k}}_{\ast}
{\sf D}_k((\Psi _1^{\omega _l})^{\ast}\otimes \Ll (-\omega _l))=\left\{
   \begin{array}{l}
    \Ll (-\omega _{\beta}), \quad k=l, \quad \alpha _k=\alpha,\\
    \\
    \Ll (-3\omega _{\alpha}), \quad k=l, \quad \alpha _k=\beta,\\
    \\
    (\Psi _1^{\omega _l})^{\ast}\otimes \Ll (-\omega _l), \quad k\neq l. \\
   \end{array}
  \right.
\end{equation}
}
\end{remark}

\begin{proposition}\label{prop:pi_k_*Psi _k(-omega_k)}

\begin{equation}\label{eq:D_kPsi _lb}
{\sf D}_k(\Psi _1^{\omega _l}\otimes \Ll (-\omega _l))=\left\{
   \begin{array}{l}
    \Ll (0)[-1], \quad k=l,\\
    \\
    \Psi _1^{\omega _l}\otimes \Ll (-\omega _l), \quad k\neq l. \\
   \end{array}
  \right.
\end{equation}
\end{proposition}

\begin{proof}
Tensoring the sequence (\ref{eq:Psi_omega_i}) for $i=l$ with $\Ll (-\omega _l)$ and applying ${\sf D}_k$ to the result, and taking into account $\langle -\omega _l,\alpha _k^{\vee}\rangle =-\delta _{kl}$, the ensuing exact triangle gives the result.
\end{proof}

\begin{proposition}\label{prop:Demfunpreserveorthogonal-0}
The functors ${\sd}_w$ preserve the subcategory $\langle \Oo _{X}\rangle ^{\perp}$.
\end{proposition}

\begin{proof}
Let $\Ff$ be an object of $\langle \Oo _{X}\rangle ^{\perp}$, i.e.\ $\RHom _X(\Oo _X,\Ff)=0$. Let $w\in W$ and $s_{\alpha _i}s_{\alpha _{i_2}}\dots s_{\alpha _{i_k}}$ be a reduced expression of $w$. We need to show that $\RHom _X(\Oo _X,{\sd}_w\Ff)=0$. We argue by induction on $l(w)$. The case $l(w)=0$ holds by the assumption. Let $l(w)>0$. Setting $\Ee =\Ll (0)=\Oo _X$ in the equation (\ref{eq:Demazure_reduction}) of Proposition \ref{prop:Demazure_reduction}, we obtain an isomorphism $\RHom _X(\Oo _X,{\sd}_w\Ff)=\RHom _X(\Oo _X,{\sd}_{s_{\alpha _1}w}\Ff)$. Now $l(s_{\alpha _1}w)<l(w)$ and we conclude by induction on $l(w)$.
\end{proof}

By Theorems \ref{th:BWB_van} and \ref{th:Kodaira_van}, the collection of line bundles $\Ll(\omega _1),\dots ,\Ll(\omega _k)$ is an exceptional collection on $X$ the terms of which are completely orthogonal to each other. By Theorem \ref{th:BWB_van}, $\RHom _X(\bigoplus _{i=1}^{i=r}\Ll(\omega _i),\Oo _X)=0$ and the admissible subcategories $\langle \Oo _X\rangle$ and $\langle \Ll(\omega _1),\dots ,\Ll(\omega _k)\rangle$  form a semiorthogonal sequence in $\Dd ^b(X)$.

\begin{proposition}\label{prop:Demfunpreserveorthogonal-Pic}
The functors ${\sd}_w$ preserve the subcategory $\langle \Oo _{X},  \bigoplus _{i=1}^{i=r}\Ll(\omega _i)\rangle ^{\perp}$.
\end{proposition}

\begin{proof}
Let $\Ee$ be an object of $\langle \Oo _{X},  \bigoplus _{i=1}^{i=r}\Ll(\omega _i)\rangle ^{\perp}$, i.e.\ $\RHom _X(\Oo _{X},\Ee)=0$ and $\RHom _X(\Ll (\omega _i),\Ee)=0$ for $i=1,\dots ,r$. It is sufficient to prove the statement for a simple reflection $w=s_k$. Proposition \ref{prop:Demfunpreserveorthogonal-0} reduces the statement to verifying 
$\RHom _X(\Ll (\omega _i),{\sd}_k\Ee)=0$ for $i=1,\dots ,r$. Let first $k=i$; 
then $\RHom _X(\Ll (\omega _k),{\sd}_k\Ee)=\RHom _X(\Oo _X,{\sd}_k\Ee\otimes \Ll (-\omega _k))=0$ as ${\pi _k}_{\ast}\Ll (-\omega _k)=0$. By Proposition \ref{prop:alpha_i-triangle}, for $k\neq i$ consider the exact triangle 	

$$
\dots \rightarrow {\sd }_k(\Ee \otimes \Ll (-\omega _k))\otimes \Ll (\omega _k-\alpha_k)\rightarrow {\sd}_k\Ee\rightarrow \Ee \rightarrow \dots 
$$

By the assumption, the required orthogonality will follow from $\RHom _X(\Ll (\omega _i),{\sd }_k(\Ee \otimes \Ll (-\omega _k))\otimes \Ll (\omega _k-\alpha_k))=0$. The latter equality follows from ${\pi _k}_{\ast}\Ll (\omega _k-\alpha _k-\omega _i)=0$ as $\langle \omega _k-\alpha _k-\omega _i,\alpha _k^{\vee}\rangle =-1$. We conclude by induction on $l(w)$ as in Proposition \ref{prop:Demfunpreserveorthogonal-0}.
\end{proof}

\begin{corollary}\label{cor:Demazure_preserve_rightperpO}
The functors ${\sd}_w$ preserve the subcategory $\langle \bigoplus _{i=1}^{i=r}\Psi _1^{\omega _i}, \Oo _{X}\rangle ^{\perp}$
\end{corollary}

\begin{proof}
The categories $\langle \Oo _{X}, \bigoplus _{i=1}^{i=r}\Ll(\omega _i)\rangle$ and $\langle \bigoplus _{i=1}^{i=r}\Psi _1^{\omega _i}, \Oo _{X}\rangle$ coincide, hence so do their right orthogonals.
\end{proof}

\begin{proposition}\label{prop:purity_D_iPsi_i}
Let $w\in W$ and ${\sd }_w$ the associated Demazure functor. Then ${\sd }_w\Psi _1^{\omega _i}$ is a vector bundle for any $i=1,\dots, r$.
\end{proposition}

\begin{proof}
Applying ${\sd }_w$ to (\ref{eq:dualPsi_omega_i}), we obtain the exact triangle 
\begin{equation}\label{eq:D_wtoPsi-omega_i}
\dots \rightarrow {\sd }_w\Psi _1^{\omega _i}\rightarrow \nabla _{\omega _i}\otimes \Ll (0)\rightarrow {\sd }_w\Ll (\omega _i)\rightarrow \dots 
\end{equation}

Consider ${\sd }_w\Ll (\omega _i)$ which {\it a priori} belongs to $\Dd ^{\geq 0}(X)$, the functor ${\sd }_w$ being a right derived functor. In fact, ${\sd }_w\Ll (\omega _i)$ is a coherent sheaf, i.e.\ a pure object of $\Dd ^b(X)$: this follows from Definition \ref{def:Joseph_functor}, Proposition \ref{prop:Joseph-to-Demazure} and the higher cohomology vanishing ${\rm H}^k(X_w,\Ll (\omega _i))=0$ for $k>0$. Moreover, ${\sd }_w\Ll (\omega _i)$ is locally free, being a ${ \bg}$-equivariant coherent sheaf on $X$. 

The space of global sections 
of
$\Ll (\omega _i)$ is also isomorphic to $\nabla _{\omega _i}$ and remains to show that the morphism $\nabla _{\omega _i}\otimes \Ll (0)\rightarrow {\sd }_w\Ll (\omega _i)$ is surjective; this also follows from the surjectivity of the above restriction map on the global sections $\nabla _{\omega _i}\rightarrow {\rm H}^0(X_w,\Ll (\omega _i))$ (cf. \cite[Proposition A.2.6]{WvdKTata}) and applying Proposition \ref{prop:Joseph-to-Demazure} to that surjection. Now ${\sd }_w\Psi _1^{\omega _i}$ is the kernel of a surjective map of vector bundles, hence is also a vector bundle.
\end{proof}

\begin{proposition}\label{prop:D_iPsi_i}
The vector bundle ${\sd }_i\Psi _1^{\omega _i}$ fits into the following short exact sequence:
\begin{equation}
0\rightarrow {\sd }_i\Psi _1^{\omega _i}\rightarrow \Psi _1^{\omega _i}\rightarrow \Ll (\omega _i-\alpha _i)\rightarrow 0.
\end{equation}
\end{proposition}

\begin{proof}
Writing the exact triangle (\ref{eq:alpha_i-triangle}) from Proposition \ref{prop:alpha_i-triangle} for $\Psi _1^{\omega _i}$, we obtain 
\begin{equation}
\dots \rightarrow {\sd }_i(\Psi _1^{\omega _i}\otimes \Ll (-\omega _i))\otimes \Ll (\omega _i-\alpha _i)\rightarrow {\sd }_i\Psi _1^{\omega _i}\rightarrow \Psi _1^{\omega _i}\rightarrow \dots 
\end{equation}

Proposition \ref{prop:pi_k_*Psi _k(-omega_k)} computes ${\sd }_i(\Psi _1^{\omega _i}\otimes \Ll (-\omega _i))\otimes \Ll (\omega _i-\alpha _i)$ to be isomorphic to $\Ll (\omega _i-\alpha _i)[-1]$, so the above exact triangle gives a short exact sequence.

\begin{equation}\label{eq:D_iPsi_1^omega_i}
0\rightarrow {\sd }_i\Psi _1^{\omega _i}\rightarrow \Psi _1^{\omega _i}\rightarrow \Ll (\omega _i-\alpha _i)\rightarrow 0.
\end{equation}

The map $\Psi _1^{\omega _i}\rightarrow \Ll (\omega _i-\alpha _i)$ is given by an isomorphism $\RHom _X(\Psi _1^{\omega _i},\Ll (\omega _i-\alpha _i)
)=\RGamma (X,\Ll (-\alpha _i)[1])={\sf k}$ as can be seen from the long exact cohomology sequence associated to the tensor product of sequence (\ref{eq:dualPsi_omega_i}) with $ \Ll (\omega _i-\alpha _i)$.
\end{proof}.

\begin{proposition}\label{prop:Psi_iD_wPsi_j-orth}
Let $\omega _i,\omega _j$ be two fundamental weights and $w\in W$. Then $\RHom _X(\Psi _1^{\omega _i},{\sf D}_w\Psi _1^{\omega _j})=0$ if $l(w)>0$ and a reduced decomposition of $w$ contains $s_{\alpha _j}$ 
or $i\neq j$ and $l(w)\geq 0$.
\end{proposition}

\begin{proof}
We proceed by induction on the length $l(w)$ of $w\in W$. We assume first $i\neq j$, so the base of induction is $l(w)=0$. Then ${\sf D}_w={\sf id}$ and we need to show $\RHom _X(\Psi _1^{\omega _i},\Psi _1^{\omega _j})=0$. The pair $\Psi _1^{\omega _i},\Psi _1^{\omega _j}$ is the left mutation of the pair $\Ll (\omega _i),\Ll (\omega _j)$ through $\Ll (0)$ and the line bundles in the latter pair are mutually orthogonal, hence the bundles $\Psi _1^{\omega _i},\Psi _1^{\omega _j}$ are also mutually orthogonal.

Let $l(w)>0$. There is an isomorphism
\begin{equation}
\RHom _X(\Psi _1^{\omega _i},{\sf D}_w\Psi _1^{\omega _j})=\RGamma (X,{\sf D}_w\Psi _1^{\omega _j}\otimes (\Psi _1^{\omega _i})^{\ast})
\end{equation}

The base of induction is $l(w)=1$, thus $w=s_{\alpha _l}$ for a simple root $\alpha _l$ and by the assumption on $w$ we must have $l=j$. If $i=j$ then by projection formula we obtain $\RGamma (X,{\sf D}_i\Psi _1^{\omega _i}\otimes (\Psi _1^{\omega _i})^{\ast})=\RGamma (Y_i,{\pi _{\alpha _i}}_{\ast}\Psi _1^{\omega _i}\otimes {\pi _{\alpha _i}}_{\ast}(\Psi _1^{\omega _i})^{\ast})$
and by Proposition \ref{prop:pi_k_*Psi _l} the last group is isomorphic to 
$\RGamma (Y_1,{\pi _{\alpha _i}}_{\ast}\Psi _1^{\omega _i})\otimes (\nabla _{\omega _i})^{\ast}=\RGamma (X,\Psi _1^{\omega _i})\otimes (\nabla _{\omega _i})^{\ast}=0$. If $i\neq j$ then $\RHom _X(\Psi _1^{\omega _i},{\sf D}_j\Psi _1^{\omega _j})=\RHom _X(\Psi _1^{\omega _i},\Psi _1^{\omega _j})=0$ where the last isomorphism follows from the previous paragraph and the first isomorphism 
is explained in detail below in the equation (\ref{eq:refeq}).

Let $w = s_{\alpha _1}s_{\alpha _2}\dots s_{\alpha _k}, k=l(w)>0$ be a reduced expression of $w\in W$, so ${\sf D}_w={\sf D}_1{\sf D}_{s_{\alpha _1}w}$. Then by projection formula 

\begin{equation}
\RGamma (X,{\sf D}_w\Psi _1^{\omega _j}\otimes (\Psi _1^{\omega _i})^{\ast})=\RGamma (Y_1,{\pi _{\alpha _1}}_{\ast}{\sf D}_{s_{\alpha _1}w}\Psi _1^{\omega _j}\otimes {\pi _{\alpha _1}}_{\ast}(\Psi _1^{\omega _i})^{\ast}).
\end{equation}

Assume first $i=1$. Then using Proposition \ref{prop:pi_k_*Psi _l}, we obtain
\begin{equation}
\RGamma (Y_1,{\pi _{\alpha _1}}_{\ast}{\sf D}_{s_{\alpha _1}w}\Psi _1^{\omega _j}\otimes {\pi _{\alpha _1}}_{\ast}(\Psi _1^{\omega _1})^{\ast})=\RGamma (Y_1,{\pi _{\alpha _1}}_{\ast}{\sf D}_{s_{\alpha _1}w}\Psi _1^{\omega _j})\otimes (\nabla _{\omega _i})^{\ast}.
\end{equation}

Now $\RGamma (Y_1,{\pi _{\alpha _1}}_{\ast}{\sf D}_{s_{\alpha _1}w}\Psi _1^{\omega _j})=\RGamma (X,{\sf D}_w\Psi _1^{\omega _j})=0$ where the last isomorphism follows from Proposition \ref{prop:Demfunpreserveorthogonal-0}. 

Assume $i\neq 1$, then using again Proposition \ref{prop:pi_k_*Psi _l}, we obtain

\begin{eqnarray}\label{eq:refeq}
& \RGamma (Y_1,{\pi _{\alpha _1}}_{\ast}{\sf D}_{s_{\alpha _1}w}\Psi _1^{\omega _j}\otimes {\pi _{\alpha _1}}_{\ast}(\Psi _1^{\omega _i})^{\ast})=\RGamma (Y_1,{\pi _{\alpha _1}}_{\ast}{\sf D}_{s_{\alpha _1}w}\Psi _1^{\omega _j}\otimes (\Psi _1^{\omega _i})^{\ast})= \\ \nonumber & \RHom _X(\Psi _1^{\omega _i},{\sf D}_{s_{\alpha _1}w}\Psi _1^{\omega _j}).
\end{eqnarray}

Now $l(s_{\alpha _1}w)<l(w)$ and by the induction assumption $\RHom _X(\Psi _1^{\omega _1},{\sf D}_{s_{\alpha _1}w}\Psi _1^{\omega _j})=0$. 
\end{proof}

\begin{proposition}\label{prop:orthogon_omega-to-Demazure-Psis}
Let $\omega _i, \omega _j$ be two distinct fundamental weights of $ \bg$. Then 

$\RHom _X(\Ll (\omega _i),{\sf D}_w\Psi _1^{\omega _j})=0$

\end{proposition}

\begin{proof}
Consider the sequence (\ref{eq:Psi_omega_i}) and apply to it $\RHom _X(-,{\sf D}_w\Psi _1^{\omega _j})$. We have $\RGamma (X,{\sf D}_w\Psi _1^{\omega _j})=0$ by Proposition \ref{prop:Demfunpreserveorthogonal-0} and $\RHom _X(\Psi _1^{\omega _i},{\sf D}_w\Psi _1^{\omega _j})=0$ by the previous Proposition \ref{prop:Psi_iD_wPsi_j-orth}. Hence, $\RHom _X(\Ll (\omega _i),{\sf D}_w\Psi _1^{\omega _j})=0$.
\end{proof}

%----------------------------------------------------------------------------
%\subsection{ Semiorthogonal sequences via Demazure functors}
%----------------------------------------------------------------------------

\begin{proposition}\label{prop:exceptionality_D_i}
The rank two vector bundles ${\sd}_i\Ll (\omega _i), i=1,\dots ,r$ are  exceptional, except for the two exceptional cases (see the footnote).
\end{proposition}

\begin{proof}
Applying $\RHom _X(-,{\sd}_i\Ll (\omega _i))$ to the short exact sequence (\ref{eq:Demazureses}), and using Proposition \ref{prop:Demfunpreserveorthogonal-Pic}, we obtain an isomorphism $\RHom _X({\sd}_i(\Ll (\omega _i),{\sd}_i(\Ll (\omega _i))=\RHom _X(\Ll (\omega _i-\alpha _i),{\sd}_i(\Ll (\omega _i))$. In turn, applying 
$\RHom _X(\Ll (\omega _i-\alpha _i),-)$ to the same sequence (\ref{eq:Demazureses}), and observing that $\RHom _X(\Ll (\omega _i-\alpha _i),\Ll (\omega _i))=\RGamma (X,\Ll (\alpha _i))=0$, since $\langle \alpha _i,\alpha _j^{\vee}\rangle =-1$ for at least one $j\in [1,\dots ,r],j\neq i$ (cf. Cartan matrix), obtain $\RHom _X({\sd}_i(\Ll (\omega _i),{\sd}_i(\Ll (\omega _i))=k.\footnote{$\rm NB$: This breaks down in the ${\bf G}_2$ case (cf. the rank 2 vector bundle on ${\bf G}_2/{\bp}_{\beta}={\sf Q}_5$ obtained by the tautological representation of the Levi factor ${\bp}_{\beta}$) and for the root $\alpha _n$ in type ${\fb}_n$.}$
\end{proof}

\begin{corollary}\label{cor:Psi_omega-exc}
In the assumption of the previous proposition, the bundles ${\sd }_i\Psi _1^{\omega _i}, i=1,\dots ,r$ are exceptional. 
\end{corollary}

\begin{proof}
Recall that $\Psi _1^{\omega _i}, i=1,\dots ,r$ are defined to be the left mutations of $\Ll (\omega _i)$ through $\Oo _X$. Hence ${\sd }_i\Psi _1^{\omega _i}$ is the cokernel of the evaluation morphism $\nabla _{\omega _i}\otimes \Oo _X\rightarrow {\sd}_i\Ll (\omega _i)$. By the previous Proposition \ref{prop:exceptionality_D_i}, the bundles ${\sd}_i\Ll (\omega _i), i=1,\dots ,r$ are exceptional, hence so are ${\sd }_i\Psi _1^{\omega _i}$.
\end{proof}

\begin{proposition}\label{prop:Demazure_fullness}
Let the rank of group $ \bg$ be equal to 2, and let ${\bp}_{\alpha _i}$ denote the minimal parabolic generated by ${\fb}$ and $\bu_{\alpha _i}$ for a simple root $\alpha _i, i=1,2$. Consider the subcategory $\pi _{\alpha _1}^{\ast}\Dd ^b(Y_1)\diamond \pi _{\alpha _2}^{\ast}\Dd ^b(Y_2)$, i.e.\ the minimal thick strict full triangulated subcategory containing both subcategories $\pi _{\alpha _i}^{\ast}\Dd ^b(Y_i)$ for $i=1,2$. Then 
\begin{equation}\label{eq:leftperprho}
\langle \pi _{\alpha _1}^{\ast}\Dd ^b(Y_1)\diamond \pi _{\alpha _2}^{\ast}\Dd ^b(Y_2)\rangle={^{\perp}\langle \Ll (-\rho)\rangle}.
\end{equation}

\end{proposition}

\begin{proof}
Note that the subcategory $\langle \Ll (-\rho)\rangle$ is admissible by Theorem \ref{th:Kodaira_van} and Section \ref{subsec:FEC}.
We will prove an equivalent statement: 

\begin{equation}\label{eq:rightperp0}
^{\perp}\langle (\pi _{\alpha _1}^{\ast}\Dd ^b(Y_1))\otimes \Ll (-\rho)\diamond (\pi _{\alpha _2}^{\ast}\Dd ^b(Y_2))\otimes\Ll (-\rho)\rangle=\langle \Ll (0)\rangle.
\end{equation}

Indeed, having established the above equality, we obtain (\ref{eq:leftperprho}) by tensoring (\ref{eq:rightperp0}) with $\Ll (\rho)$ and mutating the subcategory $\langle \Ll (\rho)\rangle$ to the left past its right orthogonal. By \cite[Lemma 2.11]{Kuzrat}, the mutation is given by the Serre functor which is isomorphic to tensoring with $\omega _X=\Ll (-2\rho)$. 

To prove (\ref{eq:rightperp0}), observe that $^{\perp}\langle (\pi _{\alpha _1}^{\ast}\Dd ^b(Y_1))\otimes \Ll (-\rho)\diamond (\pi _{\alpha _2}^{\ast}\Dd ^b(Y_2))\otimes\Ll (-\rho)\rangle\subsetq \bigcap _{\alpha _i}\pi _{\alpha}^{\ast}\Dd ^b(Y_i)$. Indeed, by Theorem \ref{th:Orlovth} any object $\mathcal \bg$ of $^{\perp}\langle (\pi _{\alpha _1}^{\ast}\Dd ^b(Y_1))\otimes \Ll (-\rho)\diamond (\pi _{\alpha _2}^{\ast}\Dd ^b(Y_2))\otimes\Ll (-\rho)\rangle$ belongs to $\pi _{\alpha _i}^{\ast}\Dd ^b(Y_i)$ for each simple root $\alpha _i$ and by Serre duality

\begin{equation}
\RHom  _X(\mathcal G,\pi _{\alpha _i}^{\ast}{\pi _{\alpha _i}}_{\ast}\Ff \otimes \Ll ({-\rho})) = \RHom  _{Y_i}({\pi _{\alpha _i}}_{\ast}\Ff,{\pi _{\alpha _i}}_{\ast} (\mathcal G\otimes \Ll ({-\rho}))[{\rm dim}(X)])^{\ast}=0.
\end{equation}

Since $\Ff$ is arbitrary and the functor ${\pi _{\alpha _i}}_{\ast}$ is surjective, it follows that ${\pi _{\alpha _i}}_{\ast} (\mathcal G\otimes \Ll ({-\rho}))=0$. This implies $\mathcal G\in \pi _{\alpha _i}^{\ast}\Dd ^b(Y_i)$, the line bundle $\Ll ({-\rho})$ having degree $-1$ along $\pi _{\alpha}$. It follows that $\sd _{\alpha _i}\mathcal G=\mathcal \G$ for $i=1,2$ for such an object $\mathcal \G$. Choosing a reduced expression for the longest element $w_0$, we obtain $\sd _{w_0}\mathcal G=\mathcal G$, and 
hence $\pi ^{\ast}\pi _{\ast}\mathcal G=\mathcal G$, since $\sd _{w_0}=\pi ^{\ast}\pi _{\ast}$ by Theorem \ref{th:Demw0}. Now $\pi ^{\ast}\pi _{\ast}\mathcal G=\Ll (0)\otimes \pi _{\ast}\mathcal G=\Ll (0)\otimes \RGamma(X,\mathcal G)\in \langle \Ll (0)\rangle$, and the statement follows.
\end{proof}

\begin{corollary}\label{cor:Demazurefullness}
Let the rank of group $ \bg$ be equal to 2, and $X$ the flag variety of $ \bg$. Then $\Dd ^b(X)$ has a semiorthogonal decomposition:

\begin{equation}
\Dd ^b(X)=\langle \langle \Ll (-\rho)\rangle, \langle \pi _{\alpha _1}^{\ast}\Dd ^b(Y_1)\diamond \pi _{\alpha _2}^{\ast}\Dd ^b(Y_2)\rangle\rangle .
\end{equation}

\end{corollary}

\begin{remark}\label{rem:-rhoSOD}
{Proposition \ref{prop:Demazure_fullness} and Corollary \ref{cor:Demazurefullness} hold more generally for groups of arbitrary rank.}
\end{remark}

%--------------------------------------------------------------
\subsection{Baby case: rank one}
%--------------------------------------------------------------

In rank one we are dealing with $\bg=\mathrm{SL}_2$, $W=\{s,\id\}$.
Let $\rho$ be the fundamental weight.
One has $e_s=-\rho$, $e_\id=0$, $P(0)=\sk$, $Q(-\rho)=\sk_{-\rho}$, ${\sf Q}_{\succeq \id}={\sf P}_{\preceq s}=
\Dd ^b(\rep(\fb))$, $X_s=Y_s=Q(-\rho)=\sk_{-\rho}$, $X_\id=P(0)^*=\sk$, $\tX_s=\Ll_{-\rho}=\Oo(-1)$, $\tX_\id=\Oo$,
where $\Oo$ is the structure sheaf of the projective line
$\bg/\fb$.
One gets the familiar full
exceptional collection $\Oo(-1)$, $\Oo$ on the projective line.

%--------------------------------------------------------------
\subsection{ Type ${\bf A}_2$}\label{subsec:typeA_2}
%--------------------------------------------------------------

Let $\omega _1$ and $\omega _2$ be the two fundamental weights for ${\bf SL}_3$ and $\alpha _1,\alpha _2$ the simple roots. As before, set $X:={\bf SL}_3/{\fb}$ and $Y_i:={\bf SL}_3/{\bp}_{\alpha _i}, i=1,2$.

\begin{lemma}\label{lem:A_2-Demazure-FEC}

The following block collection is a full strong exceptional collection in $\Dd ^b(X)$:

%\vspace{0.3cm}

%\begin{figure}[\bh]
\begin{equation}\label{eq:FrobExcCollonSl_3/B}
\xymatrix{
& \A _{-3}& \A _{-2}& \A _{-1}& \A _0\\
& || & || & || & || \\
&*++<10pt>[F]\txt{$\Ll (-\rho)$}
&*++<10pt>[F]\txt{${\sf D}_{\alpha _1}(\Psi _1^{\omega _1})$  \\ \\ ${\sf D}_{\alpha _2}(\Psi _1^{\omega _2})$}
&*++<10pt>[F]\txt{$\Psi _1^{\omega _1}$  \\ \\ $\Psi _1^{\omega _2}  $}
 &*++<10pt>[F]\txt{$\Ll (0)$}
}
\end{equation}
% \end{figure}
\end{lemma}

That is, $\A _{i}\subset {^{\perp}\A _{j}}$ for $-3\leq j<i\leq 0$,
each of the vector bundles is exceptional, the vector bundles in a box are completely orthogonal, and the triangulated hull of $\langle 
{\sf A}_{-3},\dots,{\sf A}_0\rangle$ is the whole $\Dd ^b(X)$. 

\begin{lemma}\label{lem:dualA_2-Demazure-FEC}
Consider the following collection:

%\begin{figure}[\bh]
\begin{equation}\label{eq:dualFrobExcCollonSl_3/B}
\xymatrix{
& \B _{0}& \B _{1}& \B _{2}& \B _3\\
& || & || & || & || \\
&*++<10pt>[F]\txt{$\Ll (0)$}
&*++<10pt>[F]\txt{${\sf D}_{\alpha _1}(\Psi _1^{\omega _1})\otimes \Ll (\rho)$  \\ \\ ${\sf D}_{\alpha _2}(\Psi _1^{\omega _2})\otimes \Ll (\rho)$}
&*++<10pt>[F]\txt{$\Psi _1^{\omega _1}\otimes \Ll (\rho)$  \\ \\ $\Psi _1^{\omega _2}\otimes \Ll (\rho)$}
 &*++<10pt>[F]\txt{$\Ll (\rho)$}
}
\end{equation}

\noindent
% \end{figure}
in which the bundles from the block ${\B _i}, i=0,\dots,3$ are shifted by $[-i]$ in the category $\Dd ^b(X)$. In other words, the above collection is obtained from the collection (\ref{eq:FrobExcCollonSl_3/B}) by tensoring with $\Ll (\rho)$ and applying the shift functor relative to the corresponding block. Then the collection (\ref{eq:dualFrobExcCollonSl_3/B}) is left dual to (\ref{eq:FrobExcCollonSl_3/B}).
\end{lemma}

\begin{remark}\label{rem:SteinbergA_2}
{The Steinberg weights for ${\bf A}_2$ are given by (see Section \ref{subsec:Steinberg})
$$\ -\rho, \ -\omega _1, \ -\omega _2, \ -\omega _1+\omega _2, \ \omega _1-\omega _2, \ 0;$$
these coincide with the weights of the decomposition factors of the collection (\ref{eq:FrobExcCollonSl_3/B}).
}
\end{remark}

%--------------------------------------------------------------
\subsection{ The quiver for ${\bf A}_2$}
%--------------------------------------------------------------

The picture below depicts the path quiver algebra of the above full exceptional collection. There are only zero-degree $\Hom$-spaces in the non-trivial direction in the collection (\ref{eq:FrobExcCollonSl_3/B}). The underlying graph of the quiver is the Bruhat graph for $SL_3/{\fb}$; the quiver path algebra is generated in degree one and the relations are quadratic that one reads off the $\Hom$--spaces next to the arrows of the quiver.
The labels assume $p$ is not too small. For small $p$ the characters are still correct and the $\Hom$--spaces have a good filtration.
%\vspace{0.3cm}

\begin{equation}
\xymatrix{
& {\sd }_{\alpha _1}(\Psi _1^{\omega _1})\ar@/^6pc/[rrrrd]^{\nabla _{\omega _1}}\ar[rrrdd]\ar[rrr]^{\nabla _{\omega _2}} & & & \Psi _1^{\omega _1}\ar[rd]_{\nabla _{\omega _2}}  \\
\Ll (-\rho) \ar@/^6pc/[rrrru]^{{\sf Ker}(\nabla _{\omega _1}\otimes \nabla _{\rho}\rightarrow \nabla _{\omega _1+\rho})
}\ar@/_6pc/[rrrrd]_{{\sf Ker}(\nabla _{\omega _2}\otimes \nabla _{\rho}\rightarrow \nabla _{\omega _2+\rho})}\ar[rd]^{\nabla _{\omega _1}}\ar[ru]_{\nabla _{\omega _2}}  \quad  &   & &  &  &  \quad \Ll (0) \\
& {\sd }_{\alpha _2}(\Psi _1^{\omega _2})\ar@/_6pc/[rrrru]_{\nabla _{\omega _2}}\ar[rrruu]\ar[rrr]_{\nabla _{\omega _1}}  &   & & \Psi _1^{\omega _2}\ar[ru]^{\nabla _{\omega _1}} \\
}
\end{equation}\\
The Weyl group elements that we associate to the nodes are
as follows: The leftmost node belongs with the longest element $w_0$, the rightmost with the identity element. In the top row we have from left to right 
$s_2 s_1$, $s_1$, 
and in the bottom row from left to right
$s_1 s_2$, 
$s_2 $. 

We list some of the arrows:

\begin{itemize}

\item[(1)] $\RHom (\Ll(-\rho),\Psi _1^{\omega _1})={\sf Ker}(\nabla _{\omega _1}\otimes \nabla _{\rho}\rightarrow \nabla _{\omega _1+\rho})$. Note that $\nabla$-filtration on both ${\sf Ker}(\nabla _{\omega _1}\otimes \nabla _{\rho}\rightarrow \nabla _{\omega _1+\rho})$ and on $\nabla _{\omega _2}\otimes \nabla _{\omega _2}$ coincide and the associated graded is isomorphic to $\nabla _{\omega _1}\oplus \nabla _{2\omega _2}$.

\item[(2)] $\RHom (\Ll(-\rho),\Psi _1^{\omega _2})={\sf Ker}(\nabla _{\omega _2}\otimes \nabla _{\rho}\rightarrow \nabla _{\omega _2+\rho})$. Note that $\nabla$-filtration on both ${\sf Ker}(\nabla _{\omega _2}\otimes \nabla _{\rho}\rightarrow \nabla _{\omega _2+\rho})$ and on $\nabla _{\omega _1}\otimes \nabla _{\omega _1}$ coincide and the associated graded is isomorphic to $\nabla _{\omega _2}\oplus \nabla _{2\omega _1}$.

\item[(3)] $\RHom (\Ll(-\omega _2),\Psi _1^{\omega _1})={\sf k}$

\item[(4)] $\RHom (\Ll(-\omega _1),\Psi _1^{\omega _2})={\sf k}$

\end{itemize}

The $\nabla$-filtrations on some tensor products are given by:

\begin{itemize}

\item[(1)] $\nabla _{\omega _1}\otimes \nabla _{\omega _1}$ is filtered by $\nabla _{2\omega _1}$ and $\nabla _{\omega _2}$;

\item[(2)] $\nabla _{\omega _2}\otimes \nabla _{\omega _2}$ is filtered by $ \nabla _{\omega _1}$ and $\nabla _{2\omega _2}$

\item[(3)] $\nabla _{\omega _1}\otimes \nabla _{2\omega _1}$ is filtered by 
$\nabla _{3\omega _1}$ and $\nabla _{\rho}$;

\item[(4)] $\nabla _{\omega _2}\otimes \nabla _{2\omega _2}$ is filtered by
$\nabla _{3\omega _2}$ and $\nabla _{\rho}$;

\item[(5)] $\nabla _{\omega _1}\otimes \nabla _{\omega _2}$ is filtered by $\nabla _{\rho}$ and ${\sf k}$;

\item[(6)] $\nabla _{\omega _1}\otimes \nabla _{\rho}$ is filtered by $\nabla _{\omega _1+\rho},\nabla _{\omega _1}$ and $\nabla _{2\omega _2}$.

\end{itemize}

%--------------------------------------------------------------
\subsection{ Type ${\bf B}_2$}\label{subsec:typeB_2}
%--------------------------------------------------------------
We assume here that the prime $2$ is invertible in $\sk$. 
Let $ \bg$ be a simple group of type ${\fb}_2$. Let $\alpha$ and $\beta$ be the simple roots for ${\fb}_2$ with $\beta$ being the long root. Let $\omega _{\alpha}$ and $\omega _{\beta}$ be the two fundamental weights. Let ${\bp}_{\alpha}$ (resp., ${\bp}_{\beta}$) be the minimal parabolic generated by ${\fb}$ and $\bu_{\alpha}$ (resp., ${\fb}$ and $\bu_{\beta}$). The Grassmannian ${ \bg}/{\bp}_{\alpha}$ is embedded into $  \Pp (\nabla _{\omega _{\beta}}^{\ast})$ with the help of line bundle $\Ll (\omega _{\beta})$ and ${\rm dim}\nabla _{\omega _{\beta}}=5$. The Grassmannian ${ \bg}/{\bp}_{\beta}$ is embedded into $  \Pp (\nabla _{\omega _{\alpha}}^{\ast})$ with the help of line bundle $\Ll (\omega _{\alpha})$ and ${\rm dim}\nabla _{\omega _{\alpha}}=4$ (it is in fact isomorphic to $  \Pp ^3=  \Pp (\nabla _{\omega _{\alpha}}^{\ast}))$.
Associated to the fundamental weights $\omega _{\alpha}$ and $\omega _{\beta}$ are the vector bundles $\Psi _1^{\omega _{\alpha}}$ and $\Psi _1^{\omega _{\beta}}$ defined as in the sequence (\ref{eq:Psi_omega_i}).

\begin{lemma}\label{lem:B_2-Demazure-FEC}

The following collection is a full exceptional collection in $\Dd ^b( \bg/{\fb})$:

%\begin{figure}[\bh]
\begin{equation}\label{eq:FrobExcCollonSp_4/B}
\xymatrix{
\A _{-4} & \A _{-3}& \A _{-2}& \A _{-1}& \A _0\\
|| & || & || & || & || \\
*++<10pt>[F]\txt{$\Ll (-\rho)$}
&*++<10pt>[F]\txt{${\sf D}_{\beta\alpha}(\Psi _1^{\omega _{\alpha}})$  \\ \\ ${\sf D}_{\alpha\beta}(\Psi _1^{\omega _{\beta}})$}
&*++<10pt>[F]\txt{${\sf D}_{\alpha}(\Psi _1^{\omega _{\alpha}})$  \\ \\ ${\sf D}_{\beta}(\Psi _1^{\omega _{\beta}})$}
&*++<10pt>[F]\txt{$\Psi _1^{\omega _{\alpha}}$  \\ \\ $\Psi _1^{\omega _{\beta}}  $}
 &*++<10pt>[F]\txt{$\Ll (0)$}
}
\end{equation}
 %\end{figure}
That is, $\A _{i}\subset {^{\perp}\A _{j}}$ for $-4\leq j<i\leq 0$,
each of the vector bundles is exceptional, the vector bundles in a box are completely orthogonal, and the triangulated hull of $\langle 
{\sf A}_{-4},\dots,{\sf A}_0\rangle$ is the whole $\Dd ^b(X)$. 
 
\end{lemma}

\begin{remark}\label{rem:identB_2}
{
The $ \bg$-equivariant vector bundles in the above collection can be explicitly identified: 

\begin{itemize}[(1)]

\item[(1)] ${\sf D}_{\alpha}(\Psi _1^{\omega _{\alpha}}) = \Uu _2$

\item[(2)] ${\sf D}_{\beta}(\Psi _1^{\omega _{\beta}}) = \Psi _2^{\omega _{\alpha}}$

\item[(3)] ${\sf D}_{\beta\alpha}(\Psi _1^{\omega _{\alpha}}) = \Ll (-\omega _{\alpha})$

\item[(4)] ${\sf D}_{\alpha\beta}(\Psi _1^{\omega _{\beta}}) = \Ll (-\omega _{\beta})$.

\end{itemize}
Here the rank two bundle $\Uu _2$ fits into a short exact sequence 

\begin{equation}\label{eq:weights_of_U2}
0\rightarrow \Ll (-\omega _{\alpha})\rightarrow \Uu _2\rightarrow \Ll (\omega _{\alpha}-\omega _{\beta})\rightarrow 0;
\end{equation}
the corresponding extension is given by the isomorphism $\Ext ^1_{ \bg/\fb}(\Ll (\omega _{\alpha}-\omega _{\beta}),\Ll (-\omega _{\alpha}))={\rm H}^1( \bg/ \fb,\Ll (-\beta))=k$. The bundle $\Psi _2^{\omega _{\alpha}}$ is the left mutation of line bundle $\Ll (2\omega _{\alpha})$ through the admissible subcategory $\langle \Ll (0),\Ll (\omega _{\alpha})\rangle\subset \Dd ^b( \bg/ \fb)$ and fits into a short exact sequence 
\begin{equation}
0\rightarrow \Psi _2^{\omega _{\alpha}}\rightarrow \Lambda ^2\nabla _{\omega _{\alpha}}\otimes \Ll (0)\rightarrow \Psi _1^{\omega _{\alpha}}\otimes \Ll (\omega _{\alpha})\rightarrow 0.
\end{equation}
The isomorphism (1) follows from Proposition \ref{prop:D_iPsi_i}. 
Let us show how (2) is obtained as Proposition \ref{prop:exceptionality_D_i} doesn't directly lead to conclusion that the bundle ${\sf D}_{\beta}(\Psi _1^{\omega _{\beta}})$ is exceptional in the non simply-laced case which is ${\bf B}_2$. To this end, observe that the collection of bundles 
\begin{equation}
\Ll (-\omega _{\alpha}),\Psi _2^{\omega _{\alpha}},\Psi _1^{\omega _{\alpha}},\Ll (0);
\end{equation}
is a full exceptional collection on $ \bg/{\bp} _{\beta}$. To ensure that ${\sf D}_{\beta}\Psi _1^{\omega _{\beta}}=\Psi _2^{\omega _{\alpha}}$ it is sufficient to check that $\RHom (\Ll (0),{\sf D}_{\beta}\Psi _1^{\omega _{\beta}})=0$, $\RHom (\Psi _1^{\omega _{\alpha}},{\sf D}_{\beta}\Psi _1^{\omega _{\beta}})=0$ and $\RHom ({\sf D}_{\beta}\Psi _1^{\omega _{\beta}},\Ll (-\omega _{\alpha}))=0$, and the ranks of bundles in Remark \ref{rem:identB_2}, (2) are the same. That the ranks are the same is immediate from Proposition \ref{prop:purity_D_iPsi_i}. As for the orthogonality, 
the first equality holds by Proposition \ref{prop:Demfunpreserveorthogonal-0}, the second one by Proposition \ref{prop:Psi_iD_wPsi_j-orth}, and the third one is a consequence of $\RHom _X({\sf D}_{\beta}\Psi _1^{\omega _{\beta}},\Ll (-\omega _{\alpha}))=0$. The latter equality is obtained by applying $\Hom _X(-,\Ll (-\omega _{\alpha}))$ to the sequence (\ref{eq:D_wtoPsi-omega_i}) and then using Proposition \ref{prop:Demfunpreserveorthogonal-0}, the exact sequence (\ref{eq:Demazureses}), and Theorem \ref{th:BWB_van}: all weights that are needed to compute $\RHom _X({\sf D}_{\beta}(\Ll (\omega _{\beta})),\Ll (-\omega _{\alpha}))$ are singular.}
\end{remark}

\begin{lemma}\label{lem:dualB_2-Demazure-FEC}

Consider the following collection:
\begin{figure}[H]

\begin{equation}\label{eq:dualFrobExcCollonSp_4/B}
\scalebox{.8}{
\xymatrix{
\B _0 & \B _1& \B _2& \B _3& \B _4\\
|| & || & || & || & || \\
*++<10pt>[F]\txt{$\Ll (0)$}
&*++<10pt>[F]\txt{${\sf D}_{\beta\alpha}(\Psi _1^{\omega _{\alpha}})\otimes \Ll (\rho)$  \\ \\ ${\sf D}_{\alpha\beta}(\Psi _1^{\omega _{\beta}})\otimes \Ll (\rho)$}
&*++<10pt>[F]\txt{${\sf D}_{\alpha}(\Psi _1^{\omega _{\alpha}})\otimes \Ll (\rho)$  \\ \\ ${\sf D}_{\beta}(\Psi _1^{\omega _{\beta}})\otimes \Ll (\rho)$}
&*++<10pt>[F]\txt{$\Psi _1^{\omega _{\alpha}}\otimes \Ll (\rho)$  \\ \\ $\Psi _1^{\omega _{\beta}}\otimes \Ll (\rho)  $}
 &*++<10pt>[F]\txt{$\Ll (\rho)$}
}}
\end{equation}

 \end{figure}
 \noindent
in which the bundles from the block ${\B _i}, i=0,\dots,4$ are shifted by $[-i]$ in the category $\Dd ^b({\bf Sp}_4/{\fb})$. In other words, the above collection is obtained from the collection (\ref{eq:FrobExcCollonSp_4/B}) by tensoring with $\Ll (\rho)$ and applying the shift functor relative to the corresponding block. Then the collection (\ref{eq:dualFrobExcCollonSp_4/B}) is right dual to (\ref{eq:FrobExcCollonSp_4/B}). 
 
\end{lemma}

\begin{remark}\label{rem:SteinbergB_2}
{The Steinberg weights for ${\bf B}_2$ are given by (see Section \ref{subsec:Steinberg})

$$\ -\rho, \ -\omega _{\alpha}, \ -\omega _{\beta}, \ -2\omega _{\alpha}+\omega _{\beta}, \ \omega _{\alpha}-\omega _{\beta}, \ -\omega _{\alpha}+\omega _{\beta}, \ 2\omega _{\alpha}-\omega _{\beta}, \ 0;$$

these coincide with the weights of the decomposition factors of the collection (\ref{eq:FrobExcCollonSp_4/B}).
}
\end{remark}

%--------------------------------------------------------------
\subsection{ The quiver for ${\bf B}_2$}
%--------------------------------------------------------------

%\vspace*{0.5cm}

The picture below depicts the path quiver algebra of the above full exceptional collection. There are only zero-degree $\Hom$-spaces in the non-trivial direction in the collection (\ref{eq:FrobExcCollonSp_4/B}). The underlying graph of the quiver is the Bruhat graph for $Sp_4/{\fb}$; the quiver path algebra is generated in degree one and the relations are quadratic that one reads off the $\Hom$--spaces next to the arrows of the quiver. 
This diagram is over a field of characteristic at least seven or characteristic zero.  The $\RHom$'s have a good filtration  for $p\geq3$. For $p=3,\;5$ we checked that the characters are still correct (cf. \cite[Section 2.6, (7)]{Donkin2}).
 

\begin{equation}\label{eq:quiverFrobExcCollonSp_4/B}
\scalebox{.9}{\xymatrix{
& {\sd}_{\beta\alpha}(\Psi _1^{\oa})\ar[rrdd]\ar[rr]^{\nabla _{\oa}}\ar@/^6pc/[rrrr]^{\Lambda ^2\nabla _{\oa}=\nabla _{\ob}\oplus {\sf k}}  & & {\sd}_{\alpha}(\Psi _1^{\oa})\ar[rr]^{\nabla _{\oa}} \ar[rrdd]\ar@/^6pc/[rrrd]^{\Lambda ^2\nabla _{\oa}=\nabla _{\ob}\oplus {\sf k}} & & \Psi _1^{\oa}\ar[rd]_{\nabla _{\oa}} \\
\Ll (-\rho) \ar@/^6pc/[rrru]^{\nabla _{\oa}\otimes \nabla _{\ob}= \nabla _{\rho}\oplus \nabla _{\oa}}\ar@/_6pc/[rrrd]_{\nabla _{2\oa}\oplus \nabla _{\ob}}\ar[rd]^{\nabla _{\oa}}\ar[ru]_{\nabla _{\ob}} \quad  &   & &  &  & &  \quad \Ll (0)\\
& {\sd}_{\alpha\beta}(\Psi _1^{\ob})\ar[rr]_{\nabla _{\oa}}\ar[rruu]\ar@/_6pc/[rrrr]_{\nabla _{2\oa}\oplus {\sf k}} & &  {\sd}_{\beta}(\Psi _1^{\ob})\ar[rr]_{\nabla _{\oa}}\ar[rruu]\ar@/_6pc/[urrr]_{\nabla _{\oa}}  & & \Psi _1^{\ob}\ar[ru]^{\nabla _{\ob}} \\
}
}
\end{equation}
The Weyl group elements that we associate to the nodes are
as follows: The leftmost node belongs with the longest element $w_0$, the rightmost with the identity element. In the top row we have from left to right 
$s_\alpha s_\beta s_\alpha $, 
$s_\beta  s_\alpha   $, 
$s_\alpha  $, 
and in the bottom row from left to right
$s_\beta s_\alpha s_\beta$, 
$s_\alpha s_\beta $, 
$s_\beta$. 

\

The $\nabla$-filtrations on the tensor products of fundamental representations are given by:

 \begin{enumerate}
\item $\nabla _{\omega _{\alpha}}\otimes \nabla _{\omega _{\alpha}}$ is filtered by $\nabla _{2\omega _{\alpha}}, \nabla _{\omega _{\beta}}$, and $\sk$.
\item  $\nabla _{\omega _{\alpha}}\otimes \nabla _{\omega _{\beta}}$ is filtered by $\nabla _{\rho}$ and $\nabla _{\omega _{\alpha}}$.
\end{enumerate}


%--------------------------------------------------------------
\subsection{ Type ${\bf G}_2$}\label{subsec:typeG_2}
%--------------------------------------------------------------

Let $p\geq 5$ be a prime. Let $ \bg$ be a simple group of type ${ \bg}_2$. Let $\alpha$ and $\beta$ be the simple roots for ${ \bg}_2$ with $\beta$ being the long root. Let $\omega _{\alpha}$ and $\omega _{\beta}$ be the two fundamental weights. Let ${\bp}_{\alpha}$ (resp., ${\bp}_{\beta}$) be the minimal parabolic generated by ${\fb}$ and ${ \bp}_{\alpha}$ (resp., ${\fb}$ and ${ \bp}_{\beta}$). The Grassmannian ${ \bg}/{\bp}_{\beta}$ is embedded into $  \Pp (\nabla _{\omega _{\alpha}}^{\ast})$ via the line bundle $\Ll (\omega _{\alpha})$ and ${\rm dim}\nabla _{\omega _{\alpha}}=7$. The Grassmannian ${ \bg}/{\bp}_{\alpha}$ is embedded into $  \Pp (\nabla _{\omega _{\beta}}^{\ast})$ via the line bundle $\Ll (\omega _{\beta})$ and ${\rm dim}\nabla _{\omega _{\beta}}=14$. As usual, denote $\pi _{\alpha}:{ \bg}/{\fb}\rightarrow { \bg}/{\bp}_{\alpha}$ (resp., $\pi _{\beta}:{ \bg}/{\fb}\rightarrow { \bg}/{\bp}_{\beta}$) the projection. Associated to $\pi _{\alpha}$ and $\pi _{\beta}$ are Demazure operators ${\sf D}_{\alpha}:=\pi _{\alpha}^{\ast}{\pi _{\alpha}}_{\ast}$ and ${\sf D}_{\beta}:=\pi _{\beta}^{\ast}{\pi _{\beta}}_{\ast}$. Associated to the fundamental weights $\omega _{\alpha}$ and $\omega _{\beta}$ are the vector bundles $\Psi _1^{\omega _{\alpha}}$ and $\Psi _1^{\omega _{\beta}}$ defined as in the sequence (\ref{eq:Psi_omega_i}).

\begin{lemma}\label{lem:G_2-Demazure-FEC}

The following collection is a full exceptional collection in $\Dd ^b({ \bg}_2/{\fb})$:



\qquad\qquad\scalebox{.67}{
$$
\centerline{
\xymatrix{
 \A _{-6} & \A _{-5}& \A _{-4} & \A _{-3}& \A _{-2}& \A _{-1}& \A _0\\
 || & || & || & || & || & || & || \\
*++<1pt>[F]\txt{$\Ll (-\rho)$}
&*++<10pt>[F]\txt{${\sf D}_{\beta\alpha\beta\alpha}(\Psi _1^{\omega _{\alpha}})$  \\ \\ ${\sf D}_{\alpha}((\Psi _1^{\omega _{\alpha}})^{\ast}(-\omega _{\alpha}))$}
&*++<10pt>[F]\txt{$(\Psi _1^{\omega _{\alpha}})^{\ast}(-\omega _{\alpha})$  \\ \\ ${\sf D}_{\alpha\beta\alpha}(\Psi _1^{\omega _{\alpha}})$}
&*++<10pt>[F]\txt{${\sf D}_{\beta\alpha}(\Psi _1^{\omega _{\alpha}})$  \\ \\ ${\sf D}_{\alpha}(\Psi _2^{\omega _{\alpha}})$}
&*++<10pt>[F]\txt{$\Psi _2^{\omega _{\alpha}}$  \\ \\ ${\sf D}_{\alpha}(\Psi _1^{\omega _{\alpha}})$}
&*++<10pt>[F]\txt{$\Psi _1^{\omega _{\alpha}}$ \\ \\ $\Psi _1^{\omega _{\beta}}  $}
 &*++<1pt>[F]\txt{$\Ll (0)$}
}
}
$$
}


%\end{equation}
% \end{figure}

\

\

That is, $\A _{i}\subset {^{\perp}\A _{j}}$ for $-6\leq j<i\leq 0$,
each of the vector bundles is exceptional, the vector bundles in a box are completely orthogonal, and the triangulated hull of $\langle 
{\sf A}_{-6},\dots,{\sf A}_0\rangle$ is the whole $\Dd ^b(X)$. 
\end{lemma}

Here $\Psi _2^{\omega _{\alpha}}$ is the left mutation of $\Ll (2\omega _{\alpha})$ through the admissible subcategory $\langle \Oo ,\Ll (\omega _{\alpha})\rangle$ and $=(\Psi _1^{\omega _{\alpha}})^{\ast}\otimes \Ll (-\omega _{\alpha})$ is the right mutation of $\Ll (-2\omega _{\alpha})$ through $\Ll (2\omega _{\alpha})$. Specifically, $\Psi _2^{\omega _{\alpha}}$ is defined by the following exact sequence:
\begin{equation}\label{eq:star}
0\rightarrow \Psi _2^{\omega _{\alpha}}\rightarrow {\sf Ker}(\nabla _{\omega _{\alpha}}^{\otimes 2}\rightarrow \nabla _{2\omega _{\alpha}})\otimes \Oo _{{ \bg}_2/{\fb}}\rightarrow \nabla _{\omega _{\alpha}}\otimes \Ll (\omega _{\alpha})\rightarrow \Ll (2\omega _{\alpha})\rightarrow 0.
\end{equation}

The bundle $(\Psi _1^{\omega _{\alpha}})^{\ast}\otimes \Ll (-\omega _{\alpha})$ is defined by the following short exact sequence:

\begin{equation}\label{eq:starstar}
0\rightarrow \Ll (-2\omega _{\alpha})\rightarrow \nabla _{\omega _{\alpha}}\otimes \Ll (-\omega _{\alpha})\rightarrow (\Psi _1^{\omega _{\alpha}})^{\ast}\otimes \Ll (-\omega _{\alpha})\rightarrow 0.
\end{equation}

The above sequence is obtained by dualizing (\ref{eq:dualPsi_omega_i}) for $\omega _{\alpha}$ and tensoring the result with $\Ll (-\omega _{\alpha})$ (note that $\nabla _{\omega _{\alpha}}$ is self-dual).

\begin{remark}
{The case of ${ \bg}_2$ differs from other types in rank two in that the bundle $\Psi _2^{\omega _{\alpha}}$ appears 
in the second block of the exceptional collection for ${ \bg}_2$ in Lemma \ref{lem:G_2-Demazure-FEC} (as opposed to ${\sf D}_{\beta}\Psi _1^{\omega _{\beta}}$ suggested by the cases of $A_2$ and ${\fb}_2$). The reason is that the bundle ${\sf D}_{\beta}\Psi _1^{\omega _{\beta}}$ is not exceptional (one of the two deviations in Proposition \ref{prop:exceptionality_D_i}): ${\sf D}_{\beta}\Psi _1^{\omega _{\beta}}$ is the kernel of the surjective map $\nabla _{\omega _{\beta}}\otimes \Ll (0)\rightarrow {\sf D}_{\beta}\Ll (\omega _{\beta})$ and the latter vector bundle ${\sf D}_{\beta}\Ll (\omega _{\beta})$ is not exceptional. In turn, ${\sf D}_{\beta}\Psi _1^{\omega _{\beta}}$ is not exceptional either.
}
\end{remark}

\begin{remark}\label{rem:identG_2}
{
The  $ \bg$-equivariant vector bundles in the above collection can be explicitly identified: 

\begin{itemize}[(1)]

\item[(1)] ${\sf D}_{\alpha}(\Psi _1^{\omega _{\alpha}}) = \Uu _2^{\perp}$ which is the kernel of canonical surjective morphism $\nabla _{\omega _{\alpha}}\otimes \Ll (0)\rightarrow \Uu _2^{\ast}$; the bundle $\Uu _2$ is defined by the short exact sequence (\ref{eq:weights_of_U2_G2}) below. 

\item[(2)] ${\sf D}_{\beta\alpha}(\Psi _1^{\omega _{\alpha}}) = {\sf S}$ where ${\sf S}$ is a rank 4 vector bundle called the spinor bundle defined by the short exact sequence (\ref{eq:S-to-U_2}) below. 

\item[(3)] ${\sf D}_{\alpha\beta\alpha}(\Psi _1^{\omega _{\alpha}}) = \Uu _2$ which is defined by the short exact sequence (\ref{eq:weights_of_U2_G2}) below. 

\item[(4)] ${\sf D}_{\beta\alpha\beta\alpha}(\Psi _1^{\omega _{\alpha}}) = \Ll (-\omega _{\alpha})$.
 
\item[(5)] ${\sf D}_{\alpha}(\Psi _2^{\omega _{\alpha}})=\Ee$ is a rank 11 vector bundle defined by the short exact sequence (\ref{eq:d_alphaPsi^2_alpha}) below.

\item[(6)] ${\sf D}_{\alpha}((\Psi _1^{\omega _{\alpha}})^{\ast}\otimes \Ll (-\omega _{\alpha}))=\Ll (-\omega _{\beta})$, by Remark \ref{rem:prop_pi_k_*Psi _l(-l)_othertypes}, (\ref{eq:pi_k_*Psi _l(-l)d}).

\end{itemize}

Let us give concrete descriptions of the high rank $ \bg$-equivariant vector bundles in the above list. The rank two bundle $\Uu _2$ fits into a short exact sequence 

\begin{equation}\label{eq:weights_of_U2_G2}
0\rightarrow \Ll (-\omega _{\alpha})\rightarrow \Uu _2\rightarrow \Ll (\omega _{\alpha}-\omega _{\beta})\rightarrow 0;
\end{equation}
the corresponding extension is given by the isomorphism $\Ext ^1_{ \bg/\fb}(\Ll (\omega _{\alpha}-\omega _{\beta}),\Ll (-\omega _{\alpha}))={\rm H}^1( \bg/ \fb,\Ll (-\beta))=k$. 
}
\end{remark}

The spinor bundle $\sf S$ fits into the short exact sequence:

\begin{equation}\label{eq:S-to-U_2}
0\rightarrow \Uu _2\rightarrow {\sf S}\rightarrow \Uu _2^{\ast}\otimes \Ll (-\omega _{\alpha})\rightarrow 0.
\end{equation}

Applying $\sd _{\alpha}$ to (\ref{eq:star}), and taking into account the isomorphisms $\sd _{\alpha}(\Ll (\omega _{\alpha}))=\Uu _2^{\ast}, \sd _{\alpha}(\Ll (2\omega _{\alpha}))={\sf S}^2\Uu _2^{\ast}$, 
we obtain:

\begin{equation}\label{eq:d_alphaPsi^2_alpha}
0\rightarrow \sd _{\alpha}(\Psi _2^{\omega _{\alpha}})\rightarrow {\sf Ker}(\nabla _{\omega _{\alpha}}^{\otimes 2}\rightarrow \nabla _{2\omega _{\alpha}})\otimes \Oo _{{ \bg}_2/{\fb}}\rightarrow \nabla _{\omega _{\alpha}}\otimes \Uu _2^{\ast}\rightarrow {\sf S}^2\Uu _2^{\ast}\rightarrow 0.
\end{equation}

Now $\sd _{\alpha}(\Psi _2^{\omega _{\alpha}})=:\Ee$ is a vector bundle on $ \bg _2/ \fb$ since ${\rm R}^1{\pi _{\alpha}}_{\ast}\Psi _2^{\omega _{\alpha}}=0$; the latter fact follows from the short sequence 

\begin{equation}
0\rightarrow \Omega ^2_{  \Pp (\nabla _{\omega _{\alpha}}^{\ast})}\otimes \Ll (2\omega _{\alpha})|_{ \bg/P _{\beta}}\rightarrow \Psi _2^{\omega _{\alpha}}\rightarrow \Ll (0)\rightarrow 0,
\end{equation}
and from the splitting type of $\Omega ^2_{  \Pp (\nabla _{\omega _{\alpha}}^{\ast})}\otimes \Ll (2\omega _{\alpha})|_{ \bg/P _{\beta}}$ to a line $l$ on $ \bg/P _{\beta}$ (cf. \cite[Section 4]{Str}) as $(\underbrace{-1,\dots,-1,}_5\underbrace{0,\dots,0}_{10})$. It also follows from the above that $\Ee$ fits into the short exact sequence 

\begin{equation}\label{eq:E}
0\rightarrow \Lambda ^2\Uu _2^{\perp}\rightarrow \Ee\rightarrow \Ll (0)\rightarrow 0;
\end{equation}

These exact sequences will be of help in cohomological computations related to the collection of Lemma \ref{lem:G_2-Demazure-FEC}.

\begin{lemma}\label{lem:dualG_2-Demazure-FEC}

Consider the following collection:



\qquad\qquad\qquad\quad\scalebox{.53}{
$$
\centerline{
\xymatrix{
 \A _{-6} & \A _{-5}& \A _{-4} & \A _{-3}& \A _{-2}& \A _{-1}& \A _0\\
 || & || & || & || & || & || & || \\
*++<1pt>[F]\txt{$\Ll (0)$}
&*++<10pt>[F]\txt{${\sf D}_{\beta\alpha\beta\alpha}(\Psi _1^{\omega _{\alpha}})\otimes \Ll (\rho)$  \\ \\ ${\sf D}_{\alpha}((\Psi _1^{\omega _{\alpha}})^{\ast}(-\omega _{\alpha}))\otimes \Ll (\rho)$}
&*++<10pt>[F]\txt{$(\Psi _1^{\omega _{\alpha}})^{\ast}\otimes \Ll (\omega _{\beta})$  \\ \\ ${\sf D}_{\alpha\beta\alpha}(\Psi _1^{\omega _{\alpha}})\otimes \Ll (\rho)$}
&*++<10pt>[F]\txt{${\sf D}_{\beta\alpha}(\Psi _1^{\omega _{\alpha}})\otimes \Ll (\rho)$  \\ \\ ${\sf D}_{\alpha}(\Psi _2^{\omega _{\alpha}})\otimes \Ll (\rho)$}
&*++<10pt>[F]\txt{$\Psi _2^{\omega _{\alpha}}\otimes \Ll (\rho)$  \\ \\ ${\sf D}_{\alpha}(\Psi _1^{\omega _{\alpha}})\otimes \Ll (\rho)$}
&*++<10pt>[F]\txt{$\Psi _1^{\omega _{\alpha}}\otimes \Ll (\rho)$ \\ \\ $\Psi _1^{\omega _{\beta}}\otimes \Ll (\rho)  $}
 &*++<1pt>[F]\txt{$\Ll (\rho)$}
}
}
$$
}\\\\
in which the bundles from the block ${\B _i}, i=0,\dots,6$ are shifted by $[-i]$ in the category $\Dd ^b( \bg _2/{\fb})$. In other words, the above collection is obtained from the collection Lemma \ref{lem:G_2-Demazure-FEC} by tensoring with $\Ll (\rho)$ and applying the shift functor relative to the corresponding block. Then the collection above is right dual to that of Lemma \ref{lem:G_2-Demazure-FEC}.
 \end{lemma}

\begin{remark}\label{rem:SteinbergG_2}
{The Steinberg weights for ${\bf G}_2$ are given by (see Section \ref{subsec:Steinberg})

{
$$\ -\rho, \ -\omega _{\alpha}, \ -\omega _{\beta}, \  -\omega _{\alpha}+\omega _{\beta},\ \omega _{\alpha}-\omega _{\beta}, \  2\omega _{\alpha}-\omega _{\beta}, \ -2\omega _{\alpha}+\omega _{\beta},$$

$$3\omega _{\alpha}-2\omega _{\beta},\ -3\omega _{\alpha}+2\omega _{\beta}, \ 3\omega _{\alpha}-\omega _{\beta}, \  -3\omega _{\alpha}+\omega _{\beta}, \ 0;$$
the weights of the decomposition factors of the collection of Lemma \ref{lem:G_2-Demazure-FEC} are the Steinberg weights and the dominant weight $\omega _{\alpha}$.
}
}
\end{remark}

%--------------------------------------------------------------
\subsection{ The quiver for ${\bf G}_2$}
%--------------------------------------------------------------

The picture below depicts the path quiver algebra of the above full exceptional collection. There are only zero-degree $\Hom$-spaces in the non-trivial direction in the collection of Lemma \ref{lem:G_2-Demazure-FEC}. The underlying graph of the quiver is the Bruhat graph for $ \bg _2/{\fb}$; the quiver path algebra is generated in degree one and the relations are quadratic that one reads off the $\Hom$--spaces next to the arrows of the quiver. This diagram is over a field of characteristic at least eleven or characteristic zero. Or over $\Z[1/6]$. Or rather, the $\RHom$'s have a good filtration over $\Z[1/6]$ and the labels are correct for $p>7$. For $p= 5,\;7$ we checked that the characters are correct. 

\vspace{0.4cm}

\qquad\qquad\qquad\scalebox{.58}{
 $$
\centerline{
\xymatrix{
& {\sd}_{\beta\alpha\beta\alpha}(\Psi _1^{\oa})\ar[rr]^{\nabla _{\oa}}\ar[rrdd]\ar@/^6pc/[rrrr]^{\nabla _{\oa}\oplus \sf k} & & (\Psi _1^{\omega _{\alpha}})^{\ast}(-\oa) \ar@/^6pc/[rrrr]^{\nabla _{\oa}\otimes\nabla _{\oa} \oplus \nabla _{\oa}\oplus \sf k}  \ar[rr]^{\nabla _{\oa}\oplus \sf k}\ar[rrdd] & & {\sd}_{\beta\alpha}(\Psi _1^{\oa}) \ar[rr]^{\nabla _{\oa}\oplus \sf k}\ar[rrdd]\ar@/^6pc/[rrrr]^{\nabla _{\oa}\oplus \sf k} & &  \Psi _2^{\oa} \ar[rr]^{\nabla _{\oa}}\ar[rrdd]\ar@/^6pc/[rrrd]^{\Lambda ^2 \nabla _{\oa}\oplus {\sf k}}  & & \Psi _1^{\oa}\ar[rd]_{\nabla _{\oa}} \\
\Ll (-\rho) \ar@/^6pc/[rrru]^{\nabla _{\oa}\otimes \nabla _{\ob}}\ar@/_6pc/[rrrd]_{\nabla _{2\oa}\oplus \nabla _{\ob}} \ar[rd]^{\nabla _{\oa}}\ar[ru]_{\nabla _{\ob}}  \quad  &   & &  & &  & &  &  & &  \quad \Ll (0)\\
& {\sd}_{\alpha}((\Psi _1^{\omega _{\alpha}})^{\ast}(-\oa))\ar[rr]^{\nabla _{\oa}}\ar[rruu]\ar@/_6pc/[rrrr]_{\nabla _{\oa}\otimes \nabla _{\oa}}  & & {\sd}_{\alpha\beta\alpha}(\Psi _1^{\oa})\ar[rr]^{\nabla _{\oa}\oplus \sf k}\ar[rruu] \ar@/_6pc/[rrrr]_{\nabla _{\oa}\oplus \sf k} & & {\sd}_{\alpha}(\Psi _2^{\oa})\ar[rr]^{\nabla _{\oa}\oplus \sf k}\ar[rruu]\ar@/_6pc/[rrrr]_{\nabla _{\oa}\otimes \nabla _{\oa}}  & &  {\sd}_{\alpha}(\Psi _1^{\oa})\ar[rruu]\ar@/_6pc/[rrru]_{\nabla _{\oa}}   \ar[rr]^{\nabla _{\oa}} & & \Psi _1^{\ob}\ar[ru]^{\nabla _{\ob}} \\
}
}
$$
}

\vspace{0.4cm}

The $\nabla$-filtrations on the tensor products of fundamental representations are given by:

\begin{itemize}

\item[(1)] $\nabla _{\omega _{\alpha}}\otimes \nabla _{\omega _{\beta}}$ is filtered by $\nabla _{\rho}, \nabla _{2\omega _{\alpha}}$, and $\nabla _{\omega _{\alpha}}$.

\item[(2)]  $\nabla _{\omega _{\alpha}}\otimes \nabla _{\omega _{\alpha}}$ is filtered by $\nabla _{2\omega _{\alpha}},  \nabla _{\omega _{\beta}}, \nabla _{\omega _{\alpha}}$, and $\sk$.

\end{itemize}

The Weyl group elements that we associate to the nodes are
as follows: The leftmost node belongs with the longest element $w_0$, the rightmost with the identity element. In the top row we have from left to right 
$s_\alpha s_\beta s_\alpha s_\beta s_\alpha$, 
$s_\beta  s_\alpha s_\beta s_\alpha $, 
$s_\alpha s_\beta s_\alpha $, 
$ s_\beta s_\alpha$, 
$s_\alpha $, 
and in the bottom row from left to right
$s_\beta s_\alpha s_\beta s_\alpha s_\beta$, 
$s_\alpha s_\beta s_\alpha s_\beta$, 
$s_\beta s_\alpha s_\beta$, 
$s_\alpha s_\beta$, 
$s_\beta$. 
%--------------------------------------------------------------
\subsection{ Type ${\bf G}_2$ over $\Z$}
%--------------------------------------------------------------

%{
The non-vanishing statements of Lemmas \ref{lem:dualA_2-Demazure-FEC}, \ref{lem:dualB_2-Demazure-FEC}, and \ref{lem:dualG_2-Demazure-FEC} boiled down to non-vanishing cohomology groups of line bundles either in degree zero, or in the top degree, or in degree one. By \cite[Corollary 2.6]{AndChev} and \cite[II, Section 8.8, (6)]{Jan}, those cohomology groups are free $\mathbb Z$-modules considered as cohomology groups of line bundles over $\Gm/\Bm$. Thus, the results of the above lemmas hold over $\Z$.
%}



\section{ The left dual exceptional collection to the $X_v$'s, conjecturally}\label{sec:dualFEC_conj}
%--------------------------------------------------------------

Recall the notation $\tP_v=\Ll(P(-e_v))$ and $\tQ_v=\Ll(Q(e_v))$ for $v\in W$.
%Unknown notation. Either $\tP_v^*$ or $\Ll(P(-e_v)^*)$.
Some evidence suggests that the following could be true:
\begin{equation}\label{eq:indPQ-rho} R\Gamma (\bg/\fb,\tP_v\otimes \tQ_{w_0v}\otimes \Ll (-\rho))= {\sf k}[-\ell(w_0v)].
\end{equation}
We have in any case, provided $0\in\overline C _{\mathbb Z}$,
\begin{equation}\label{eq:indStSt-rho}
 R\Gamma (\bg/\fb,\Ll(-e_v)\otimes \Ll(e_{w_0v})\otimes \Ll (-\rho))= {\sf k}[-\ell(w_0v)] .
\end{equation}
This uses that 
$-e_v+e_{w_0v}-\rho=(w_0v)^{-1}\rho-\rho$ and Theorem \ref{th:Bott-Demazure_th}.

There is also evidence that 
\begin{equation}\label{eq:QQ-rho}R\Gamma ({\bg/\fb},\tQ_v^{\ast}\otimes \Ll (-\rho)\otimes \tQ_w) \mbox{ vanishes unless } w_0 v \geq w,\end{equation}
and 
\begin{equation}\label{eq:QQ-rho=}R\Gamma ({\bg/\fb},\tQ_v^{\ast}\otimes \Ll (-\rho)\otimes \tQ_{w_0 v}))={\sf k}[-\ell(w_0 v)],
\end{equation}
and\begin{equation}\label{eq:PP-rho} R\Gamma ({\bg/\fb}, \tP_v\otimes \Ll (-\rho)\otimes \tP_w^{\ast}) \mbox{ vanishes unless } w_0 v \leq w,\end{equation}
and\begin{equation}\label{eq:PP-rho=} R\Gamma ({\bg/\fb}, \tP_v\otimes \Ll (-\rho)\otimes \tP_{w_0 v}^{\ast})={\sf k}[-\ell(w_0 v)].\end{equation}

Assuming (\ref{eq:QQ-rho}), (\ref{eq:QQ-rho=}), (\ref{eq:PP-rho}), (\ref{eq:PP-rho=}), it is not difficult to see with equation (\ref{eq:X in Q and P}) that 
$\RHom_{\bg/\fb}(\tX_v\otimes \Ll (\rho), \tX_w)$ vanishes unless $w_0 v = w$.
And that $\RHom_{\bg/\fb}(\tX_{w_0 v}\otimes \Ll (-\rho),\tX_{v}))={\sf k}[-\ell(w_0v)]$.

We can now state:

\begin{conjecture}\
The left dual exceptional collection to the $\tX_v$'s consists of the $\tX_{w_0v}\otimes\Ll(\rho)[-\ell(w_0v)]$.
\end{conjecture}



\begin{thebibliography}{10}

\bibitem{Achar}
P.Achar, {\it Staggered $t$-structures on derived categories of equivariant coherent sheaves}, Int. Math. Res. Not. IMRN 2009, no. 20, 3843--3900.

\bibitem{AchTr}
P.Achar, D.Treumann, {\it Baric structures on triangulated categories and coherent sheaves}, Int. Math. Res. Not. IMRN 2011, no. 16, 3688--3743.

%\bibitem{Ach}
%P. Achinger, {\it Frobenius push-forwards on quadrics}, Communications in Algebra, 40:8 (2012)
\bibitem{Ana}
A.Ananyevskiy, {\it On the algebraic K-theory of some homogeneous varieties}. Doc. Math. 17 (2012),  167--193

\bibitem{AAGZ}
A.Ananyevskiy, A.Auel, S.Garibaldi, K.Zainoulline, {\it 
Exceptional collections of line bundles on projective homogeneous varieties}, Adv. Math. 236 (2013), 111--130

\bibitem{AndChev}
H.H.Andersen, {\it Filtrations of cohomology modules for Chevalley groups}, Ann. Sci. \'Ecole Norm. Sup. (4) 16 (1983), no. 4, 495--528. 

\bibitem{And}
H.H.Andersen, {\it Schubert varieties and Demazure's character formula}, 
Invent. Math. 79 (1985), no. 3, 611--618.

\bibitem{APW}
H.H.Andersen, P.Polo, K.Wen, {\it Representations of quantum algebras}, Invent. Math. 104 (1991), no. 1, 1--59.

\comment{
\bibitem{BBM}
A.Beilinson, R.Bezrukavnikov, I.Mirkovi$\acute c$, {\it Tilting exercises}, Mosc. Math. J. 4 (2004), no. 3, 547--557, 782.
}

\bibitem{BGG}
I.Bern\u{s}te\u{i}n, I.M.Gel'fand, S.I.Gel'fand, {\it Schubert cells, and the cohomology of the spaces G/P}, Uspehi Mat. Nauk 28 (1973), no. 3(171), 3--26. 

\bibitem{Bondal}
A.Bondal, {\it Representations of associative algebras and coherent
sheaves}, (Russian) Izv. Akad. Nauk SSSR Ser. Mat. 53 (1989), no. 1,
pp. 25--44; translation in Math. USSR Izv. 34 (1990), no. 1,
pp. 23--42.

\bibitem{BK}
M.Brion, S.Kumar, {\it Frobenius splitting methods in geometry and representation theory}, Progress in Mathematics, 231, 2005

\bibitem{BLVdB}
R.-O.Buchweitz, G.J.Leuschke, M.Van den Bergh, {\it On the derived category of Grassmannians in arbitrary characteristic}. Compos. Math. 151 (2015), no. 7, 1242--1264.

\bibitem{Buehler} T.B\"uhler, {\it Exact categories}, Expositiones Mathematicae
28 (2010), Issue 1,  1--69.

\bibitem{CPSWvdK}
E.Cline, B.Parshall, L.Scott, W. van der Kallen, {\it Rational and generic cohomology}, Invent. Math. 39 (1977), no. 2, 143--163.

\bibitem{CPS}
E.Cline, B.Parshall, L.Scott, {\it Induced modules and extensions of representations},   Invent. Math. 47 (1978), no. 1, 41--51.

\bibitem{CPSMackey}
E.Cline, B.Parshall, L.Scott, {\it A Mackey imprimitivity theory for algebraic groups}. Math. Z. 182 (1983), no. 4, 447--471.

\bibitem{CPSHW}
E.Cline, B.Parshall, L.Scott, {\it Finite dimensional algebras and highest
weight categories},   J. reine angew. Math. 391 (1988), 85--99.

\comment{
\bibitem{CPSPerv}
E.Cline, B.Parshall, L.Scott, {\it Simulating perverse sheaves in modular representation theory}, Algebraic groups and their generalizations: classical methods, 63--104, Proc. Sympos. Pure Math., 56, Part 1, Amer. Math. Soc., Providence, RI, 1994.
}

\bibitem{Daw}
S.Dawydiak, {\it A coherent categorification of the based ring of the lowest two-sided cell}, 
\href{https://doi.org/10.48550/arXiv.2111.15648}{arXiv:2111.15648}

\bibitem{Demres}
M.Demazure, {\it D\'esingularisation des vari\'et\'es de Schubert g\'en\'eralis\'ees}, Ann. Sci. \'Ecole Norm. Sup. (4) 7 (1974), 53--88.

\bibitem{Dem}
M.Demazure, {\it A very simple proof of Bott's theorem}, 
Invent. Math. 33 (1976), no. 3, 271--272.

\bibitem{Donkin1}
S.Donkin, {\it A filtration for rational modules}, Math. Zeit. 177 (1981), 1--8

\bibitem{Donkin2}
S.Donkin, {\it Finite resolutions of modules for reductive algebraic groups}, J. Algebra 101 (1986), no. 2, 473--488. 

\bibitem{Donsur}
S.Donkin, {\it Tilting modules for algebraic groups and finite dimensional algebras}, Handbook of tilting theory, LMS vol. 332 (2007)

\bibitem{Ef}
A.Efimov, {\it Derived categories of Grassmannians over integers and modular representation theory}, Adv. Math. 304 (2017), 179--226. 

\bibitem{FKM}
E.Feigin, A.Khoroshkin, and Ie.Makedonskyi, {\it Parahoric Lie algebras and parasymmetric Macdonald polynomials},
\href{https://doi.org/10.48550/arXiv.2311.12673}{arXiv.2311.12673}

\bibitem{Fon}
A.Fonarev, {\it Derived categories of Grassmannians: a survey}, 
\href{https://doi.org/10.48550/arXiv.2407.07455}{arXiv.2407.07455}

\bibitem{FvdK}
V.Franjou, W. van der Kallen, 
{\it Power reductivity over an arbitrary base}, 
Documenta Mathematica, Extra Volume Suslin (2010) , 171--195.

\bibitem{Frid}
E.Friedlander, {\it A canonical filtration for certain rational modules}, Math. Z. 188 (1985), no. 3, 433--438

\bibitem{GrahamKumar}
W.Graham, S.Kumar. {\it On positivity in T-equivariant K-theory of flag varieties}, Int. Math. Res. Not. IMRN (2008),  DOI: 10.1093/imrn/rnn093.

%\bibitem{Hab}
%W. Haboush, {\it Homogeneous vector bundles and reductive subgroups of reductive algebraic groups}. Amer. J. Math. 100 (1978), no. 6, 1123--1137

\bibitem{Huyb}
D.Huybrechts, {\it Fourier--Mukai transforms in algebraic geometry},
Oxford Mathematical Monographs,  (2006)

\bibitem{Humphreys}
J.E.Humphreys, {\it Reflection Groups and Coxeter Groups}, Cambridge University Press, 1992 

\bibitem{Hummodular}
J.E.Humphreys, {\it Modular representations of finite groups of Lie type}, London Mathematical Society Lecture Note Series, 326. Cambridge University Press, Cambridge, 2006

\bibitem{Jan}
J.C.Jantzen, {\it Representations of algebraic groups}, 2nd edition, Mathematical Surveys and Monographs, 107, (2003)

%\bibitem{Ka}
%M. Kaneda, {\it Exceptional collections of sheaves on quadrics in positive characteristic}, 
%S$\tilde a$o Paulo Journal of Mathematical Sciences 8, 1 (2014), 117--156

%\bibitem{Kan}
%M. Kaneda, {\it On the Frobenius direct image of the structure sheaf of a homogeneous projective variety}, J. Algebra 512 (2018) 160--188

%\bibitem{KK}
%B. Kostant, S. Kumar, {\it T--equivariant K--theory of generalized flag varieties}, J. Differential Geometry 32 (1990), 549--603.

\comment{
\bibitem{Krpol}
H.Krause, {\it Koszul, Ringel and Serre duality for strict polynomial functors}. Compos. Math. 149 (2013), no. 6, 996--1018.
}

\bibitem{Kr}
H.Krause, {\it Homological theory of representations}. Cambridge Studies in Advanced Mathematics, 195. Cambridge University Press, Cambridge, 2022.

\bibitem{KuzHPD}
A.Kuznetsov, {\it Homological projective duality}, Publ. Math. Inst. Hautes \'Etudes Sci. 105 (2007), 157--220.

\bibitem{Kuzrat}
A.Kuznetsov, {\it Derived categories of cubic fourfolds}, Cohomological and geometric approaches to rationality problems, 219--243, Progr. Math., 282, (2010)

\bibitem{Kuzbasech}
A.Kuznetsov, {\it Base change for semiorthogonal decompositions},
Compos. Math. 147 (2011), no. 3, 852--876. 

\bibitem{KuzPol}
A.Kuznetsov, A.Polishchuk, {\it Exceptional collections on isotropic Grassmannians}, J. Eur. Math. Soc. (JEMS) 18 (2016), no. 3, 507--574

\bibitem{Math}
O.Mathieu, {\it Filtrations of B-modules}, Duke Math. J. 59 (1989), no. 2, 421--442

\bibitem{MathG}
O.Mathieu, {\it Filtrations of G-modules}, Annales scientifiques de l'\'Ecole Normale Sup\'erieure, Serie 4, Volume 23 (1990) no. 4, 625-644.

\bibitem{MR}
C.Mautner, S.Riche, {\it On the exotic $t$--structure in positive characteristic}, International Mathematics Research Notices, Vol. 2016, No. 18, pp. 5727--5774.

\bibitem{Neem}
A.Neeman, {\it The Grothendieck duality theorem via Bousfield's techniques and Brown representability},  J. Amer. Math. Soc. 9 (1996), no. 1, 205--236.

%\bibitem{Neemtr}
%A. Neeman, {\it Triangulated categories}, Annals of Mathematics Studies, 148 (2001)

%\bibitem{Per}
%A. Perry, {\it Noncommutative homological projective duality}, Adv. Math. 350 (2019), 877--972

\bibitem{Polishchuk}
A.Polishchuk, {\it K-theoretic exceptional collections at roots of 
unity}, J. K-Theory 7, (2010), 169--201.

\comment{
\bibitem{PoloCM}
P.Polo, {\it On Cohen-Macaulay posets, Koszul algebras and certain modules associated to Schubert varieties}, Bull. London Math. Soc. 27 (1995), no. 5, 425--434.
}

\bibitem{Polo3}
P.Polo, {\it Vari\'et\'es de Schubert et excellentes filtrations}, Ast\'erisque 173--174 (1989), 281--311.

\bibitem{Sam}
A.Samokhin, {\it The Frobenius morphism on flag varieties, I}, {arXiv:1410.3742}

\bibitem{Serre}
J.-P.Serre, {\it Groupe de Grothendieck des sch\'emas en groupes r\'eductifs d\'eploy\'es},
Publ. Math. IH\'ES, 34 (1968), pp. 37-52.

\bibitem{St}
R.Steinberg, {\it On a theorem of Pittie}, Topology vol. 14 (1975) 173--117.

\bibitem{Str}
E.Strickland, {\it Lines in $ G/P$}, Math.Z. 242 (2002), 227--240.

%\bibitem{Tho}
%R.W.Thomason, {\it The classification of triangulated subcategories},
%Compositio Math. 105 (1997), no. 1, 1--27.

\bibitem{WvdK}
W. van der Kallen, {\it Longest weight vectors and excellent filtrations}, Math. Z. 201 (1989), 19--31.

\bibitem{WvdKTata}
W. van der Kallen, {\it Lectures on Frobenius splittings and B-modules}. Notes by S. P. Inamdar. Published for the Tata Institute of Fundamental Research, Bombay; by Springer-Verlag, Berlin (1993)

\end{thebibliography}



\end{document}
